% ============================================================================
%  BURST / neutt manual
%  Standard-Model + Delta-N_eff sector of the JAX/GPU port of BURST, and (Part II) the dark sectors.
%  Compile:  module load texlive/2024 && latexmk -pdf burst_manual.tex
%       or:  pdflatex burst_manual.tex   (run twice for the ToC / refs)
% ============================================================================
\documentclass[11pt]{article}

\usepackage[a4paper,margin=1in]{geometry}
\usepackage{amsmath,amssymb}
\usepackage{booktabs}
\usepackage{longtable}
\usepackage{array}
\usepackage{framed}
\usepackage{xcolor}
\usepackage{listings}
\usepackage{seqsplit}
\usepackage{enumitem}
\usepackage{titlesec}
\usepackage[colorlinks=true,linkcolor=blue!55!black,urlcolor=teal!65!black,
            citecolor=blue!55!black]{hyperref}

% ---- inline literal code: \detokenize prints _, %, # ... verbatim in tt;
%      \seqsplit lets a long identifier wrap rather than run off the page.
%      \expanded forces \detokenize to expand BEFORE \seqsplit iterates it.
%      Spaces are kept: \code@split runs \seqsplit word by word (seqsplit drops spaces).
\makeatletter
\def\code@split#1 #2\@nil{\seqsplit{#1}\ifx\relax#2\relax\else\ \code@split#2\@nil\fi}
\newcommand{\code}[1]{\texttt{\expandafter\code@split\detokenize{#1} \@nil}}
\makeatother

% ---- colours ----
\definecolor{codebg}{rgb}{0.97,0.97,0.95}
\definecolor{kw}{rgb}{0.10,0.25,0.62}
\definecolor{cmt}{rgb}{0.32,0.45,0.32}
\definecolor{str}{rgb}{0.58,0.20,0.10}
\definecolor{frbg}{rgb}{0.93,0.95,0.99}
\definecolor{ntbg}{rgb}{0.99,0.96,0.90}

% ---- code listings ----
\lstdefinestyle{py}{
  language=Python,
  basicstyle=\ttfamily\footnotesize,
  backgroundcolor=\color{codebg},
  keywordstyle=\color{kw}\bfseries,
  commentstyle=\color{cmt}\itshape,
  stringstyle=\color{str},
  showstringspaces=false,
  breaklines=true,
  breakatwhitespace=false,
  breakindent=10pt,
  postbreak=\mbox{\textcolor{gray}{$\hookrightarrow$}\space},
  columns=fullflexible,
  frame=single,
  framerule=0.2pt,
  rulecolor=\color{gray!40},
  xleftmargin=4pt,
  literate={~}{{$\sim$}}1 {>=}{{$\geq$}}1 {<=}{{$\leq$}}1,
}
\lstdefinestyle{sh}{
  language=bash,
  basicstyle=\ttfamily\footnotesize,
  backgroundcolor=\color{codebg},
  keywordstyle=\color{kw},
  commentstyle=\color{cmt}\itshape,
  showstringspaces=false,
  breaklines=true,
  breakatwhitespace=false,
  breakindent=10pt,
  postbreak=\mbox{\textcolor{gray}{$\hookrightarrow$}\space},
  columns=fullflexible,
  frame=single,
  framerule=0.2pt,
  rulecolor=\color{gray!40},
  xleftmargin=4pt,
}
\lstset{style=py}

% ---- call-out boxes -------------------------------------------------------
% Fortran <-> JAX mapping box
\newenvironment{fmap}
  {\par\medskip\begin{shaded*}\small\noindent\textbf{[Fortran $\leftrightarrow$ JAX]}\ }
  {\end{shaded*}\medskip}
% Difference / missing-feature box
\newenvironment{diffbox}
  {\par\medskip\begin{shaded*}\small\noindent\textbf{[Port note]}\ }
  {\end{shaded*}\medskip}
% User-approved default / decision that must not be changed silently
\newenvironment{decisionbox}
  {\par\medskip\begin{shaded*}\small\noindent\textbf{[Deliberate choice]}\ }
  {\end{shaded*}\medskip}
% shaded* with a colour
\colorlet{shadecolor}{frbg}

% convenient physics macros
\newcommand{\Neff}{N_{\mathrm{eff}}}
\newcommand{\Tcm}{T_{\mathrm{cm}}}
\newcommand{\Tpl}{T_{\mathrm{pl}}}
\newcommand{\Tg}{T_\gamma}
\newcommand{\Tnu}{T_\nu}
\newcommand{\eps}{\varepsilon}
\newcommand{\dd}{\mathrm{d}}
\newcommand{\GF}{G_{\mathrm F}}
\newcommand{\me}{m_e}
\newcommand{\rhonu}{\rho_\nu}
\newcommand{\rhog}{\rho_\gamma}
\newcommand{\Yp}{Y_{\mathrm P}}
\newcommand{\DoverH}{\mathrm{D/H}}
\newcommand{\dNeff}{\Delta\Neff}

\titleformat{\paragraph}[runin]{\bfseries}{}{0pt}{}[.]

\setlength{\parskip}{0.35em}
\setlength{\parindent}{0pt}
\setlength{\emergencystretch}{3em}   % let TeX relax tight lines instead of overflowing

% reusable wrapping column types for the wide reference tables
\newcolumntype{R}[1]{>{\raggedright\arraybackslash}p{#1}}

% ============================================================================
\begin{document}

\begin{titlepage}
\centering
\vspace*{2.2cm}
{\Huge\bfseries BURST \,/\, \texttt{neutt}}\\[0.6cm]
{\LARGE A Manual for the JAX/GPU Port of\\[0.2cm]
Neutrino Boltzmann Transport in Big Bang Nucleosynthesis}\\[1.1cm]
{\large\itshape Standard Model, $\dNeff$, flavour mixing, and the dark sectors}\\[2.0cm]

\begin{minipage}{0.82\textwidth}\centering\small
This manual documents the neutrino-transport sector of \texttt{neutt}, the
differentiable JAX/GPU re-implementation of the Fortran code \textbf{BURST}
(Grohs, Fuller, Kishimoto, Paris \& Vlasenko 2016).  Part one covers the
Standard-Model weak-decoupling calculation of $\Neff$, the neutrino spectral
distortions, finite-temperature quantum-electrodynamics (QED) corrections (through
the next-to-leading-order (NLO) $\mathcal O(e^3)$ term, on by default), the coupling
to the \textbf{LINX} BBN abundance
code, and the neutrino flavour-mixing refinement (the QKE, or quantum kinetic
equation, density-matrix
transport).  Part two covers the interacting dark sectors of the \code{dark}
branch: the dark boson decay sectors (vector $Z'$ / majoron), the thermal-majoron
$\dNeff$, and the $\chi\leftrightarrow\nu$ contact sector.  The emphasis throughout
is on the reasons behind the difficult numerical design choices.

Throughout, the JAX implementation is mapped back to its Fortran BURST
ancestor, and every place where the port omits a BURST feature, or takes a
different numerical route, is flagged in a \textbf{[Port note]} box.
\end{minipage}

\vfill
{\small Generated \today.}
\end{titlepage}

\tableofcontents
\newpage

% ===========================================================================
\section{Introduction and scope}
% ===========================================================================

\subsection{What BURST is}

BURST is a Fortran Big Bang Nucleosynthesis (BBN) code written by Evan Grohs,
George Fuller, and collaborators.  Its defining feature, relative to the
standard BBN codes, is that it does not approximate the neutrino seas as
a single thermal bath with one temperature $\Tnu$.  Instead it solves the
\textbf{full momentum-resolved Boltzmann transport equation} for each neutrino
species, with the complete $2\!\to\!2$ weak collision integrals, simultaneously
with the nuclear reaction network and the plasma thermodynamics.  This yields
the out-of-equilibrium neutrino spectral distortions that develop during
weak decoupling, and a correspondingly self-consistent effective number of
neutrino species $\Neff$.  The methodology is described in

\begin{quote}\small
E.~Grohs, G.~M.~Fuller, C.~T.~Kishimoto, M.~W.~Paris, A.~Vlasenko,
``Neutrino energy transport in weak decoupling and big bang
nucleosynthesis,'' Phys.\ Rev.\ D \textbf{93}, 083522 (2016),
\href{https://arxiv.org/abs/1512.02205}{arXiv:1512.02205}.
\end{quote}

We refer to this throughout as \textbf{the BURST paper}.  The companion review
E.~Grohs \& G.~M.~Fuller, ``Big Bang Nucleosynthesis''
(\href{https://arxiv.org/abs/2301.12299}{arXiv:2301.12299}) is the best
single reference for the surrounding cosmology, and we assume the reader is
comfortable with its level.

\subsection{What \texttt{neutt} is}

\texttt{neutt} is a from-scratch re-implementation of the BURST
neutrino sector in JAX, written to run on GPU and to be
differentiable.  (The package's former names are \texttt{jaxburst} and \texttt{newt}.)  The
motivation is twofold:

\begin{enumerate}[nosep]
\item \textbf{Speed.}  The Fortran BURST takes $\mathcal{O}(600\text{--}1000)$
  CPU-hours per cosmology, almost all of it in the transport collision
  integrals.  Recasting those integrals as dense tensor contractions
  (Section~\ref{sec:kerneltensor}) lets a single GPU evaluate them, and lets
  many cosmologies be batched with \code{vmap} (JAX's automatic vectorization,
  Section~\ref{sec:jaxprimer}).
\item \textbf{Differentiability and embeddability.}  A transport solver that JAX
  can compile (\code{jit}) and differentiate (\code{grad}), with a clean
  \texttt{transport(background)$\to$spectra} seam, can be dropped into a
  gradient-based Markov-chain Monte Carlo (MCMC) sampler, or used to replace the
  thermal-$\Tnu$ approximation inside the differentiable BBN code \textbf{LINX}
  (Giovanetti et al.\ 2024, \href{https://arxiv.org/abs/2408.14538}{arXiv:2408.14538}).
\end{enumerate}

The guiding principle of the port is \textbf{accuracy parity first}: the
collision operator is a near-cancellation of large gain and loss terms, so the
port runs in float64, reproduces BURST's conservation ``sum rules,'' and is
validated bit-for-bit against an instrumented build of the original Fortran
wherever possible.  Speed is pursued only where it costs no accuracy.

\subsection{Scope of this manual}

The manual is in two parts.  \textbf{Part one (Sections~\ref{sec:physics}--%
\ref{sec:extend})} covers the
Standard-Model neutrino transport and everything needed
to drive it: the background plasma thermodynamics, finite-temperature QED corrections,
the weak collision kernels, the coupled decoupling solver that produces $\Neff$
and the per-species spectra, the exponential-midpoint hybrid solver, the coupling
into LINX for light-element abundances, and the flavour-mixing refinement
(the quantum-kinetic-equation (QKE) density-matrix transport, which is
Standard-Model neutrino-oscillation
physics), together with the simplest beyond-Standard-Model handle, the
inert dark-radiation parameter $\dNeff$ (extra relativistic energy density
that redshifts like radiation and does not interact).

\textbf{Part two (Sections~\ref{sec:dark}--\ref{sec:darkextend})} covers the
interacting dark sectors, which live on the repository's \code{dark} branch only:
the dark boson decay sectors (a vector $Z'$ and a majoron, coupled to neutrinos
through $X\leftrightarrow\nu\bar\nu$), the thermal-majoron $\dNeff$ case study, and
the $\chi\leftrightarrow\nu$ contact sector (a dark-matter fermion scattering on
neutrinos), the last documented up to its in-progress lepton-kernel source fix
(Section~\ref{sec:chinu}).  These are the original research purpose of BURST, and
the emphasis there is on the rationale behind their delicate numerical design.  With
the dark sectors off (the default) every solver is the Standard-Model one of Part
one, operation for operation.

\textbf{Genuinely out of scope}, because the JAX port does not implement them, are
BURST's internal nuclear reaction network and its recombination / sound-horizon
(\code{sdratio}) machinery; abundances instead come from LINX (see
Appendix~\ref{sec:fortranmap} for the full ported-vs-not table).

\subsection{How the Fortran mapping is presented}

Two boxed annotations recur throughout:

\begin{fmap}
identifies the Fortran BURST routine, file, or variable that a piece of
\texttt{neutt} corresponds to, so the reader can cross-read the two trees.
\end{fmap}

\begin{diffbox}
flags a place where the port omits a BURST feature, or deliberately
takes a different numerical approach than the Fortran (and why).
\end{diffbox}

A note on provenance: \texttt{neutt} is validated against a single, read-only
Standard-Model Fortran reference, a snapshot of the BURST code (modules
\code{mainvar}, \code{transvar}, \ldots) that ships in \code{fortran_reference/},
which compiles, runs, and serves as the oracle the port is checked against.
Wherever the port deliberately departs from it numerically --- the corrected
wide-regime lepton-kernel quadrature (Section~\ref{sec:lepkernels}) and the
default-on NLO $\mathcal O(e^3)$ QED term (Section~\ref{sec:nloqed}) --- a
\textbf{[Port note]} box says so where the code appears.  Port notes also mark
choices where the port deliberately matches the reference bit-faithfully, such as
the high-occupation interpolation clamp (Section~\ref{sec:interpnum}) and the
combined precision gate with its noise floor (Section~\ref{sec:gatenum}).

\subsection{Notation and conventions}
\label{sec:notation}

These conventions follow the BURST paper.  The most important point, and
the most common cross-group misreading, is the \textbf{two-temperature split}:

\begin{itemize}[nosep]
\item $\Tcm$: the \textbf{comoving temperature parameter}.  It is
  not a physical neutrino temperature; it is a redshift clock,
  $\Tcm(a)=\Tcm^{\rm init}\,a_{\rm init}/a$ ($a$ the scale factor; $a_{\rm init}$ and
  $\Tcm^{\rm init}$ the values at the start of the solve), i.e.\ an energy scale that redshifts
  exactly as $1/a$.  Comoving momenta are frozen in units of $\Tcm$.
\item $\Tpl\equiv\Tg$: the physical \textbf{plasma} (photon $+\,e^\pm\,+$
  baryon) temperature.  After $e^\pm$ annihilation transfers entropy to the
  photons, $\Tpl/\Tcm\to(11/4)^{1/3}\simeq1.401$.
\item $\Tnu$: a physical neutrino temperature; only meaningful in the
  thermal-fluid approximation (LINX), where the spectrum is assumed Fermi--Dirac.
  BURST/\texttt{neutt} resolve $f(\eps)$ and never assume it.
\end{itemize}

\begin{itemize}[nosep]
\item $\eps\equiv E_\nu/\Tcm$: the dimensionless \textbf{comoving energy}
  variable that labels the momentum grid.  Defined with the energy $E$;
  for the relativistic neutrinos in the transport window it equals $p/\Tcm$.
  The grid is linear in $\eps$ on $[0,\eps_{\max}]$.
\item $f_{\nu_i}(\eps)$: the neutrino occupation number (occupation
  probability) for species $i$.  Equilibrium value
  $f^{\rm eq}(\eps)=1/(e^{\eps}+1)$ (Fermi--Dirac with chemical potential $\mu=0$,
  at $\Tnu=\Tcm$).
\item $\delta f\equiv\big[f(\eps)-f^{\rm eq}(\eps)\big]/f^{\rm eq}(\eps)$:
  the fractional \textbf{spectral distortion}.  The absolute distortion
  $f-f^{\rm eq}=f^{\rm eq}\,\delta f$ carries energy with an $\eps^3$ weight:
  $\Delta(\dd\rho/\dd\eps)/\rho=(120/7\pi^4)\,\eps^3[f-f^{\rm eq}]$, with $\rho$ the
  equilibrium energy density of that single species.
\item $\Neff$: the effective number of neutrino species, defined through the
  radiation energy density at photon decoupling,
  \begin{equation}
    \rho_{\rm rad}=\Big[2+\tfrac{7}{4}\big(\tfrac{4}{11}\big)^{4/3}\Neff\Big]
    \frac{\pi^2}{30}\,\Tg^4 .
    \label{eq:neffdef}
  \end{equation}
  Equivalently $\Neff=\frac{8}{7}\big(\frac{11}{4}\big)^{4/3}\rhonu/\rhog$
  evaluated late, after $e^\pm$ annihilation completes
  (Section~\ref{sec:neff}).
\item $\Yp$: the $^4$He \textbf{mass} fraction, $\Yp\simeq2(n/p)/(1+n/p)$ with $n/p$
  the neutron-to-proton ratio when nucleosynthesis begins;
  $\DoverH$ is the deuterium abundance ratio.  These are the BBN observables.
\item $\GF$, $\sin^2\theta_W$, $\alpha$, $\me$: Fermi constant, weak mixing
  angle, fine-structure constant ($\alpha=e^2/4\pi$ with $e$ the elementary charge,
  so $\mathcal O(e^3)$ counts powers of $e$), electron mass.
\end{itemize}

The integration variable for the decoupling ordinary differential equation (ODE)
is $\ln a$ (the number of
$e$-folds since the start); $\Tcm=\Tcm^{\rm init}e^{-\ln a}$ analytically.  All
energies are in MeV and energy densities in MeV$^4$ unless noted.

\subsection{JAX concepts used in this manual}
\label{sec:jaxprimer}

The port is written in JAX, and a handful of ideas from JAX and its compiler, XLA
(Accelerated Linear Algebra), recur throughout and
explain many of the design choices and timings.  This is a one-page glossary for
a reader fluent in physics but new to JAX; nothing here is BBN-specific.

\begin{description}[leftmargin=1.4em,style=nextline,nosep]
\item[Trace and compile.]  When you first call a function wrapped in \code{jax.jit}
  (just-in-time compilation), JAX runs your
  Python once symbolically (``tracing'') to record the array operations,
  hands that graph to the XLA compiler, and caches the compiled executable.
  Every later call reuses it.  So the first call is slow (it compiles) and later
  calls are fast (they just run the cached graph).
\item[Cold / maiden vs.\ warm.]  We call the first, compile-paying call
  cold (or maiden); a later, graph-reusing call warm.  The
  cold cost is one-time; the warm cost is what matters for a scan or an MCMC.
  This is why every solver follows a build-once, reuse pattern: build the
  solver, then call it many times (Section~\ref{sec:knobs}).
\item[Traced vs.\ build-time (static).]  A parameter is traced if it
  enters the graph as data: you can change it on every warm call with
  no recompile (MCMC-safe).  It is build-time / static if changing
  it changes the shape of the computation, forcing a recompile.  Which
  physical inputs are which is the subject of Section~\ref{sec:knobs}.
\item[PyTree.]  A nested container (dict / tuple / list) of arrays that JAX
  treats as one structured object.  The evolved state $\{\Tpl,\ f\}$ is a PyTree;
  so are the returned output dicts.
\item[Immutable arrays.]  JAX arrays cannot be mutated in place; an ``edit''
  returns a new array.  The demos use \code{x.at[i].add(v)} / \code{.set(v)} for
  this (it reads ``$x$ with index $i$ incremented by $v$'').
\item[\code{vmap} / \code{pmap}.]  Automatically vectorize a function over a new
  leading axis (a batch of cosmologies, say), and \code{pmap} spreads that
  batch across devices.  This is the ``across-calls'' GPU win.
\item[Host vs.\ device.]  Host $=$ CPU (where Python and some setup run);
  device $=$ GPU (where the compiled kernels run).  The coupled-BBN driver
  deliberately splits work across the two (Section~\ref{sec:linx}).
\item[\code{lax.scan}.]  A compiled fixed-length loop (a \code{for} loop XLA can
  see through).  The hybrid head is a \code{lax.scan} of a fixed number of steps.
\item[\code{einsum}.]  Einstein-summation contraction of several arrays in one
  fused operation, the form the collision integral is recast into
  (Section~\ref{sec:kerneltensor}).
\item[Fusion and constant-folding.]  XLA fuses adjacent array operations
  into one GPU kernel; the grouping it chooses can change the floating-point
  reduction order, so two builds of the ``same'' code can round slightly
  differently (this is behind the executable-to-executable $\Neff$ spread of
  Section~\ref{sec:hybrid}).  Constant-folding is XLA evaluating an
  input-independent sub-computation at compile time; usually harmless, but it
  can mishandle a near-cancellation (Section~\ref{sec:gatenum}).
\item[Adjoint.]  The strategy a differentiable ODE solver uses to back-propagate
  gradients through the integration.  It affects gradient cost and compile
  time but not the forward solution (Section~\ref{sec:extend}).
\item[\code{diffrax} PID controller.]  \code{diffrax} is the JAX ODE-solver
  library used here; its adaptive step size is set by a PID
  (proportional--integral--derivative) controller on the
  Runge--Kutta embedded error estimate (Section~\ref{sec:odenum}).
\end{description}

A reader who wants to skip this can: the terms are re-explained briefly where
they first do real work.

\subsection{The canonical Standard-Model result}
\label{sec:canonical}

The benchmark this whole code exists to compute is the Standard-Model (SM)
$\Neff$.  The \textbf{default} includes the NLO $\mathcal O(e^3)$
finite-temperature QED correction (Section~\ref{sec:nloqed}), and
\texttt{neutt}'s full Standard-Model decoupling produces
\[
  \boxed{\ \Neff^{\rm SM}=3.04332\ }\qquad(\text{NLO QED, the default}),
\]
which reconciles with the NLO precision calculation of Cielo, Escudero,
Mangano \& Pisanti (\href{https://arxiv.org/abs/2306.05460}{arXiv:2306.05460},
``$\Neff$ in the Standard Model at NLO is $3.043$'').  Turning the NLO term
off (\code{qed_nlo=False}) recovers the leading-order (LO) value
\[
  \Neff = 3.0442761,
\]
and evaluating that at Fortran BURST's own $\sin^2\theta_W=0.23$ (\code{SIN2THW_BURST}; production
uses the more precise $0.23122$, which raises $\Neff$ by $+5.1\times10^{-5}$) gives $3.0442252$,
the bit-faithful BURST replica (it matches a collaborator's BURST oracle
to $5.5\times10^{-7}$, and reproduces the precision results of
Akita \& Yamaguchi, \href{https://arxiv.org/abs/2005.07047}{arXiv:2005.07047};
Froustey, Pitrou \& Volpe, \href{https://arxiv.org/abs/2008.01074}{arXiv:2008.01074};
and Bennett et al., \href{https://arxiv.org/abs/2012.02726}{arXiv:2012.02726},
all $\simeq3.0440$).  The NLO term shifts $\Neff$ by $-9.5\times10^{-4}$.  Turning
the finite-$T$ QED correction to the plasma equation of state (EoS) off entirely
(\code{use_renorm=False}) gives the
``decoupling-only'' $\Neff=3.034029$.  These numbers recur as pins throughout the
manual; keep them in mind as the correctness anchors.

\subsection{Roadmap: how to read this manual}
\label{sec:roadmap}

The manual is layered so a reader can stop at the depth they need.

\begin{itemize}[nosep]
\item \textbf{Just want $\Neff$ or a spectrum.}  Skim the conventions
  (Section~\ref{sec:notation}) and the canonical result
  (Section~\ref{sec:canonical}), then run Quickstart~1 (Section~\ref{sec:install}).
\item \textbf{Want light-element abundances.}  Go to the LINX coupling
  (Section~\ref{sec:linx}) and Quickstart~3; the SM transport runs underneath.
\item \textbf{Driving an MCMC or a parameter scan.}  Read which arguments are
  traced vs.\ build-time (Section~\ref{sec:knobs}) and the ``build once, reuse the
  compiled closure'' pattern (Quickstart~2;
  Sections~\ref{sec:hybrid}, \ref{sec:linx}).
\item \textbf{Cross-reading the Fortran.}  The inline
  \textbf{[Fortran $\leftrightarrow$ JAX]} and \textbf{[Port note]} boxes map each
  piece to its BURST ancestor; the full enumeration is the mapping appendix
  (Appendix~\ref{sec:fortranmap}).
\item \textbf{Want the numerical reasoning first} (why the kernel tensors, the
  precision gate, the hybrid solver): read the numerical-methods background
  (Section~\ref{sec:numerics}) before the API.
\item \textbf{Hit a problem.}  See the troubleshooting quick reference
  (Section~\ref{sec:troubleshoot}) and the caveats
  (Section~\ref{sec:knownissues}).
\end{itemize}

The body proceeds physics (Section~\ref{sec:physics}) $\to$ numerics
(Section~\ref{sec:numerics}) $\to$ install and quickstarts
(Section~\ref{sec:install}) $\to$ a module-by-module reference
(Sections~\ref{sec:core}--\ref{sec:kernels}) $\to$ the coupled and hybrid solvers
(Sections~\ref{sec:solver}, \ref{sec:hybrid}) $\to$ the LINX coupling
(Section~\ref{sec:linx}) $\to$ parameters (Section~\ref{sec:knobs}), validation
(Section~\ref{sec:validation}), and caveats (Section~\ref{sec:knownissues}); then
the QKE flavour-mixing refinement (Section~\ref{sec:qke})
and directions for extending the code (Section~\ref{sec:extend}); then Part~II on the
interacting dark sectors (Sections~\ref{sec:dark}--\ref{sec:darkextend}); closing with
the mapping appendix, glossary, and references.

% ###########################################################################
\part{Standard-Model neutrino transport}
% ###########################################################################

% ===========================================================================
\section{Physics background}
\label{sec:physics}
% ===========================================================================

This section sets out the physics the code implements, at the level of the
BURST paper.  Readers fluent in weak decoupling can skim to the conventions in
Section~\ref{sec:notation} and the collision integral in
Section~\ref{sec:collint}.

\subsection{Weak decoupling and neutrino spectral distortions}

In the early universe ($T\gtrsim$ few MeV) neutrinos are kept in equilibrium
with the $e^\pm/\gamma$ plasma by charged- and neutral-current weak
interactions, which have rate $\Gamma\sim\GF^2 T^5$.  As the universe expands,
$\Gamma/H\propto T^3$ ($H$ the Hubble expansion rate) falls below unity near
$T\simeq2\text{--}3$ MeV and the
neutrinos decouple.  In the idealized ``instantaneous decoupling''
picture, decoupling happens before $e^\pm$ annihilation, the neutrinos retain a
perfect Fermi--Dirac spectrum, and the subsequent $e^+e^-\to\gamma\gamma$
heats only the photons, giving $\Tg/\Tnu=(11/4)^{1/3}$ and exactly $\Neff=3$.

Decoupling is not instantaneous.  At the few-MeV epoch where the weak rates are
shutting off, $e^\pm$ annihilation has already begun, so a small amount
of the entropy released by annihilating pairs leaks into the
still-marginally-coupled neutrinos before the channel closes.  Because the weak
cross-section grows with energy, this leakage preferentially heats the
high-energy tail of the neutrino spectrum: the distribution develops a
non-thermal \textbf{spectral distortion} $\delta f(\eps)$, positive and growing
with $\eps$.  The extra energy carried by the distorted seas raises the
radiation density slightly above the instantaneous-decoupling value, i.e.\
$\Neff>3$.  Capturing $\delta f(\eps)$ (not just an averaged $\Tnu$) is
exactly what BURST and \texttt{neutt} are for.

\begin{fmap}
The distortion is what the Fortran writes to its per-species output units
\code{nu_72.dat}\,\ldots\,\code{nu_75.dat}.  In \texttt{neutt} the same
information is the trajectory array \code{f_traj} returned in memory by
the \code{solve} of \code{make_decoupling_solver} (Section~\ref{sec:evolve}); the port writes no files (Section~\ref{sec:io}).
\end{fmap}

\subsection{The Boltzmann transport equation and the collision integral}
\label{sec:collint}

Each neutrino occupation evolves under
\begin{equation}
  \Big[\frac{\partial}{\partial t}-H(a)\,p\,\frac{\partial}{\partial p}\Big]
  f_{\nu_i}(p,t)=C_{\nu_i}\big[\{f_j\}\big],
  \label{eq:boltz}
\end{equation}
where $H$ is the Hubble rate and $C_{\nu_i}$ is the collision integral summing
all $2\!\to\!2$ weak processes the species participates in.  In the comoving
variable $\eps=p/\Tcm$ the free-streaming (redshift) term on the left is
absorbed: a decoupled neutrino sits at fixed $\eps$, so in $\eps$-space
Eq.~\eqref{eq:boltz} becomes simply $\dd f/\dd t=C$ (and the code integrates in
$\ln a$, dividing $C$ by $H$).

For a representative process $r$, $1+2\to3+4$, with four-momenta
$P_k=(E_k,\mathbf p_k)$ of the legs $k=1..4$, the collision integral is
\begin{equation}
  C^{(r)}_{\nu_1}=\frac{1}{2E_1}\int\!\frac{\dd^3p_2}{(2\pi)^3 2E_2}
  \frac{\dd^3p_3}{(2\pi)^3 2E_3}\frac{\dd^3p_4}{(2\pi)^3 2E_4}
  (2\pi)^4\delta^4(P_1{+}P_2{-}P_3{-}P_4)\,S_r\,\langle|M_r|^2\rangle\,F_r,
  \label{eq:collint}
\end{equation}
with the Pauli-blocked \textbf{gain $-$ loss} phase-space factor
\begin{equation}
  F_r=\underbrace{(1-f_1)(1-f_2)f_3 f_4}_{\rm gain}
     -\underbrace{f_1 f_2(1-f_3)(1-f_4)}_{\rm loss}.
  \label{eq:Fr}
\end{equation}
$\langle|M_r|^2\rangle$ is the spin-summed squared matrix element, $S_r$ a
symmetry factor.  At equilibrium ($f_i=f^{\rm eq}$) detailed balance makes
$F_r=0$ identically, so $C=0$: the equilibrium spectrum is a fixed point.

The processes are:
\begin{itemize}[nosep]
\item $\nu\nu\to\nu\nu$ and $\bar\nu\bar\nu\to\bar\nu\bar\nu$ elastic scattering
  (neutral current among neutrinos);
\item $\nu\bar\nu\to\nu\bar\nu$ scattering and the flavour-changing annihilation
  $\nu_i\bar\nu_i\to\nu_j\bar\nu_j$;
\item $\nu+e^\pm\to\nu+e^\pm$ scattering (two angular reductions, called
  R1 and R2 in BURST);
\item $\nu\bar\nu\to e^-e^+$ annihilation (and its inverse).
\end{itemize}
The first two families exchange energy only among the neutrino seas; the
$\nu$--$e^\pm$ families exchange energy with the plasma and therefore drive the
distortion.  This split is reflected in the code's \code{leg4} bookkeeping, BURST's
$e^\pm$ partner slot, not a momentum leg (Section~\ref{sec:assemble}).

The angular ($\delta^4$) integrals in Eq.~\eqref{eq:collint} can be done
analytically, reducing each process to a two-dimensional integral over outgoing
energies for every output bin $\eps_1$.  Energy conservation fixes the fourth
leg ($E_4=E_1+E_2-E_3$, i.e.\ $p_4=p_1+p_2-p_3$ for the relativistic legs),
which generally lands off the $\eps$-grid and must be interpolated.  How
the port turns this structure into a GPU tensor contraction is the subject of
Section~\ref{sec:kerneltensor}.

\subsection{The two-temperature description and $\Neff$}
\label{sec:neff}

After $e^\pm$ annihilation completes, the neutrino energy density is frozen in
comoving units and the photon temperature has been boosted.  $\Neff$ packages
the total neutrino (plus any extra-radiation) energy density relative to
photons through Eq.~\eqref{eq:neffdef}, equivalently
\begin{equation}
  \Neff=\frac{8}{7}\Big(\frac{11}{4}\Big)^{4/3}\frac{\rhonu}{\rhog},
  \qquad
  \rhonu=\frac{\Tcm^4}{2\pi^2}\sum_{i\in\{\rm flavour\},\,\rm cp}
        \int_0^{\eps_{\max}}\!\eps^3 f_i(\eps)\,\dd\eps ,
  \qquad
  \rhog=\frac{\pi^2}{15}\Tpl^4 .
  \label{eq:neffspec}
\end{equation}
The sum runs over the three flavours and both charge-parity (particle /
antiparticle) states; each contributes $g=1$.  Instantaneous decoupling gives
exactly $3$; the spectral distortions and finite-$T$ QED corrections push it to
$\approx3.044$ (LO), or $3.0433$ with the NLO QED term that is on by default
(Section~\ref{sec:nloqed}).

\begin{fmap}
Eq.~\eqref{eq:neffspec} is \code{decoupling.neff_from_spectrum(grid, occ, T_cm,
T_pl)} (the second argument \code{occ} is the spectrum $f[3,2,N]$: 3 flavours
$\times$ 2 cp states $\times$ $N$ grid bins; and
\code{rho_nu_from_spectrum(grid, occ, T_cm)}); the prefactor $1/2\pi^2$ matches
BURST's \code{drho} moment normalization.  In Fortran the analogous number is
written to unit~82 (\code{params_82.dat}).
\end{fmap}

\paragraph{Initial vs.\ end-point evaluation}
LINX (and Escudero-style fluid codes) report a $\dNeff$ defined at the
initial time, $\dNeff^{\rm init}=\frac{8}{7}(\frac{11}{4})^{4/3}
\rho_{\rm ext}^{\rm init}/\rho_\gamma^{\rm init}$ (with $\rho_{\rm ext}$ the
extra-radiation energy density), distinct from the late-time
$\Neff$; the
coupling layer (Section~\ref{sec:linx}) is careful about which is which, and
the public $\dNeff$ knob uses LINX's initial-time convention.

\subsection{Finite-temperature QED corrections}
\label{sec:qedphys}

At $\mathcal{O}(\alpha)$ the surrounding $e^\pm/\gamma$ plasma shifts the
electron and photon dispersion relations (thermal masses $\delta m_e^2$,
$\delta m_\gamma^2$) and therefore the plasma equation of state $\rho_{\rm pl}$,
$P_{\rm pl}$, $\dd\rho_{\rm pl}/\dd T$.  Because $\Neff\propto\Tpl^{-4}$ at fixed
$\rhonu$, this $\sim$per-cent correction to the plasma energy budget moves
$\Neff$ at the third decimal; it is the difference between the
``decoupling-only'' $3.034$ and the LO-QED $3.044$ (a further NLO term, on by default,
trims it to $3.0433$; Section~\ref{sec:nloqed}).  BURST computes the LO
correction
as an additive $\delta\rho/\delta P$ built from the thermal self-energies
$\delta m_e^2$, $\delta m_\gamma^2$ (which enter BURST's seven Gauss--Laguerre
plasma-moment integrals \code{glag1}\,\ldots\,\code{glag7}, the $e^\pm$/$\gamma$
energy density, pressure, $\dd\rho/\dd T$, and their finite-$T$ QED corrections),
not as a mass shift inside the free-gas integrand; the
port follows this faithfully
(Section~\ref{sec:qedeos}).

\subsection{BBN observables and the $n\leftrightarrow p$ rates}
\label{sec:bbnphys}

The light-element abundances ($\Yp$, $\DoverH$, $^3$He/H, $^7$Li/H) are set by the
neutron-to-proton ratio at freeze-out and the subsequent nuclear network.  The
$n\leftrightarrow p$ interconversion rates ($n+\nu_e\leftrightarrow p+e^-$,
$n+e^+\leftrightarrow p+\bar\nu_e$, and neutron decay) depend on the
electron-neutrino occupation at physical energy, precisely the
distorted $f_{\nu_e}(\eps)$ that transport produces.  A distortion in
$f_{\nu_e}$ thus feeds directly into $\Yp$ and $\DoverH$.  In the \texttt{neutt}
$+$ LINX coupling, this is the role of \code{f_e_occupation}
(Section~\ref{sec:linx}): it evaluates the transport-distorted $\nu_e$ spectrum
at the energies LINX's weak-rate integrals request.

\begin{diffbox}
Fortran BURST contains its own nuclear reaction network and computes
$\Yp$, $\DoverH$ internally.  The JAX port does not re-implement the network;
it ports only the neutrino transport and hands the distorted spectrum to LINX's
abundance model.  See Appendix~\ref{sec:fortranmap}.
\end{diffbox}

% ===========================================================================
\section{Numerical methods background}
\label{sec:numerics}
% ===========================================================================

The port reproduces BURST's physics but reorganizes its numerics for the GPU
and for differentiability.  This section explains the non-trivial numerical
ideas at a level the reader will want before reading the API.

\subsection{The kernel-tensor reformulation}
\label{sec:kerneltensor}

The collision integral~\eqref{eq:collint}, after the analytic angular
reduction, has the schematic form (illustrated for $\nu$--$\nu$)
\begin{equation}
  \mathrm{net}[i]=\mathrm{prefac}\sum_j \mathrm{kinem}(i,j)
  \sum_{p_3}\mathrm{glw}(i,j,p_3)\,\mathrm{w_3}(i,j,p_3)\,
  F\big(f_1,f_2,f_3,f_4\big),
  \label{eq:net}
\end{equation}
where $f_1=f(\eps_i)$, $f_2=f(\eps_j)$ are on-grid (so $p_1\equiv\eps_i$,
$p_2\equiv\eps_j$), $f_3=f(p_3)$ and $f_4=f(p_1{+}p_2{-}p_3)$ are off-grid, and
$F$ is the gain$\,-\,$loss tensor of Eq.~\eqref{eq:Fr}.  The schematic symbols
($\mathrm{prefac}$ the process prefactor, $\mathrm{kinem}$ the kinematic weight,
$\mathrm{glw}$ the Gauss--Legendre weights, $\mathrm{w_3}$ the leg-3 quadrature
measure) stand for the host-static factors enumerated below.  The key
observation (anticipated by the Fortran authors
in \code{trans_init.f90}) is that every factor except $F$ depends only
on the $\eps$-grid and the quadrature nodes, never on the occupation numbers.
The prefactor, kinematic weights, Gauss weights, segment boundaries, and node
positions are all \textbf{host-static}, computable once.

Moreover the off-grid reads are done by interpolating $\log f$ with a
fixed-order polynomial, which is linear in the sampled $\log f$ values:
$\log f(p)=\sum_a L[p,a]\,\log f_a$ over a 6-point window, with static Lagrange
weights $L$.  So the only nonlinear step is the final pointwise $\exp$.  The
runtime therefore collapses to:
\begin{enumerate}[nosep]
\item gather $\log f$ through two static sparse interpolation tensors to get
  $f_3,f_4$ (a \code{jnp.take} plus weighted sum, then $\exp$);
\item form $F$ pointwise (multilinear in $f_1,f_2,f_3,f_4$);
\item contract $F$ against one static weight tensor $W$ with an \code{einsum}.
\end{enumerate}
The depth-3 Fortran loops over (bin $\times$ partner-bin $\times$ node) become a
single batched contraction over the (flavour $\times$ cp $\times$ bin) axis,
exactly the axis the Fortran's MPI (Message Passing Interface) build already
splits across ranks.  This is the
central GPU idea and the reason the port can be both faithful and fast.

\paragraph{Static objects (grid-only kernels)}
With $N$ bins, $n_{\rm GL}=5$ Gauss nodes, $M=6$ interpolation points, and $K$
quadrature slots per bin-pair, the precomputed tensors of Eq.~\eqref{eq:net}
are the weight tensor $W$ of shape $(N,N,K)$ and the interpolation index/weight
tensors
$(\mathrm{idx}_3,\mathrm{lw}_3,\mathrm{idx}_4,\mathrm{lw}_4)$ of shape
$(N,N,K,M)$.  Sparse 6-point windows keep each interpolation tensor at $\sim$7~MB;
a dense $(N,N,K,N)$ matrix would be $\sim$120~MB.

\paragraph{Kinematic case-branches $\to$ union grid $+$ masks}
The energy integral runs over segments split at kinematic boundaries, and the
$\nu$--$e^\pm$ kernels further split the bin range into ``regimes'' with
different segment structure.  Rather than data-dependent control flow (illegal
under \code{vmap}/\code{jit}), every segment/regime is materialized into a fixed
slot/tile dimension and the inactive ones are masked, baked to $W=0$ for the
grid-only kernels, and selected by runtime regime masks for the lepton kernels.

\begin{diffbox}
This ``union grid $+$ mask'' is the port's structural departure from the
Fortran, which uses ordinary \code{if}-branches on the kinematic regime.  The
two agree in value; the JAX version simply evaluates a fixed superset and
zeroes what the Fortran would have skipped.  A faster variant
(Section~\ref{sec:sched}) gathers each bin's active tiles instead of masking.
\end{diffbox}

\subsection{Quadrature: what rule is used where}
\label{sec:quadnum}

Three quadrature families appear, exactly as in the BURST paper:
\begin{itemize}[nosep]
\item \textbf{Gauss--Legendre} for inner energy integrals over finite intervals;
\item \textbf{Gauss--Laguerre} for semi-infinite $[a,\infty)$ tails (the
  exponentially-decaying Boltzmann factors).  The port uses the plain variant
  ($\alpha=0$ in the generalized weight $x^\alpha e^{-x}$; this $\alpha$ is not the
  fine-structure constant) with $e^{x}$ folded into the weights so the rule
  integrates a bare $\int_0^\infty f\,\dd x$ (\code{grid.gauss_laguerre_plain});
  the genuinely generalized ($\alpha\ne0$) routine exists but is unused;
\item \textbf{composite Boole's rule} (5-point Newton--Cotes) for the outer
  integral over the neutrino energy grid (this is why $N$ must be $4n{+}1$ for
  integer $n$).
\end{itemize}
The classical Gauss rules are tiny and fixed, so the port precomputes
their nodes and weights on the host with NumPy/SciPy and carries them as
constants.

\begin{diffbox}
This is the one deliberate algorithmic substitution in the auxiliaries: BURST
generates Gauss nodes at runtime via the Golub--Welsch eigensolve
(\code{gauss.f90:imtqlx}).  The port does not port the eigensolve; it
precomputes nodes with \code{numpy.polynomial.legendre.leggauss} and
\code{scipy.special.roots_genlaguerre}, and ports only the affine interval
rescale \code{scqf}, which the kernels apply inline to the canonical $[-1,1]$ nodes
(\code{grid.rescale_legendre} is the standalone form), because the
collision-integral limits move with the running momentum and must stay jittable.
\end{diffbox}

\subsection{Off-grid interpolation in log space}
\label{sec:interpnum}

The off-grid legs $f_3,f_4$ are read by BURST's \code{interp_occfrac}: a
5th-order (6-point) polynomial interpolation of $\log f$, clamped to avoid
$\log$ of a non-positive number.  Working in $\log f$ matches the
near-exponential shape of the Fermi--Dirac tail.  The port reproduces this two
ways:
\begin{itemize}[nosep]
\item for static query points (the grid-only kernels, whose $p_3,p_4$
  positions are fixed) the 6-point window indices and Lagrange weights are
  precomputed on the host (\code{interp.lagrange_window}) and the runtime is a
  gather $+$ weighted sum $+$ clamp (\code{interp.apply_log_interp}, which
  returns clamped $\log f$), followed by the pointwise $\exp$ in the kernel;
\item for dynamic query points (the lepton kernels, whose nodes move with
  $\me/\Tcm$) the window is found at runtime with \code{searchsorted}
  (\code{interp.interp_log_occ_dynamic}).
\end{itemize}

\begin{diffbox}
The high-occupation clamp follows the Fortran reference (\code{transvar.f90}): an
interpolation overshoot to $\log f>0$ is sent to the floor ($\log f=-300$, i.e.\
$f\to0$) rather than clamped to $f=1$.  This matters only for unphysical
overshoots; on the smooth SM spectrum it never fires.
\end{diffbox}

\begin{diffbox}
\textbf{Port-only off-grid closures.}  \code{interp.py} carries three optional
off-grid tail closures, the Fermi--Dirac envelope cap
$f(p)\le C\cdot f^{\rm eq}(p)$ (\code{fd_cap_logC}$\,=\ln C$), \code{fd_decay} and
\code{monotone_read}, used by the flavour-mixing (QKE) solve to tame a tail
amplification (Section~\ref{sec:qkecap}).  All three default off
(\code{fd_cap_logC}$=+\infty$ is a bitwise no-op), so the classical SM/$\dNeff$
solve does not use them.
\end{diffbox}

\subsection{The net/frs decomposition and the precision gate}
\label{sec:gatenum}

The gain$\,-\,$loss factor~\eqref{eq:Fr} is a catastrophic cancellation near
equilibrium: gain and loss are each $\mathcal O(1)$ but their difference is
$\mathcal O(10^{-9})$ or smaller.  BURST tracks both the signed difference and
the cancellation-free sum of the forward (gain) and reverse (loss) rates:
\begin{equation}
  \mathrm{net}=\mathrm{forward}-\mathrm{reverse},\qquad
  \mathrm{frs}=\mathrm{forward}+\mathrm{reverse}.
\end{equation}
\code{net} is the physical rate; \code{frs} is the total throughput, used to
normalize the cancellation.  Because the off-grid interpolation is only
$\mathcal O(10^{-9}\text{--}10^{-13})$ accurate (a degree-5 polynomial is not
exact for $\log$ of a Fermi--Dirac), the raw \code{net} near equilibrium is
dominated by interpolation noise.  BURST suppresses it with a \textbf{precision
gate}: for each bin it forms a noise floor \code{prec} (the value of
$|{\rm net}|/{\rm frs}$ at exact equilibrium, where the true rate is zero) and
keeps a rate only if it stands out above the floor,
\begin{equation}
  \mathrm{test\_ratio}=\frac{|\mathrm{net}|}{\mathrm{frs}\cdot\mathrm{prec}}
  \ \ge\ \mathrm{nuratiotol}=3.0
  \ \Rightarrow\ \text{keep};\qquad \text{else set the rate to }0.
  \label{eq:gate}
\end{equation}
The floor (\code{prec_trans} in the code) is a one-time static precompute at the equilibrium
spectrum.  This gate is essential to reproduce BURST's $\Neff$; it is what stops
interpolation noise from random-walking the spectrum during the long
near-equilibrium phase.

\begin{diffbox}
\textbf{The \code{optimization_barrier}.}  The noise floor is computed from a
fixed (input-independent) equilibrium spectrum, so under full \code{jit} XLA
would recognize it as a compile-time constant and evaluate it through its
constant-folding path.  That path reorders/aggregates the large gain and
loss terms differently from the runtime kernel, and so would lose the $10^{-9}$
cancellation that is the floor, collapsing \code{prec}$\to1$, inverting
the gate, and shifting $\Neff$ (in the decoupling-only configuration) from
$3.034$ to $3.004$.  (A deterministic precompute can give two answers because
the two evaluation paths use different floating-point arithmetic, and a
near-cancellation is exactly where that matters.)  Wrapping the equilibrium
occupation in \code{jax.lax.optimization_barrier} marks it opaque, forcing XLA to
use the same runtime arithmetic in both places.  At runtime the barrier is an
identity, so the floor stays bit-faithful (Section~\ref{sec:knownissues}).
\end{diffbox}

\subsection{ODE integration: diffrax, Tsit5, Kvaerno5, PID control}
\label{sec:odenum}

The coupled state is the pair $\{\Tpl,\;f[3,2,N]\}$ (a JAX PyTree), integrated
in $\ln a$ by \texttt{diffrax}.  Two solvers are offered:
\begin{itemize}[nosep]
\item \textbf{Tsit5}: an explicit 5th-order Runge--Kutta (Tsitouras), the
  default.  Accurate and cheap per step, but stability-limited in the stiff
  high-$T$ phase, where the thermally averaged collision loss rate is $\Gamma/H\sim10^3$
  at 10~MeV (the stiffest bin's is $\approx8.6\times10^3$ per $e$-fold,
  Section~\ref{sec:hybridnum}).
\item \textbf{Kvaerno5}: an L-stable implicit ESDIRK (explicit-first-stage,
  singly diagonally implicit Runge--Kutta) method, an alternative
  for a stiff high-$T$ start.  Note the default explicit Tsit5 with
  $\code{atol}=10^{-30}$ already completes the $10$~MeV start (it does not
  blow up there); it is merely stability-limited (slow, not unstable) in
  the stiff head, which is precisely what the hybrid of
  Section~\ref{sec:hybridnum} cures.
\end{itemize}

\begin{diffbox}
BURST integrates in time $t$ with an adaptive Cash--Karp RK5(4) of
$\sim3\times10^4$ steps.  The port integrates in $\ln a$ (with
$\Tcm=\Tcm^{\rm init}e^{-\ln a}$ analytic) using diffrax's adaptive solvers and
a PID step controller.  The physics window is the same.
\end{diffbox}

Step size is set by a \textbf{PID controller} on the embedded error estimate,
the difference between the high- and low-order solutions the Runge--Kutta pair
produces for free at each step, which estimates the local error the controller
keeps near the requested tolerance.  The defaults are a damped PI law,
$(p_{\rm coeff},i_{\rm coeff},d_{\rm coeff})=
(0.4,0.3,0.0)$, with a relative tolerance $\code{rtol}=10^{-6}$ and an
absolute tolerance $\code{atol}=10^{-30}$.  Two non-obvious facts about these
defaults (both measured, both documented in Section~\ref{sec:knobs}):
\begin{itemize}[nosep]
\item $\code{rtol}=10^{-6}$ is faster than the looser $10^{-4}$, because
  at $10^{-4}$ the explicit controller runs right at the stability edge and
  spends $\sim25\%$ of attempts on rejected steps;
\item $\code{atol}=10^{-30}$ makes the error norm effectively pure-relative,
  matching BURST.  Raising it toward $10^{-7}$ leaves the stiff $\bar\nu$-tail
  bins error-uncontrolled and silently biases $\Neff$ by $+1.5\times10^{-3}$
  (Section~\ref{sec:knownissues}).  Do not raise it.
\end{itemize}

These tolerances govern the classical (diagonal) transport ODE only.  The
QKE flavour-coherence sector (Section~\ref{sec:qkesplit}) is deliberately
not integrated this way: its off-diagonal substep is an exact matrix
exponential carrying no per-component tolerance (the stiff flavour
oscillation is exponentiated exactly rather than error-controlled), so
\code{rtol}/\code{atol} never act on the coherences, and the adaptive-integrator
failure mode (a tiny imaginary off-diagonal drifting against a loose, trace-scaled
tolerance) cannot arise.

\subsection{The exponential-midpoint hybrid solver}
\label{sec:hybridnum}

In the stiff $10\!\to\!2.2$~MeV ``head'' of the solve the collision loss rate is
very large (at the 10~MeV start the largest $B$ over all modes and bins is
$\approx8.6\times10^3$ per $e$-fold, where $B$ is the per-bin
relaxation coefficient $B=(g{+}l)/H$ of the affine right-hand side (RHS) below,
with $g$ and $l$ the bin's gated gain and loss rates per unit time).  The spectrum
therefore relaxes to the instantaneous
collision fixed point (the \textbf{quasi-static-equilibrium (QSS) attractor}
$f^\star=A/B$, the value at which the collision RHS $\dd f/\dd\ln a=A-Bf$ vanishes)
far faster than the universe expands, and sits on it to one part in
$10^{17}$.  An explicit method is stability-limited there, not accuracy-limited:
it burns $65$--$92\%$ of all steps fighting a stiffness that carries no
information (the attractor it is chasing barely moves).  The remedy exploits the
exact affine structure of the gated collision RHS,
\begin{equation}
  \frac{\dd f}{\dd\ln a}=A-B\,f,\qquad A=g/H\ge0,\quad B=(g+l)/H\ge0,
  \label{eq:affine}
\end{equation}
which is pure algebra on the kernels' forward/reverse rates (every forward rate
carries a Pauli factor $1-f_1$, every reverse carries $f_1$, so
$\mathrm{forward}=(1-f)g$, $\mathrm{reverse}=fl$).  An exponential integrator
treats the linear part exactly:
\begin{equation}
  f' = f\,e^{-Bh}+\frac{A}{B}\big(1-e^{-Bh}\big),
  \label{eq:expstep}
\end{equation}
for a step $h$ in $\ln a$: a convex combination of $f$ and the attractor $A/B$,
unconditionally
stable and positivity-preserving, with no step-size ceiling.  (Equation
\eqref{eq:expstep} is the 1-stage ``exponential-Euler'' building block; the head
uses its 2-stage midpoint form ($A$ and $B$ re-evaluated at the step
midpoint) to reach second order, while both treat the stiff $-Bf$ linear part
exactly.)  The \textbf{hybrid} runs this fixed-step exponential-midpoint
integrator over the stiff head (a \code{lax.scan} of $\sim1600$ steps), then
hands off to the production adaptive Tsit5 for the accuracy-demanding
annihilation tail, where stiffness is gone and order-5 earns its keep
($\sim96$ tail steps, zero rejections).

The handoff is \textbf{qualified}: at the switch $\ln a$ the solve checks that
the spectrum is on the attractor ($\max|f-A/B|\le\code{qual_tol}=10^{-4}$) and
sane (finite, no negative or super-unity bins).  If the check fails, it
silently falls back to the full adaptive solve from $\ln a=0$; both
modes share one compiled executable (the fallback is selected by a traced
\code{where} on the tail solve's initial time and state, $t_0,y_0$).  A mis-qualified configuration therefore costs speed, never
accuracy.  Measured: the hybrid's $\Neff$ matches the tight-tolerance
($\code{rtol}=10^{-8}$) arbiter to $-6.9\times10^{-8}$ (comfortably inside
the $|\dNeff|\le2\times10^{-6}$ acceptance bar), at $2.36\times$ the speed
of the plain transport solve.  The full
rationale and the no-go for a pure exponential solve (the order-2 energy
budget would need $\sim18{,}000$ steps; high-order exponential schemes go
unstable against the discontinuous precision gate) is in
Section~\ref{sec:hybrid}.

\subsection{float64, NaN guards, and differentiability}
\label{sec:float64}

float64 is \textbf{mandatory}: the collision operator is a forward$-$reverse
near-cancellation that float32 cannot resolve (the conservation sum rules fail).
The package enables it at import:
\begin{lstlisting}
import neutt   # runs jax.config.update("jax_enable_x64", True) on import
\end{lstlisting}
\textbf{Import \code{neutt} before creating any JAX array}, or x64 will not
take effect for arrays made earlier in the process.

Two robustness guards deserve mention because they are visible in the API:
\begin{itemize}[nosep]
\item \code{nanrobust_rms_norm} maps a NaN error estimate to $+\infty$ so
  diffrax's controller routes it into the designed reject-and-shrink path
  instead of freezing at $\dd t=\mathrm{NaN}$ (the diffrax NaN freeze,
  Section~\ref{sec:knownissues}); bit-identical on every finite step;
\item the kernels use double-\code{where} guards on \code{sqrt} and on masked
  denominators so a masked-off tile cannot emit a NaN that $0\cdot\mathrm{NaN}$
  would propagate.  These make every kernel \code{grad}-safe, a capability
  the Fortran does not have.
\end{itemize}

\subsection{The GPU-throughput strategy}
\label{sec:gpunum}

There are two axes of parallelism, used differently.

\paragraph{Within a call (automatic)}  A single \code{solve()} already runs the
whole computation on the GPU: the collision kernels are one batched tensor
contraction over the (flavour $\times$ cp $\times$ bin) axes
(Section~\ref{sec:kerneltensor}), so the parallelism lives inside each
solve and there is nothing for the user to arrange.

\paragraph{Across calls (batch a scan with \texttt{vmap})}  A single solve is
small and dispatch-bound: its wall time is dominated by the per-call GPU
launch overhead, not by arithmetic, so the way to use the GPU efficiently for a
scan is to \code{vmap} the built solver over a parameter array, so one
launch does many solves at once:
\begin{lstlisting}
import jax, jax.numpy as jnp
from neutt.hybrid import make_hybrid_solver
solve, _ = make_hybrid_solver()                       # build / compile ONCE
dneff = jnp.linspace(0.0, 0.5, 1024)                  # 1024 dark-radiation values
neffs = jax.vmap(lambda d: solve(delta_neff_init=d)["neff_final"])(dneff)  # ONE batched launch
\end{lstlisting}
The background ODE alone drops from $\sim$60~ms/solve to $\sim$96~$\mu$s/cosmology
this way; the full transport solve amortizes its per-call launch overhead the same
way, at a larger absolute cost.  (\code{pmap} instead spreads the batch across several GPUs; reach for it
only when a batch is too large for one device; a single solve is far too
small to shard usefully, Section~\ref{sec:extend}.)

\paragraph{Coupling to LINX: the device split, and collecting the outputs}
LINX's background and abundance models compile prohibitively slowly on the GPU, so
the coupled driver runs LINX on the CPU and the transport solve on the GPU
in one process (Section~\ref{sec:linx}).  For a single cosmology \code{run()}
handles the hand-off internally.  For a batch, the GPU results must be
\textbf{collected to the host before the CPU LINX side can use them}; you cannot
call a CPU model from inside a GPU \code{vmap}.  The pattern is: batch the transport
on the GPU, pull the whole array back to the host in one transfer (optionally save
it), then drive LINX over the stored spectra:
\begin{lstlisting}
import numpy as np, jax
specs   = jax.vmap(transport_solve)(params)   # batched on the GPU
specs_h = jax.device_get(specs)               # ONE GPU -> host transfer
np.savez("transport_batch.npz", **specs_h)    # optional: persist for a CPU-only LINX run
for s in specs_h:                             # LINX (CPU) consumes them one at a time
    abundances = linx_abundances(s)
\end{lstlisting}
Here \code{params} is the batched array of per-call inputs, \code{transport_solve}
the GPU spectrum solver and \code{linx_abundances} the CPU abundance step; the one rule is that the GPU batch is materialized to the
host (one transfer, or a saved file) before LINX reads it.

% ===========================================================================
\section{Installation, environment, and a first run}
\label{sec:install}
% ===========================================================================

\subsection{Where the code lives}

\texttt{neutt} is a pure-Python package; there is nothing to build (no compiled
extension).  It is installed from PyPI with \code{pip install neutt} (CPU) or
\code{pip install "neutt[cuda12]"} (NVIDIA GPU), or used from a checkout by putting the
repository root on \code{PYTHONPATH} (or \code{pip install -e .}).  The API reference is a
Sphinx site built from \code{docs/sphinx/} (\code{pip install "neutt[docs]"}, then
\code{sphinx-build docs/sphinx docs/sphinx/_build}).
Throughout this manual \texttt{\$BURST\_ROOT} denotes your checkout of the BURST
repository (the \texttt{neutt} package is \code{neutt/} inside it) and
\texttt{\$LINX\_ROOT} the LINX checkout used for the abundance coupling
(Section~\ref{sec:linx}).  A read-only Standard-Model Fortran reference (a
snapshot of the Fortran BURST code, against which the port is validated) ships in
\code{fortran_reference/}.

\paragraph{The package at a glance}
The modules layer from foundation to coupling (a reader's first map of the code):
\begin{center}\small
\begin{tabular}{@{}R{0.30\textwidth}R{0.62\textwidth}@{}}
\toprule
\textbf{layer / module} & \textbf{role} \\
\midrule
\code{neutt} (\code{__init__}) & enables float64 at import; the compile banner \code{compiling()} \\
\code{constants}, \code{grid}, \code{interp} & foundation: physical constants, the $\eps$-grid with its Gauss/Boole rules, equilibrium occupations and momentum moments, log-interpolation \\
\code{thermo} & the plasma: free-gas $e^\pm/\gamma$ EoS, the charge-neutral electron degeneracy $\phi_e=\mu_e/\Tpl$ ($\mu_e$ the electron chemical potential), finite-$T$ QED, Hubble / decoupling clock \\
\code{transport/} (\code{nunu}, \code{nunubar}, \code{nuer1}, \code{nuer2}, \code{nuepma}, \code{assemble}) & the SM weak collision kernels (each lepton kernel with its reference and scheduled implementation, Section~\ref{sec:sched}) and the gated RHS assembly \\
\code{decoupling}, \code{hybrid} & the coupled solvers: the production RHS, \code{make_decoupling_solver}, the exponential-midpoint hybrid \\
\code{transport/qke/} & the QKE density-matrix solver, the SM flavour-mixing refinement (Section~\ref{sec:qke}) \\
\code{dark}, \code{transport/} (\code{dv_nu}, \code{chi_nu}, \code{chi_brute}) & the dark sectors (Part~II): the sector data types and the decay and $\chi\leftrightarrow\nu$ kernels \\
LINX adapter (separate repo) & \code{transport_adapter} / \code{coupled_bbn}: the seam that hands the spectra to LINX for abundances \\
\bottomrule
\end{tabular}
\end{center}

\subsection{Environment (Perlmutter)}

Everything that touches JAX runs on a Perlmutter GPU compute node, never
the login node.  Request an interactive GPU allocation and activate the conda
environment:
\begin{lstlisting}[style=sh]
salloc --nodes=1 --qos=interactive --time=04:00:00 \
       --constraint=gpu --gpus=1 --account=<your-account>
module load conda && conda activate <your-jax-env>
export OMP_NUM_THREADS=1
export XLA_PYTHON_CLIENT_PREALLOCATE=false
export MPICH_GPU_SUPPORT_ENABLED=1
export BURST_ROOT=/path/to/BURST          # your BURST checkout
export LINX_ROOT=/path/to/LINX            # your LINX checkout (coupled BBN only)
export PYTHONPATH=$BURST_ROOT
\end{lstlisting}
For the coupled-BBN demos you additionally need the LINX checkout on the path
(Section~\ref{sec:linx}).  The Python environment must carry JAX with GPU support
and \texttt{diffrax}.

\subsection{Quickstart 1: standalone $\Neff$}

The most useful entry point.  This evolves the coupled (plasma $+$
neutrino-spectrum) system through $e^\pm$ annihilation and returns $\Neff$:
\begin{lstlisting}
# on a GPU node, with PYTHONPATH=$BURST_ROOT
from neutt.decoupling import make_decoupling_solver
solve, grid = make_decoupling_solver()   # builds the kernels and the noise floor
out = solve()                            # compiles on the first call, then solves
print(float(out["neff_final"]))          # -> 3.043322 (NLO QED default; 3.044276 with qed_nlo=False)
\end{lstlisting}
Building the solver and its first call pay a one-time kernel-build $+$ XLA-compile
cost (the slow ``cold''/``maiden'' call, $\sim$40--80 s), and the compile prints
the banner (a newt, ``neutt is compiling\dots'').  Later calls of the same
\code{solve} reuse the compiled executable (the fast ``warm'' call, no banner)
with new values of its per-call arguments; Section~\ref{sec:knobs} says which
arguments are which.  The result is the canonical Standard-Model number of
Section~\ref{sec:canonical}.

\begin{diffbox}
This single Python call replaces the transport sector of a $600$--$1000$
CPU-hour Fortran BURST run (the full Fortran run also does the nuclear network
and recombination; the transport collision integrals dominate its cost).  It
computes $\Neff$ and the spectra, but not abundances; for those,
couple to LINX (Quickstart~3).
\end{diffbox}

\subsection{Quickstart 2: a repeated transport solve}

For repeated solves (a scan, an MCMC) use the production hybrid solver, which is
$2.4\times$ faster than the adaptive \code{make_decoupling_solver}
(Section~\ref{sec:hybrid}).  Every solver has the same shape: a factory
\code{make_*_solver(...)} that builds the kernels and returns a jitted
\code{solve}, whose arguments are the per-call physics inputs:
\begin{lstlisting}
from neutt.hybrid import make_hybrid_solver
solve, grid = make_hybrid_solver()     # builds kernels; compiles on the first call
out = solve()                          # ~2.8 s warm
print(float(out["neff_final"]), bool(out["qualified"]))   # qualified=True: fast hybrid path used
out2 = solve(delta_neff_init=0.1)      # extra dark radiation, traced: later values reuse this compile
\end{lstlisting}
(\code{qualified=False} would mean the hybrid silently fell back to the full solve,
still correct, just slower; Section~\ref{sec:hybrid}.)

\subsection{Quickstart 3: coupled BBN (abundances $+$ $\Neff$)}

This couples transport to LINX.  It needs both repos on the path:
\begin{lstlisting}[style=sh]
export PYTHONPATH=$BURST_ROOT:$LINX_ROOT
\end{lstlisting}
\begin{lstlisting}
import neutt                                   # FIRST: enables x64
from linx.coupled_bbn import make_coupled_bbn

run, info = make_coupled_bbn(max_steps=8192)   # build ONCE, reuse for a whole chain
res, neff = run(delta_neff_init=0.0)                    # ~80 s maiden (cold), ~4.2 s warm
Yp = float(res[5] * 4.0)                           # res[5] = Y(He4); Yp = mass fraction
DH = float(res[2] / res[1])                        # res[2]/res[1] = Y(D)/Y(p) = D/H
print(neff, Yp, DH)        # -> N_eff 3.043322 (NLO); Yp ~0.2470, D/H ~2.44e-5 (network-dependent)
\end{lstlisting}
\code{res} is LINX's \code{AbundanceModel} yield vector, indexed
$(0{:}n,\ 1{:}p,\ 2{:}\mathrm D,\ 3{:}\mathrm T,\ 4{:}{}^3\mathrm{He},\ 5{:}{}^4\mathrm{He},\dots)$
as number abundances $Y_i=n_i/n_b$ ($n_b$ the baryon number density); the helium
mass fraction is
$\Yp=4\,Y({}^4\mathrm{He})$ and $\DoverH=Y(\mathrm D)/Y(p)$.

\subsection{Building the Fortran oracle (optional, for cross-checks)}

The Standard-Model Fortran reference (\code{fortran_reference/}) builds with
\code{gfortran} and is what the port is validated against.  You only need to
build it to produce a new reference output: the committed snapshot outputs already
cover the standard checks and are loaded by the test suite
(Section~\ref{sec:validation}).

% ===========================================================================
\section{Core module reference}
\label{sec:core}
% ===========================================================================

This section documents the four ``core'' modules, \code{constants},
\code{grid} (the grid, its quadrature, the occupation numbers, and the moments),
\code{interp} (the interpolation), and \code{thermo} (the background plasma
thermodynamics), everything the transport kernels
(Section~\ref{sec:kernels}) and the decoupling solver (Section~\ref{sec:solver})
rest on.  Each subsection gives the purpose, the public API and the Fortran origin,
and most add a short runnable demo.  A complete routine-by-routine mapping index is
in Appendix~\ref{sec:fortranmap}.

\subsection{\texttt{constants}: physical constants}
\label{sec:constants}

\paragraph{Purpose}
Module-level physical and mathematical constants, ported bit-for-bit from BURST
(except $\sin^2\theta_W$ and the last digits of $\zeta(3)$, below) so that parity checks are never polluted by a
constant mismatch.  Pure values, no functions.

The Standard-Model--relevant constants:

\begin{center}\small
\begin{tabular}{@{}llp{0.45\textwidth}@{}}
\toprule
\textbf{symbol} & \textbf{value} & \textbf{meaning / Fortran source} \\
\midrule
\code{M_ELECTRON} & $0.510998928$ & electron mass [MeV] (\code{xmelec}, \code{mainvar.f90}) \\
\code{GFHBAR}     & $0.20668190937$ & $\GF^2/\hbar$ [s$^{-1}$MeV$^{-5}$], weak collision-rate scale (\code{transvar.f90}) \\
\code{SIN2THW}    & $0.23122$ & $\sin^2\theta_W$ ($\overline{\rm MS}$ scheme at the $Z$ mass $M_Z$); enters the $\nu$--lepton kernels (and their QKE matrix counterparts) only. Fortran BURST uses $0.23$ (\code{SIN2THW_BURST}), which the Fortran-parity tests pin \\
\code{ALPHA_FS}   & $7.2973525698\!\times\!10^{-3}$ & fine-structure $\alpha$ \\
\code{MPL}        & $1.220932\!\times\!10^{22}$ & Planck mass [MeV] \\
\code{HBAR}       & $6.582119514\!\times\!10^{-22}$ & $\hbar$ [MeV$\cdot$s] \\
\code{PI}       & $3.1415926535897932$ & $\pi$ \\
\code{ZETA3}      & $1.202056903159594$ & $\zeta(3)$ \\
\code{BR_COEFFS}  & $(64,24,64,28)/45$ & interior Boole's-rule weights (\code{brcoeffs}, \code{mainvar.f90}) \\
\code{BR_END}     & $14/45$ & Boole's-rule end-point weight (\code{brends}) \\
\bottomrule
\end{tabular}
\end{center}

\begin{fmap}
Values mirror the snapshot's \code{mainvar.f90} (\code{pi}, \code{mpl}, \code{hbar},
\code{zeta3}, \code{xmelec}, \code{alphafs} at lines 25--37; \code{brcoeffs}, \code{brends} at
102--103), verified exact except \code{ZETA3}: \code{mainvar.f90:34} truncates $\zeta(3)$ to
$1.2020569031$, while the port carries it to double precision (it enters only the baryon density
behind $\phi_e$ and the QKE $n_p$).  \code{GFHBAR} and \code{SIN2THW_BURST} mirror
\code{transvar.f90:26--27} (\code{gfhbar}, \code{sin2thw}; production \code{SIN2THW} is the more
precise $0.23122$).
\end{fmap}

\begin{diffbox}
The module also defines QKE flavour-mixing constants (\code{GFERMI}$=\GF$, the
weak-potential scales \code{GFW2HBAR} $=\GF/m_W^2\hbar$ and \code{GFZ2HBAR}
$=\GF/m_Z^2\hbar$ with $m_W$, $m_Z$ the $W$ and $Z$ boson masses, and the PMNS
(Pontecorvo--Maki--Nakagawa--Sakata) mixing angles and mass-squared splittings
$\Delta m^2_{21}$, $\Delta m^2_{31}$).  These feed the QKE
flavour-mixing sector (Section~\ref{sec:qke}); the classical SM transport never
reads them.
\end{diffbox}

\subsection{\texttt{grid}: the comoving-energy grid, quadrature, occupations, and moments}
\label{sec:grid}

The module holds everything that lives on, or integrates over, the momentum
grid: the grid itself, the Gauss rules of the inner collision and plasma
integrals, the equilibrium occupations, and the Boole-weighted moments.

\subsubsection{The $\eps$-grid}

\paragraph{Purpose}
Defines the linear $\eps=p/\Tcm$ grid on which every distribution is
discretized, with baked-in composite Boole's-rule weights so a momentum integral
is just a weighted sum $\int g\,\dd\eps\approx\sum_i w_i\,g(\eps_i)$.

\paragraph{API}
\begin{lstlisting}
class EpsGrid(eqx.Module):
    abscissas : (nbins,)   # eps_i
    weights   : (nbins,)   # composite-Boole w_i ; sum(weights) == eps_range
    nbins, eps_range, eps_offset            # static

make_eps_grid(nbins=101, eps_range=25.0, eps_offset=0.0,
              bin1_weight_frac=0.0) -> EpsGrid
\end{lstlisting}
The grid has $N=$\code{nbins} points on $[\eps_{\rm offset},\eps_{\rm offset}+\eps_{\rm range}]$
(\code{eps_offset}, \code{eps_range}); \code{bin1_weight_frac} adds
$\eps_{\rm offset}\cdot$\code{bin1_weight_frac} to the first weight (BURST's offset term).
The defaults ($101$ bins, $\eps\in[0,25]$, no offset) are BURST's.  \code{nbins} must be
$4n{+}1$ (Boole requirement) or a \code{ValueError} is raised.  The spacing is
$h=\eps_{\rm range}/(N-1)$; interior weights cycle through \code{BR_COEFFS}$\times h$
and the endpoints carry \code{BR_END}$\times h$.

\begin{fmap}
\code{make_eps_grid} $\equiv$ \code{part_lib.f90::part_make_bins}; defaults from
the \code{part_bin_var} namelist (\code{nbins=101}, \code{epsrange=25}).
Bit-faithful transcription.
\end{fmap}

\paragraph{Demo}
\begin{lstlisting}
from neutt.grid import make_eps_grid
import numpy as np
g = make_eps_grid()
print(g.nbins, float(g.abscissas[-1]), float(np.sum(g.weights)))  # 101 25.0 25.0
\end{lstlisting}

\subsubsection{Equilibrium occupations}
\label{sec:occupation}

\paragraph{Purpose}
Equilibrium occupation numbers on the grid and the log-space clamp used by
interpolation.  A species in equilibrium at temperature $T_s$ and degeneracy
$\xi=\mu/T_s$ has $f(\eps)=1/(e^{\eps/r-\xi}\pm1)$ ($+$ for fermions, $-$ for bosons)
with $r=T_s/\Tcm$ (the \code{escale_ratio} argument).

\paragraph{API}
\begin{lstlisting}
fermi_dirac(eps, escale_ratio=1.0, xi=0.0)   # 1/(exp(eps/r - xi)+1)   (fermions)
bose_einstein(eps, escale_ratio=1.0, xi=0.0) # 1/(exp(arg)-1), guarded  (bosons)
log_occupation(occ)                          # clip(log(occ), -300, 0)
LOG_OCC_FLOOR = -300.0
\end{lstlisting}
\code{fermi_dirac}, the Fermi--Dirac (FD) occupation, is overflow-safe in float64;
\code{bose_einstein} returns $0$ outside $\mathrm{arg}=\eps/r-\xi\in(0,300]$ (avoiding the
$\eps\to0$ divergence) and is a
reference massless-boson occupation; the SM neutrinos use \code{fermi_dirac}.

\begin{fmap}
Maps to Fortran \code{fd_equil_calc} / \code{be_equil_calc}; the $\log$ clamp
\code{LOG_OCC_FLOOR}$=-300$ is the floor of \code{interp_occfrac} (\code{transvar.f90})
(the same value as \code{interp}'s \code{LOG_OCC_CLAMP}, Section~\ref{sec:interp}).
\end{fmap}

\subsubsection{Gauss nodes and rescaling}
\label{sec:quadrature}

\paragraph{Purpose}
Gauss--Legendre and (generalized) Gauss--Laguerre nodes/weights for the
collision-integral inner integrals and the thermo/QED momentum integrals.
Classical rules are precomputed on the host; only the affine rescale is jittable.

\paragraph{API}
\begin{lstlisting}
legendre_canonical(n)            -> (x, w)   # Gauss-Legendre on [-1,1]
gauss_laguerre(n, alpha=0.0)     -> (x, w)   # generalized GL: int x^a e^-x f
gauss_laguerre_plain(n)          -> (x, w)   # weights * exp(x): int_0^inf f  (workhorse)
rescale_legendre(x, w, a, b)     -> (x', w') # affine [-1,1] -> [a,b]  (jittable)
gauss_legendre(n, a=-1.0, b=1.0) -> (x, w)   # canonical, rescaled to static [a,b]
\end{lstlisting}
\code{gauss_laguerre_plain} ($e^{x}$ folded into the weights so
$\sum w_i f(x_i)\approx\int_0^\infty f$) is the rule used by \code{thermo}
(free-gas EoS, $\phi_e$, finite-$T$ QED); it is BURST's \code{cgqf} with rule kind 5
(generalized Laguerre) followed by \code{wts *= exp(absc)} (the node count is a
separate argument: $25$ for the $\phi_e$ and QED integrals, $100$ for the
free-gas $e^\pm$ integrals).
\code{rescale_legendre} is the \code{scqf} equivalent
and is pure arithmetic, so it is \code{jit}/\code{vmap}-friendly; \code{gauss_legendre}
uses it for a static interval, and the kernels apply the same affine map inline to
\code{legendre_canonical}'s nodes, broadcast over their moving collision-integral limits.

\begin{diffbox}
The Golub--Welsch eigensolve (\code{gauss.f90:imtqlx}) is \textbf{not ported}:
nodes come from NumPy/SciPy on the host, only the rescale is ported.  See
Section~\ref{sec:quadnum}.
\end{diffbox}

\subsubsection{Momentum-space moments}
\label{sec:moments}

\paragraph{Purpose}
Boole's-rule weighted moment integrals over an \code{EpsGrid}; the building
blocks for number/energy densities and the conservation sum rules.

\paragraph{API}
\begin{lstlisting}
integrate(grid, fvals)               # sum_i w_i f_i                = int f deps
moment(grid, occ, k)                 # sum_i w_i eps_i^k occ_i      = int eps^k occ deps
number_density(grid, occ, g=1.0, t_cm=1.0)   # n = g/(2pi^2) T^3 int eps^2 occ deps
energy_density(grid, occ, g=1.0, t_cm=1.0)   # rho = g/(2pi^2) T^4 int eps^3 occ deps
\end{lstlisting}
All are vectorized and \code{jit}/\code{vmap}-friendly.  The $k=2$ moment gives
number, $k=3$ gives energy (and feeds $\Neff$).

\paragraph{Demo: a thermal cross-check}
\begin{lstlisting}
import neutt
from neutt.grid import make_eps_grid, energy_density, fermi_dirac
import numpy as np
grid = make_eps_grid()
f_eq = fermi_dirac(grid.abscissas)          # FD at T_cm
rho  = energy_density(grid, f_eq, g=2.0, t_cm=1.0)   # the keyword g=2.0 is the dof, not the grid
# analytic relativistic fermion: (7/8)(pi^2/30)*g*T^4, here T=1, g=2
print(float(rho), 7/8 * np.pi**2/30 * 2)    # agree to ~1e-10
\end{lstlisting}

\subsection{\texttt{interp}: log-space occupation interpolation}
\label{sec:interp}

\paragraph{Purpose}
The order-5 (6-point) log-space interpolation of the occupation, the port of
BURST's \code{interp_occfrac}.  Recast as a sparse linear (Lagrange) operator on
$\log f$ for static query points, with a fully-dynamic variant for the moving
lepton-kernel queries (Section~\ref{sec:interpnum}).

\paragraph{API}
\begin{lstlisting}
ORDPOLINT = 5                                  # 6-point window
LOG_OCC_CLAMP = -300.0                         # log-occupation floor (= grid.LOG_OCC_FLOOR)
FD_ENVELOPE_C = 10.0                           # the QKE cap VALUE (env NEUTT_FD_ENVELOPE_C overrides); default fd_cap_logC=inf leaves it unused

lagrange_window(abscissas, query, ord=5) -> (idx, weights)     # HOST precompute
apply_log_interp(logf_grid, idx, weights,
        range_limit=False, fd_decay_base=None,
        bracket=None)                             -> logf      # jittable gather+clamp
interp_log_occ_dynamic(logf_grid, abscissas, query, ord=5,
        clamp_high=True, smooth_floor=0.0, range_limit=False,
        fd_cap_logC=inf, off_grid_mode="default", fd_decay=False,
        monotone_read=False)                      -> logf      # jittable, traced query
\end{lstlisting}
For the SM kernels: the grid-only kernels use \code{lagrange_window} (host) $+$
\code{apply_log_interp} (runtime); the lepton kernels use
\code{interp_log_occ_dynamic}.  An overshoot $\log f>0$ is always sent to the
$-300$ floor (the fermion clamp); the off-grid closures (\code{fd_cap_logC},
\code{fd_decay}, \code{monotone_read} in the dynamic read, and their
\code{fd_decay_base}/\code{bracket} counterparts in the static one) default off,
the SM-faithful setting.  So do the dark-sector options \code{range_limit},
\code{clamp_high}, \code{smooth_floor} and \code{off_grid_mode}
(Section~\ref{sec:darkknobs}).

\begin{diffbox}
The high-occupation clamp sends an interpolation overshoot to $\log f>0$ down to
the floor ($f\to0$) rather than clamping it to $f=1$; it fires only on unphysical
overshoots and is bit-faithful to the validated Fortran reference
(Section~\ref{sec:interpnum}).  \code{fd_cap_logC} (the envelope
$f\le C\,f^{\rm eq}$ with $\log C=$\code{fd_cap_logC}), \code{fd_decay} (off grid, $\log f$
capped at the FD shape anchored to the grid-edge value) and \code{monotone_read} (the
interpolant clamped to its two bracketing nodes; off grid both are the edge node, so the read
is pinned to the edge value, which \code{fd_decay} then lowers to the FD tail) are the tail closures used
by the QKE solve (Section~\ref{sec:qkecap}); all default off, so they do nothing in the
classical SM solve.
\end{diffbox}

\subsection{\texttt{thermo}: plasma thermodynamics, finite-$T$ QED, and the background}
\label{sec:thermo}

The module holds the $e^\pm/\gamma$ plasma: its free-gas equation of state, the
finite-$T$ QED mass renormalization, the QED equation-of-state corrections (the
LO terms built from it, and the NLO term), the decoupling-background clock, and the
charge-neutral electron degeneracy $\phi_e\equiv\mu_e/\Tpl$ ($\mu_e$ the electron chemical
potential).  \code{solve_phie} solves charge neutrality $n_{e^-}-n_{e^+}=q_{\rm baryon}\,n_b$
($n_b$ the baryon number density, set by the present baryon-to-photon ratio $\eta_0$,
argument \code{eta0}; $q_{\rm baryon}$ the net charge per baryon, argument \code{q_baryon},
default $0.5$) by fixed Newton steps from the non-degenerate $\mathrm{asinh}$ seed; it is
what \code{phie="auto"} evaluates inside the transport RHS (Section~\ref{sec:knownissues}).
\begin{lstlisting}
solve_phie(tplmev, tcmpl, eta0, q_baryon=0.5, m_e=M_ELECTRON,
           nodes=None, weights=None, n_newton=6)        # phi_e; NaN below T_pl ~ m_e/300
phie_linear(tplmev, tcmpl, eta0, q_baryon=0.5, m_e=M_ELECTRON,
            nodes=None, weights=None)                    # BURST's linearized phi_e
net_electron_number_over_T3(phi, mt, nodes=None, weights=None)   # (n_e- - n_e+)/T_pl^3, mt = m_e/T_pl
nb_over_Tcm3(eta0)                                       # n_b/T_cm^3 = eta0 * N_GAMMA * 11/4
\end{lstlisting}
Here \code{tcmpl}$=\Tcm/\Tpl$, and \code{nodes}/\code{weights} override the default 25-node plain
Gauss--Laguerre rule.  \code{phie_linear} is the single linear step from $\phi_e=0$, a validation
reference that equals \code{solve_phie} to machine precision at high $T$ ($\phi_e\sim10^{-9}$).

\subsubsection{The free-gas equation of state}

\paragraph{Purpose}
The free-gas $e^\pm/\gamma$ plasma equation of state: photon and $e^\pm$ energy
densities, pressures, and $\dd\rho/\dd T$, by Gauss--Laguerre momentum integrals.
Uses the \textbf{bare} electron mass; the finite-$T$ QED correction is added
separately (Section~\ref{sec:qedeos}).

\paragraph{API}
\begin{lstlisting}
rho_gamma(T)                          # photon rho = (pi^2/15) T^4   (g=2)
pressure_gamma(T)                     # rho_gamma/3
rho_electron(T, phi=0.0, m=M_ELECTRON, nodes=None, weights=None)      # e- + e+
pressure_electron(T, phi=0.0, m=M_ELECTRON, nodes=None, weights=None)
rho_plasma(T, phi=0.0, m=M_ELECTRON)             # rho_gamma + rho_electron
pressure_plasma(T, phi=0.0, m=M_ELECTRON)
drho_plasma_dT(T, m=M_ELECTRON)                  # analytic drho/dT at phi=0
drho_plasma_dT_phi(T, phi, m=M_ELECTRON)         # at fixed phi != 0 (autodiff)
\end{lstlisting}
The argument \code{phi} is $\phi_e$; the standard background uses
$\phi=0$ (the baryon asymmetry $\sim10^{-10}$ is negligible).  The $e^\pm$
integrals default to a 100-node plain Gauss--Laguerre rule.

\begin{fmap}
Integrands are the leading (free-gas) terms of BURST's \code{glag4} (energy) and
\code{glag7} (pressure) in \code{bbn_therm.f90}.  $\dd\rho/\dd T$ uses the
analytic $f(1-f)$ form, cross-checked against \code{jax.grad}.
\end{fmap}

\begin{diffbox}
Two intentional deltas: (i) the default uses $100$ Gauss--Laguerre nodes vs.\
BURST's $25$ (accuracy margin; pass BURST's $25$-node rule via \code{nodes}/
\code{weights} to match BURST bit-for-bit); (ii) the integrand uses the bare
mass; folding a thermal mass shift in here would double-count and send
$\Neff$ the wrong way (Section~\ref{sec:qedeos}).
\end{diffbox}

\paragraph{Demo}
\begin{lstlisting}
import neutt
from neutt import thermo
T = 2.0    # MeV
print(float(thermo.rho_plasma(T)), float(thermo.drho_plasma_dT(T)))
\end{lstlisting}

\subsubsection{Finite-$T$ QED mass renormalization}
\label{sec:renorm}

\paragraph{Purpose}
The thermal self-energies $\delta m_e^2$, $\delta m_\gamma^2$ (momentum
independent) and the momentum-dependent electron self-energy.  These feed the
QED EoS correction; physically they encode that the plasma shifts the
effective electron and photon masses.

\paragraph{API}
\begin{lstlisting}
get_eg_mass_renorm(tplmev, phie, nodes=None, weights=None, m_e=M_ELECTRON)
    -> (dmeeps2, dmedteps, dmedt2, dmgeps2, dmgdteps, dmgdt2)   # 6 dimensionless
delta_me2(tplmev, phie=0.0, ...)        # dmeeps2 * T_pl^2   [MeV^2]
delta_mg2(tplmev, phie=0.0, ...)        # dmgeps2 * T_pl^2   [MeV^2]
m_e_eff(tplmev, phie=0.0, ...)          # sqrt(m_e^2 + delta_me2)   [MeV]
get_e_mass_pdep(tplmev, phie, peps, ...)
    -> (dmep1, dmep2, dmep3, dmep4, dmep5)   # momentum-dependent, shaped like peps
\end{lstlisting}
Here \code{tplmev}$=\Tpl$ in MeV and \code{phie}$=\phi_e$.  The six outputs of
\code{get_eg_mass_renorm} are exactly the Fortran out-args, in order:
\code{dmeeps2}$=\delta m_e^2/\Tpl^2$, then \code{dmedteps} and \code{dmedt2}, its first- and
second-$T$-derivative terms, then the same three for $\delta m_\gamma^2$; they combine the
$2\pi\alpha/3$ thermal-photon piece with the $e^\pm$ thermal integral $(2\alpha/\pi)\int\dots$.
\code{get_e_mass_pdep} does the momentum-dependent piece at \code{peps}$=p/\Tpl$ (outputs
$\delta m_e^2(p)/\Tpl^2$ and its $T$, $p$, $T^2$ and $pT$ derivative terms), integrating
over the loop momentum $y=k/\Tpl$ with a reflected-Legendre rule about the singularity at
$y=$\code{peps} plus a shifted Laguerre tail (the principal-value construction).

\begin{fmap}
\code{get_eg_mass_renorm} $\equiv$ \code{renorm.f90:29--114} (verified
line-for-line); \code{get_e_mass_pdep} $\equiv$ \code{renorm.f90:119--282}.  The
$25$-node default rule matches BURST's \code{glag_bins=25}.
\end{fmap}

\subsubsection{The finite-$T$ QED equation of state}
\label{sec:qedeos}

\paragraph{Purpose}
BURST's \code{glag1}--\code{glag7} correction terms: the additive
$\delta\rho$, $\delta P$, $\delta(\dd\rho/\dd T)$ that turn the bare-mass
free-gas EoS into the $\mathcal O(\alpha)$ QED plasma EoS.  Built from the
thermal self-energies of Section~\ref{sec:renorm}, not from a shifted mass in the
free-gas integrand.

\paragraph{API}
\begin{lstlisting}
qed_eos_corrections(tplmev, phie=0.0, ...)
    -> (d_rho_gamma, d_P_gamma, d_drho_gamma_dT,  d_rho_e, d_P_e, d_drho_e_dT)
qed_eos_pdep_corrections(tplmev, phie=0.0, ..., **pdep_rules)
    -> (d_rho_e_pdep, d_P_e_pdep, d_drho_e_dT_pdep)
qed_plasma_deltas(tplmev, phie=0.0, ..., pdep=True, m_e=M_ELECTRON, nlo=False)
    -> (d_rho_plasma, d_P_plasma, d_drho_plasma_dT)        # photon + e+- summed
rho_plasma_qed_delta(..., nlo=False), pressure_plasma_qed_delta(...), drho_plasma_dT_qed_delta(...)
    -> its three components, one each
\end{lstlisting}
\code{qed_plasma_deltas} is the one-call form used in the coupled-solve RHS (it
evaluates the renorm integral once).  \code{pdep=True} (default, BURST-faithful)
includes the momentum-dependent second-loop terms; the solvers pass it as
\code{renorm_pdep}, and \code{renorm_pdep=False} (with \code{qed_nlo=False}) reproduces the
momentum-independent pin $\Neff=3.044313$.  \code{nlo=True} adds the NLO term of
Section~\ref{sec:nloqed}; it defaults to \code{False} at this level, while every solver
passes \code{qed_nlo=True}.  Add the deltas to the corresponding \code{thermo} quantities.

\begin{fmap}
The \code{glag1}--\code{glag7} correction terms in \code{bbn_therm.f90}
(momentum-independent block plus the second-loop pdep block), verified.
\end{fmap}

\begin{diffbox}
\textbf{Not ported:} the \code{glag6} ($\dd\rho_e/\dd\phi_e$) EoS-derivative term.
The degeneracy $\phi_e$ itself is carried (settable, or auto-solved from
charge neutrality; Section~\ref{sec:knownissues}); only this QED-EoS feedback on
$\phi_e$ is omitted, and the physical $\phi_e\sim10^{-10}$ makes it negligible.
\end{diffbox}

\paragraph{Demo: the QED shift in the plasma energy}
\begin{lstlisting}
import neutt
from neutt import thermo
T = 1.0
rho0 = float(thermo.rho_plasma(T))
drho = float(thermo.rho_plasma_qed_delta(T, pdep=True))
print(rho0, drho, drho/rho0)     # per-cent-level QED correction to rho_plasma
\end{lstlisting}

\subsubsection{NLO QED (the \texttt{qed\_nlo} default)}
\label{sec:nloqed}

Everything above is the leading-order $\mathcal O(\alpha)$ finite-$T$ QED
correction (the thermal self-energies and the \code{glag} terms), which is all
Fortran BURST has.  The port adds the \textbf{next-to-leading-order
$\mathcal O(e^3)$} correction on top, and it is on by default.

\paragraph{The flag}  \code{qed_nlo} (boolean) toggles the NLO term, and it
defaults \textbf{on} at every entry point and builder: \code{make_decoupling_solver},
the RHS/kernel builders (\code{make_rhs}, \code{make_split_rates}),
\code{make_hybrid_solver}, the full density-matrix QKE solver
(\code{make_full_coupled_solver}, Section~\ref{sec:qke}), and the coupled-BBN driver (next
paragraph) all run NLO by default.  Passing \code{qed_nlo=False} anywhere recovers the
bit-faithful LO BURST replica (the Fortran-parity oracle tests set it explicitly).

\paragraph{The physics}  The NLO term is the $\mathcal O(e^3)$ finite-$T$ QED
interaction pressure that the LO renormalization omits, a port of LINX's
\code{explicit_P3} (Bennett et al., \href{https://arxiv.org/abs/1911.04504}{arXiv:1911.04504},
Eq.~4.24),
\[
  \delta P_{e^3}=\frac{e^3}{12\pi^4}\,T^4\,J^{3/2},
  \qquad
  J=\sum_i w_i\,\frac{2x_i^2+(m_e/T)^2}{E_i}\,\frac{2}{e^{E_i}+1},
\]
with $e^3=(4\pi\alpha)^{3/2}$, $T=\Tpl$, $E_i=\sqrt{x_i^2+(\me/T)^2}$, and $x_i$, $w_i$ the
nodes and weights of the 25-node plain Gauss--Laguerre rule ($x=p/T$).  The energy
density and heat-capacity corrections follow from the thermodynamic relations
$\delta\rho=T\,\dd P/\dd T-P$ and $\delta(\dd\rho/\dd T)=T\,\dd^2 P/\dd T^2$, with
the $T$-derivatives taken by forward-mode autodiff (so the term stays
differentiable and varies coherently with a traced $m_e$).  The occupation carries BURST's
$E/T>300$ cutoff, as the $\mathcal O(\alpha)$ terms do, so the term is exactly zero, with zero
derivatives, once every node is cut ($T<\me/300$).  It shifts $\Neff$ by
$-9.5\times10^{-4}$ ($3.0442761\to3.0433222$; $3.0442252\to3.0432716$ at
\code{SIN2THW_BURST}), reconciling with Cielo et al.\
(\href{https://arxiv.org/abs/2306.05460}{arXiv:2306.05460}).  The implementation
is \code{thermo.qed_eos_e3_pressure} / \code{qed_eos_e3_deltas}, added by
\code{qed_plasma_deltas(..., nlo=True)}.

\paragraph{The double-count guard (coupled BBN)}  In \code{make_coupled_bbn} with \code{transport="full"} (Section~\ref{sec:linx}; under
\code{"qke"} set \code{qke_kwargs["qed_nlo"]} to match) the
transport \code{qed_nlo} is tied to LINX's own \code{BackgroundModel}
\code{NLO} flag (default \code{True}).  This matters: the transport NLO acts on
the neutrino-heating / $\Neff$ sector, while LINX's native NLO acts on the
(disjoint) charged-plasma EoS, and the coupling discards LINX's own
neutrino sector and swaps in the transport one, so the QED$\to\nu$ effect must
be supplied exactly once.  Tying the two flags guarantees that; an explicit
\code{solver_kwargs["qed_nlo"]} (or \code{qke_kwargs["qed_nlo"]} in qke mode) that disagrees
with LINX's \code{NLO} still works but emits a mismatch warning.  (The coupled $\dNeff$ between NLO and LO is
$-9.5\times10^{-4}$, the same as standalone, confirming no double-count, which
would have given $\approx-1.9\times10^{-3}$.)

\begin{diffbox}
NLO QED is a documented improvement over Fortran BURST (which carries only
the LO correction), ported from LINX.  \code{qed_nlo=False} restores exact LO
BURST parity.
\end{diffbox}

\subsubsection{The decoupling ``clock''}
\label{sec:background}

\paragraph{Purpose}
The standalone plasma-temperature and expansion evolution, with neutrinos as an
instantaneously-decoupled relativistic bath.  Photons reheat through
$e^\pm$ annihilation, recovering $\Tg/\Tnu\to(11/4)^{1/3}$ and $\Neff\to3$.  This
is the reference background without neutrino heating; the coupled solver (Section~\ref{sec:solver}) adds the
collision energy-exchange that makes $\Neff>3$.

\paragraph{API}
\begin{lstlisting}
rho_nu(T_nu, n_nu=3.0)              # n_nu (7/8)(pi^2/15) T_nu^4
hubble(rho_tot)                    # sqrt(8 pi rho /(3 MPL^2))/HBAR   [1/s]
neff(T_nu, T_gamma, n_nu=3.0)      # (8/7)(11/4)^{4/3} rho_nu/rho_gamma
evolve_background(T_init=100.0, lna_max=10.0, n_save=256,
                  rtol=1e-10, atol=1e-12) -> dict
comoving_plasma_entropy(T_pl, lna) # (rho_pl+P_pl)/T_pl * a^3   (conserved)
\end{lstlisting}
\code{evolve_background} integrates $\dd\Tpl/\dd\ln a=-3(\rho_{\rm pl}+P_{\rm pl})/
(\dd\rho_{\rm pl}/\dd T)$ (covariant energy conservation of the decoupled,
$\mu=0$ plasma; $H$ drops out of the $T$-equation) with diffrax Tsit5.  The
returned dict has \code{lna}, \code{T_pl}, \code{T_nu} ($=$\code{T_init}$\,e^{-\ln a}$, the
decoupled-neutrino temperature), \code{ratio} ($=\Tpl/\Tnu$), \code{neff_final}, \code{ratio_final}.

\begin{fmap}
\code{hubble} matches the snapshot's \code{bbn_yes_evolve_v2.f90:156}
(\code{hubcst = sqrt(8/3 pi/mpl^2 thm(10)/mev4tocgs)/hbar}, \code{thm(10)} the total energy
density), verified; the port works in
MeV$^4$ and skips BURST's \code{mev4tocgs} unit conversion.  The
$\dd\Tpl/\dd\ln a$ form replaces BURST's thermo-driven temperature stepping.
\end{fmap}

\paragraph{Demo: the instantaneous-decoupling sanity check}
\begin{lstlisting}
import neutt
from neutt.thermo import evolve_background
bg = evolve_background()
print(float(bg["ratio_final"]), float(bg["neff_final"]))  # ~1.401019, ~3.000007
\end{lstlisting}
The ratio matches $(11/4)^{1/3}=1.401020$; $\Neff\to3$ to $7\times10^{-6}$.  The
gap to $3.044$ is exactly the transport heating the coupled solver adds.

% ===========================================================================
\section{The Standard-Model transport kernels}
\label{sec:kernels}
% ===========================================================================

The five weak collision kernels and their assembly into the RHS live under
\code{neutt/transport/}.  This is the hard core of the port, the dominant
cost in BURST and the part that the kernel-tensor reformulation
(Section~\ref{sec:kerneltensor}) exists to accelerate.

\subsection{File inventory and Fortran mapping}
\label{sec:kernelmap}

\begin{center}\small
\begin{tabular}{@{}R{0.27\textwidth}R{0.42\textwidth}R{0.22\textwidth}@{}}
\toprule
\textbf{\texttt{neutt} module} & \textbf{physical channel} & \textbf{snapshot Fortran (oracle)} \\
\midrule
\code{transport/nunu.py}     & $\nu\nu,\ \bar\nu\bar\nu$ elastic scatter            & \code{trans_nunu.f90} \\
\code{transport/nunubar.py}  & $\nu\bar\nu$ scatter $+$ $\nu_i\bar\nu_i\!\to\!\nu_j\bar\nu_j$ annih & \code{trans_nunubar.f90} \\
\code{transport/nuer1.py}    & $\nu e^\pm\!\to\!e^\pm\nu$ (R1 reduction)            & \code{trans_nuer1.f90} \\
\code{transport/nuer2.py}    & $\nu e^\pm\!\to\!e^\pm\nu$ (R2 reduction)            & \code{trans_nuer2.f90} \\
\code{transport/nuepma.py}   & $\nu\bar\nu\!\to\!e^-e^+$ annihilation               & \code{trans_nuepma.f90} \\
\code{transport/assemble.py} & precision gate $+$ RHS assembly                      & \code{trans_evolve.f90} \\
\bottomrule
\end{tabular}
\end{center}

\subsection{\texttt{assemble}: precision gate and RHS}
\label{sec:assemble}

\paragraph{The main RHS}
\begin{lstlisting}
assemble_rhs(grid, occ, tcmev, tplmev, phie, prec_trans,
             nunu_kernel, nunubar_kernel, nuratiotol=3.0,
             m_e_eff=M_ELECTRON, use_sched=False, gate_width=0.0, range_limit=False,
             lag_tail_fix=True)
    -> (doccfracdt[3,2,N],  drhonudt)        # [s^-1],  [MeV^4/s]
\end{lstlisting}
Here \code{occ} is the occupation $f[\text{flavour}=3,\text{cp}=2,\text{bin}=N]$,
and this index convention is used throughout: \textbf{flavour}
$\in\{0{:}e,\ 1{:}\mu,\ 2{:}\tau\}$, \textbf{cp} $\in\{0{:}\nu,\ 1{:}\bar\nu\}$,
\textbf{bin} the $\eps$-grid index (so \code{f[0,0]} is the $\nu_e$ spectrum,
\code{f[0,1]} the $\bar\nu_e$, etc.).
\code{tcmev}$=\Tcm$, \code{tplmev}$=\Tpl$ in MeV; \code{phie}$=\phi_e$;
\code{prec_trans} the noise floor; the two \code{*_kernel} are the precomputed
grid-only kernels (Section~\ref{sec:gridkernels}); \code{nuratiotol} is the gate threshold
of Eq.~\eqref{eq:gate}, \code{use_sched} the lepton-kernel path (Section~\ref{sec:sched}),
and \code{gate_width} an optional smooth gate (below).  The internals follow
\code{trans_evolve.f90} exactly:
\begin{enumerate}[nosep]
\item form $\mathrm{meps}=\me/\Tcm$ (the \code{m_e_eff} argument is the
  bare mass, default \code{M_ELECTRON} despite its name, distinct from
  \code{thermo.m_e_eff}, the thermal effective mass, which is not fed to
  the kernels) and $\mathrm{tcmpl}=\Tcm/\Tpl$;
\item evaluate all five kernels to get, per (flavour, cp, bin),
  $\mathrm{numer}=\sum\mathrm{net}$ (all channels), $\mathrm{denom}=\sum\mathrm{frs}$,
  and $\mathrm{leg4}=\mathrm{net}_{R1}+\mathrm{net}_{R2}+\mathrm{net}_{\rm pma}$
  (only the three electron-coupled lepton kernels carry plasma exchange);
\item apply the precision gate~\eqref{eq:gate}:
  $\mathrm{doccfracdt}=\Tcm^5\,\mathrm{net}$ where the gate is open, else $0$;
\item form $\mathrm{drhonudt}=\frac{\Tcm^9}{2\pi^2}\sum_i w_i\eps_i^3
  \sum_{\rm flavour,cp}(\text{gated }\mathrm{leg4})$, the energy exchanged with
  the plasma.
\end{enumerate}
\code{doccfracdt} is $\dd f/\dd t$ in s$^{-1}$; the division by $H$ to get
$\dd f/\dd\ln a$ happens in the decoupling driver, not here.

\paragraph{The \texttt{leg4} convention}
The name comes from BURST's rate arrays \code{ratescatt(4,2,3,2,nbins)} and
\code{rateannih(4,3,2,nbins)} (\code{transvar.f90}, allocated in \code{trans_init.f90}), whose
leading index labels the other species in the process (the scattering partner, or the pair an
annihilation produces): slots 1--3 a neutrino flavour, slot~4 the $e^\pm$.  Only \code{nuer1} and \code{nuer2} ($\nu e^\pm$ scattering) and \code{nuepma}
($\nu\bar\nu\leftrightarrow e^-e^+$) fill slot~4; \code{nunu} and \code{nunubar}, whose partners
are all neutrinos, fill slots 1--3.  Slot~4 is therefore the part of the net rate that exchanges
energy with the plasma, and \code{trans_evolve.f90} sums it alone into \code{drhonudt}.  The port
encodes this by summing only the three lepton-kernel nets into \code{leg4}.  It is unrelated to the
fourth momentum leg $p_4$ of a $1+2\to3+4$ process, the off-grid leg read by interpolation
(Section~\ref{sec:kerneltensor}).

\paragraph{The affine split (for the hybrid head)}
\begin{lstlisting}
assemble_split_rhs(grid, occ, tcmev, tplmev, phie, prec_trans,
                   nunu_kernel, nunubar_kernel, nuratiotol=3.0,
                   m_e_eff=M_ELECTRON, use_sched=False, range_limit=False, fd_cap_logC=inf,
                   lag_tail_fix=True,
                   fd_decay=False, monotone_read=False)
    -> (A[3,2,N], B[3,2,N], drhonudt)        # df/dt = A - B*occ, both >= 0
\end{lstlisting}
This recasts the gated RHS as $\dd f/\dd t=A-B f$ (Eq.~\eqref{eq:affine}),
$A=\Tcm^5\,\mathrm{forward}/(1-f)$, $B=A+\Tcm^5\,\mathrm{reverse}/f$, pure algebra
on what \code{assemble_rhs} already computes (reproduces it to
$\le3.5\times10^{-12}$ relative).  The $(1-f)$ division floors at
\code{SPLIT_ONEMF_FLOOR}$=10^{-3}$ (exact for $f\le0.999$); $A,B$ are clamped
$\ge0$ so the exponential step~\eqref{eq:expstep} cannot push $f$ below the
occupation floor.  Used by the exponential-midpoint steps of the hybrid head
(Section~\ref{sec:hybrid}) and of the QKE collision substep (Section~\ref{sec:qkesplit}).

\paragraph{Equilibrium occupation and the noise floor}
\begin{lstlisting}
equilibrium_occ(grid, xi=(0.0,0.0,0.0)) -> f0[3,2,N]    # 1/(exp(eps - parfact*xi)+1)
compute_prec_trans(grid, nunu_kernel, nunubar_kernel,
                   xi=(0,0,0), use_sched=False, lag_tail_fix=True) -> prec_trans[3,2,N]
compute_prec_trans_at(grid, nunu_kernel, nunubar_kernel, occ_eq, meps, tcmpl,
                      phie=0.0, use_sched=False, lag_tail_fix=True) -> prec[3,2,N]
\end{lstlisting}
\code{equilibrium_occ} builds the Fermi--Dirac initial condition, or IC
(\code{parfact}$=+1$, i.e.\ $+\xi$, for cp$=0$; $-\xi$ for cp$=1$; the $\eps=0$ bin is exactly $0.5$ at $\xi=0$, more generally
$1/(e^{\mp\xi}+1)$).  \code{compute_prec_trans} runs all
five kernels at that equilibrium spectrum (with $\mathrm{meps}=0$,
$\mathrm{tcmpl}=1$, $\phi_e=0$) and records $|\sum\mathrm{net}|/\sum\mathrm{frs}$
per bin, the per-bin interpolation-cancellation noise floor.  It is a static
precompute (grid $+\,\xi$ only), wrapped in the \code{optimization_barrier} of
Section~\ref{sec:gatenum}.  It calls \code{compute_prec_trans_at}, which measures the same floor
at any given equilibrium occupation \code{occ_eq} and background (\code{meps}, \code{tcmpl},
\code{phie}); the $\eps=0$ bin is $0$ and bins with $\sum\mathrm{frs}\le0$ get \code{PREC_FLOOR}.

\begin{fmap}
\code{assemble_rhs} $\equiv$ snapshot \code{trans_evolve.f90:195--208} (the
combined gate); \code{compute_prec_trans} $\equiv$ \code{trans_init.f90:268--286};
\code{equilibrium_occ} $\equiv$ \code{trans_init.f90:185--196}.  Module constants:
\code{NURATIOTOL = 3.0} (from \code{trans_params.ini}), \code{OCC_FLOOR = 1e-300},
\code{PREC_FLOOR = 1e-13}.
\end{fmap}

\begin{diffbox}
\textbf{A single combined precision gate.}  The port applies one gate over
the summed \code{numer}/\code{denom} with a single \code{prec_trans} noise floor,
bit-faithful to the Fortran reference and validated to $\Neff=3.044225186$ (LO, at \code{SIN2THW_BURST}$=0.23$).
The alternative, gating each collision channel separately, is not used; the
combined gate is what the reference does and what the port reproduces.
\end{diffbox}

\begin{diffbox}
\textbf{Port-only knobs in \code{assemble}.}  \code{gate_width>0} switches the
hard gate~\eqref{eq:gate} to a smooth sigmoid weight (a continuous-RHS
robustness lever, default $0=$ BURST's hard gate); \code{fd_cap_logC} is the
FD-envelope cap (default $\infty$, a no-op).  \code{OCC_FLOOR} clamps $f$ to
$[10^{-300},1]$ before any $\log$, a positivity guard BURST does not have (it
relies on the interp clamp instead).  All three are no-ops for the SM solve.  The
module also exposes \code{decoherence_rate} (ungated $g+l$, for the QKE
off-diagonal damping) and \code{compute_sumrules} (below); both are
diagnostics/QKE, not part of the core SM RHS.
\end{diffbox}

\paragraph{The QKE damping rates}
\begin{lstlisting}
decoherence_rate(grid, occ, tcmev, tplmev, phie, nunu_kernel, nunubar_kernel,
                 m_e_eff=M_ELECTRON, use_sched=False, fd_cap_logC=inf,
                 fd_decay=False, monotone_read=False, lag_tail_fix=True) -> Gamma[3,2,N]   # [s^-1]
decoherence_rate_split(...same arguments...) -> (gamma_nue, gamma_nunu)
\end{lstlisting}
\code{decoherence_rate} is the ungated $\Gamma_a=g_a+l_a$, the $B$ of
\code{assemble_split_rhs} without the gate, which the QKE solve uses as the scalar coherence
damping (Section~\ref{sec:qkecollision}).  \code{decoherence_rate_split} splits it into the
$\nu$--$e^\pm$ part (\code{nuer1}$+$\code{nuer2}$+$\code{nuepma}) and the $\nu$--$\nu$ part
($\nu\nu$ and $\nu\bar\nu$), whose sum reproduces it on a physical state (each part is clamped at
$0$ separately); \code{collision="matrix_nue"} adds only the
$\nu$--$\nu$ part as a scalar rate.

\subsection{The grid-only kernels: \texttt{nunu} and \texttt{nunubar}}
\label{sec:gridkernels}

These two depend only on the $\eps$-grid (no $\me$, $\Tcm/\Tpl$, or $\phi_e$),
so they are precomputed once into static tensors and never rebuilt.

\begin{lstlisting}
build_nunu_kernel(grid, n_gl=5, ordpolint=5)    -> NuNuKernel       # host precompute
nunu_rates(kernel, occ, n_chunks=None, fd_cap_logC=inf, fd_decay=False, monotone_read=False) -> (net, frs)   # (3,2,N) each

build_nunubar_kernel(grid, n_gl=5, ordpolint=5) -> NuNuBarKernel
nunubar_rates(kernel, occ, n_chunks=None, fd_cap_logC=inf, fd_decay=False, monotone_read=False)
    -> (net_s, frs_s, net_a, frs_a)             # scatter + annihilation channels

nunu_dfdt(kernel, occ, t_cm_mev, n_chunks=None)      # T_cm^5 * net: ungated df/dt [s^-1]
nunubar_dfdt(kernel, occ, t_cm_mev, n_chunks=None)   # T_cm^5 * (net_s + net_a)
\end{lstlisting}
\code{NuNuKernel} carries the static $W$, $\mathrm{idx}_3/\mathrm{lw}_3/
\mathrm{idx}_4/\mathrm{lw}_4$, and the flavour factor (1 on the diagonal,
$\tfrac12$ off).  \code{nunu} has $K=15$ slots (three $p_3$ segments $\times5$
Gauss nodes); \code{nunubar} has $K=10$ (two segments) and yields both a
scattering and a flavour-changing-annihilation channel from one weight tensor,
with the $\bar\nu$ partner read at the opposite cp.  Neither has an $e^\pm$ partner, so neither
contributes to \code{leg4}.  The \code{*_dfdt} helpers return one kernel's ungated
$\dd f/\dd t$, for tests and diagnostics.

\code{n_chunks} bounds the peak memory of the $\mathcal O(N^2)$ partner-bin
gather: \code{1} is the verbatim path (bit-faithful), \code{None} auto-picks
(small $N\to1$), \code{>1} sums the partner-bin axis in blocks via \code{scan}.
Performance only; bit-identical at \code{n_chunks=1}.

\paragraph{Demo: the equilibrium sum rule}
At equilibrium every net rate must vanish to the noise floor.  This is the
cheapest correctness check and needs no Fortran run:
\begin{lstlisting}
import neutt, numpy as np
from neutt.grid import make_eps_grid
from neutt.transport.assemble import equilibrium_occ
from neutt.transport.nunu import build_nunu_kernel, nunu_rates
g  = make_eps_grid()
f0 = equilibrium_occ(g)                 # FD equilibrium
K  = build_nunu_kernel(g)
net, frs = nunu_rates(K, f0)
print(float(np.max(np.abs(net))))       # ~1e-9 or smaller: net -> 0 at equilibrium
\end{lstlisting}

\subsection{The lepton kernels: \texttt{nuer1}, \texttt{nuer2}, \texttt{nuepma}}
\label{sec:lepkernels}

These carry the energy exchange with the plasma, so they depend on the
background through $\mathrm{meps}=\me/\Tcm$, $\mathrm{tcmpl}=\Tcm/\Tpl$, and
$\phi_e$.  They are not static tensors but jittable functions:
\begin{lstlisting}
nuer1_rates(grid, occ, meps, tcmpl, phie, n_gl=5, n_lag=5, fd_cap_logC=inf, fd_decay=False, monotone_read=False, lag_tail_fix=True, lag_tail_gap=6.0) -> (net, frs)
nuer2_rates(...same keywords...) -> (net, frs)
nuepma_rates(...same keywords...) -> (net, frs)
nuer1_dfdt(grid, occ, meps, tcmpl, phie, t_cm_mev, **kw)   # T_cm^5 * net; likewise nuer2_dfdt, nuepma_dfdt
\end{lstlisting}
Here $p_1=\eps_i$ is the output bin and $E_2$, $E_3$ the incoming and outgoing $e^\pm$
energies in units of $\Tcm$ (for \code{nuepma}, the $e^-$ and $e^+$ energies $E_{\rm out}$,
$E_{\rm in}$); \code{n_gl} and \code{n_lag} are the Gauss--Legendre and Gauss--Laguerre node
counts per tile (BURST's $5$).  The double integral over these energies is organized as a
fixed table of tiles ($E_2$/$E_3$ sub-intervals using Legendre on finite stretches and
Gauss--Laguerre on $[a,\infty)$ tails) and a small number of regimes that
become runtime masks: \code{nuer1} has 2 regimes (split at $p_1=\mathrm{meps}/2$),
\code{nuer2} and \code{nuepma} have 4, R1--R4 (split at the \code{pcase} energies, the
$e^\pm$-mass thresholds $\mathrm{meps}(\sqrt5-1)/4$, $\mathrm{meps}/(2\sqrt2)$,
$\mathrm{meps}/2$ for \code{nuer2} and $\mathrm{meps}/2$, $\mathrm{meps}(1+\sqrt5)/4$,
$\mathrm{meps}$ for \code{nuepma}).  The off-grid leg $p_4=p_1+E_2-E_3$
($p_4=E_{\rm out}+E_{\rm in}-p_1$ for \code{nuepma}) is read by the dynamic interpolator
(Section~\ref{sec:interp}), because its node positions move with the background.
The coupling weights carry $\sin^2\theta_W$ and the $\nu_e$-vs-$\nu_{\mu,\tau}$
charged-current factor.

\paragraph{The quadrature fix (\texttt{lag\_tail\_fix}, default on)}
Fortran BURST has a latent quadrature bug wherever an energy cut diverges: it
applies a fixed 5-point Gauss--Legendre rule over a near-infinite interval and
mis-resolves the exponentially decaying integrand.  In \code{nuer1} and
\code{nuepma} the cut diverges as a grid momentum approaches a regime boundary, so
the per-bin rate is wrong by $2$--$4\times$ in a narrow window and the kernel is
discontinuous at the boundary (by $\sim80\%$ of peak).  In \code{nuer2} it
diverges inside regime R1 (the lowest-$p_1$ regime, not the R1 reduction) as $p_1\to0$:
with the kinematic cuts $\mathrm{ecut1}=p_1+\mathrm{meps}^2/(4p_1)$ and
$\mathrm{ecut2}=p_1+\mathrm{meps}(p_1+\mathrm{meps})/(2p_1+\mathrm{meps})$, R1's outer segment
$[\mathrm{ecut2},\mathrm{ecut1}]$ (\code{trans_nuer2.f90}) has width
$\approx\mathrm{meps}^2/(4p_1)$, tens of $\Tcm$ at the lowest bins once
$\mathrm{meps}\gtrsim5$ ($\Tcm\lesssim0.1$~MeV), and there the rate is off by
$23\%$ at $\mathrm{meps}=7$, rising to $\sim100\%$ at $\mathrm{meps}=20$ (against a
converged $40\times40$-node reference).  The port reproduces the Fortran to
$\le2.5\times10^{-12}$ per bin on the \code{nuer2} oracle, including those bins at
$\mathrm{meps}=12.7$ and $33.9$, so the defect is BURST's.

The port fixes all three, and the fix is the default (\code{lag_tail_fix=True} in every
function that takes the argument; the solvers that do not expose it, the QKE solver and the LINX
adapter's, run the kernels at that default): the wide bins switch to a Gauss--Laguerre tail.  In \code{nuer1} and
\code{nuepma} it is the rule BURST already uses on the other side of the
boundary.  In \code{nuer2}, R1 bins with $\mathrm{ecut1}-\mathrm{ecut2}\ge6\,\Tcm$
integrate $E_3\in[\mathrm{ecut2},\infty)$ with one Laguerre rule in place of
BURST's Legendre $[\mathrm{ecut2},\mathrm{ecut1}]$ plus Laguerre
$[\mathrm{ecut1},\infty)$ pair, with the pair's own $E_2$ tiles: below ecut1 the
$E_2$ integral ends at its kinematic upper limit \code{elim1}, above it \code{elim1} has
diverged and the tail
$[E_3,\infty)$ applies.  The \code{nuer1} meps-sweep discontinuity collapses
($\sim0.81\to0.025$), the lowest \code{nuer2} bins come within $0.4$--$6\%$ of the
converged rate (the accuracy of BURST's 5-point rules at those bins), detailed
balance is preserved exactly, and the scheduled kernels reproduce the reference
ones to $10^{-16}$.  The end-to-end $\Neff$ impact is $\le8.4\times10^{-7}$ for
\code{nuer1}/\code{nuepma} and below $10^{-9}$ for \code{nuer2} (the
\code{rtol=1e-8} pins agree to nine digits with and without it): local-correctness fixes
that are $\Neff$-neutral.  \code{lag_tail_fix=False} recovers the bit-faithful BURST
form, and is used only by the Fortran-parity oracle tests.

\paragraph{Quadrature accuracy}  Even with the fix, BURST's 5-point rules leave a
per-bin error of a few per cent at the lowest bins below $\Tcm\approx0.25$~MeV,
where the $e^\pm$ are Boltzmann-suppressed, and $\sim10^{-4}$ (energy-weighted)
above.  Raising the order to $(n_{\rm gl},n_{\rm lag})=(16,10)$ moves the default
$\Neff$ by $+1.3\times10^{-8}$ in \code{nuer2} and $+3.9\times10^{-6}$ in
\code{nuepma}, so the lepton-kernel quadrature is a $\sim4\times10^{-6}$ systematic
on $\Neff$, dominated by the annihilation kernel.  The pins
(Section~\ref{sec:canonical}) are at BURST's order.  The same rule makes
\code{nuepma}'s loss operator slightly negative at its lowest bin below
$\Tcm\approx0.1$~MeV (on rates that are numerically zero there).

\begin{diffbox}
At the exact regime boundary the Fortran crashes (a zero-width
\code{cscgqf} rescale); the port instead gives a zero-weight tile and stays
finite (values and gradients), so it is strictly more robust there.  That exact
coincidence is measure-zero (the continuous solve never samples it) and is
the one residual the \code{lag_tail_fix} does not address (it does not need to).
\end{diffbox}

\subsection{The scheduled (\texttt{\_sched}) variants, and what \texttt{use\_sched} is}
\label{sec:sched}

\code{use_sched} is a performance switch, not a physics one: it selects how
the lepton kernels are evaluated, changing speed and bit-reproducibility but never
the physics.  ``Never the physics'' is not the same as ``bit-identical'', and the
two scheduled modes differ on exactly that point; see the bit-reproducibility
note below.  Each of the three lepton kernels
splits a bin into kinematic regimes: depending on where the bin energy
$\eps_i$ sits relative to the $e^\pm$ mass thresholds (the \code{pcase}
energies), the energy integral is segmented differently (\code{nuer2} and
\code{nuepma} have 4 such regimes, \code{nuer1} has 2).  A given bin belongs to
exactly one regime.  But the \textbf{plain} kernels, to stay
control-flow-free under \code{jit}/\code{vmap}, evaluate all regimes for
every bin and then mask away the $\sim3/4$ (or $\sim1/2$) that do not
apply, so each bin redundantly does several times the necessary work, and the
launch is dominated by that waste.

The \textbf{scheduled} kernels (\code{nuer1_rates_sched}, \code{nuer2_rates_sched},
\code{nuepma_rates_sched}, each in the same module as its reference kernel)
remove the redundancy.  An integer schedule table
records, for each bin, which regime's quadrature tiles it needs; from there the
kernel can either gather just those active tiles (\code{select=True}: one
batched contraction over the schedule) or rebuild the same masked union of
all regimes without a gather (\code{select=False}).  ``Scheduled'' refers to that
table.  Either way it is a pure-XLA reformulation that is $\sim$6--17$\times$
cheaper than the plain kernels.

\begin{diffbox}
\textbf{Bit-reproducibility caveat: the two scheduled modes are not
equally faithful.}  This is the one place \code{use_sched} can move the last
digits of $\Neff$, and it is easy to over-state as ``bit-identical''.
\begin{itemize}
\item \code{select=False} (the \code{"union"} mode) evaluates the same
  tiles the plain kernel does and merely skips the redundant launches, so it
  reproduces the reference to $\le10^{-15}$, bit-identical for all practical
  purposes.
\item \code{select=True} (the \code{True} mode) gathers a per-bin subset, which
  re-associates the gate sum by more than a pure round-off reordering.  It
  agrees with the reference only to $\sim$few$\,\times10^{-7}$ in $\Neff$
  (measured $6.7\times10^{-7}$ versus union in the production configuration), not to the
  last bit.  This is below every physics threshold in the port, but it is not
  zero, and it is why the QKE solver defaults to \code{"union"}.
\end{itemize}
\end{diffbox}
A user normally never sets \code{use_sched}; it is exposed so the bit-identical
reference (\code{False}) and the bit-identical union path are available for
cross-checks.

\begin{lstlisting}
nuer1_rates_sched(grid, occ, meps, tcmpl, phie, select=True, ...) -> (net, frs)
# select=True  : per-bin regime GATHER  (the fast path, ~6-17x)
# select=False : regime-masked UNION    (no gather; reproduces the plain kernel exactly)
\end{lstlisting}
The three values of \code{use_sched}:
\begin{center}\small
\begin{tabular}{@{}R{0.13\textwidth}R{0.36\textwidth}R{0.42\textwidth}@{}}
\toprule
\code{use_sched} & kernel path & role \\
\midrule
\code{False}     & plain \code{nuer*_rates}            & bit-faithful reference (slowest) \\
\code{True}      & \code{*_rates_sched(select=True)}         & fast per-bin gather; faithful to $\sim$$6.7\!\times\!10^{-7}$ in $\Neff$, not bit-identical (classical default) \\
\code{"union"}   & \code{*_rates_sched(select=False)}        & regime-masked batch, no gather; bit-identical ($\le10^{-15}$) to the reference (the QKE default) \\
\bottomrule
\end{tabular}
\end{center}
Both scheduled modes perturb the gate sum relative to the plain kernel:
\code{"union"} only at round-off ($\le10^{-15}$), \code{True} more
($\sim$$10^{-7}$, per the caveat above), enough that on the adaptive (Tsit5)
solvers a full solve must cap its maximum step \code{dtmax} (in $\ln a$) at $\le0.10$
(the solver defaults handle this).  The fixed-step QKE \code{lax.scan} has no adaptive controller to
cascade, so it is immune to that hazard and is free to pick whichever mode is more
faithful rather than whichever survives the controller.

\paragraph{Defaults}  The user-facing solvers default to \code{True}
(select): \code{make_decoupling_solver}, \code{make_hybrid_solver}, and the LINX adapter's
\code{make_transport_solver} (Section~\ref{sec:adapter}).  The QKE solver
(Section~\ref{sec:qke}) defaults to \code{"union"}.
The bare low-level builders (\code{assemble_rhs}, \code{make_rhs},
\code{make_split_rates}) default to \code{False} (the reference).  In
\code{make_coupled_bbn}, \code{use_sched=None} (the default) resolves per
transport mode: \code{True} for full / semi, \code{"union"} for
\code{transport="qke"}.  The QKE choice of \code{"union"} over \code{True} is
deliberate, not a fallback: that solve is collision-bound (the lepton kernels are
not its bottleneck), so the $1.18\times$ that the gather buys is negligible, while
\code{"union"} keeps the kernel bit-identical to the reference; it is the
fast and bit-faithful path, and the fixed-step scan removes the dtmax
hazard that would otherwise argue for \code{True}.

\subsection{Conservation sum rules}
\label{sec:sumrules}

BURST monitors number, energy, and $\nu$--$\bar\nu$ (CP) symmetry through five
``sum rules'' (\code{sumrule1}$=$number $\int\eps^2\,\mathrm{net}$ and
\code{sumrule2}$=$energy $\int\eps^3\,\mathrm{net}$ of the $\nu$-only processes,
\code{sumrule3}--\code{5}$=$ the $\nu$--$\bar\nu$ symmetry of the annihilation channels;
each is returned as the ratio of its gated net to its gated frs sum).  The port reproduces
them in \code{assemble.compute_sumrules} (each as net/frs over the gated bins, with the raw \code{sumrule1net}, \code{sumrule1frs}, \ldots), and \code{make_decoupling_solver(sumrules=True)}
recomputes them post-hoc at the saved states (zero impact on the solve).  These
are the free, first-class correctness target of the port: the $\nu$-only processes
conserve neutrino number and energy, so \code{sumrule1} and \code{sumrule2} vanish to
quadrature error at any spectrum, and at equilibrium every net rate vanishes.  (Near
equilibrium the forward and reverse rates cancel; the precision gate of
Section~\ref{sec:gatenum} zeroes the bins where that cancellation leaves only noise.)

\paragraph{Demo: per-mode rates at a non-equilibrium spectrum}
\begin{lstlisting}
import neutt, numpy as np
from neutt.grid import make_eps_grid
from neutt.transport.assemble import (equilibrium_occ, compute_prec_trans,
                                         assemble_rhs, compute_sumrules)
from neutt.transport.nunu import build_nunu_kernel
from neutt.transport.nunubar import build_nunubar_kernel
g  = make_eps_grid()
Kn = build_nunu_kernel(g); Kb = build_nunubar_kernel(g)
prec = compute_prec_trans(g, Kn, Kb)
f = equilibrium_occ(g) * 1.0
f = f.at[0, 0].add(0.01 * np.exp(-(np.asarray(g.abscissas) - 5.0)**2))  # a bump in nu_e
docc, drhonu = assemble_rhs(g, f, 2.0, 2.0, 0.0, prec, Kn, Kb)
print(float(np.max(np.abs(docc))), float(drhonu))   # nonzero: the bump relaxes
sr = compute_sumrules(g, f, 2.0, 2.0, 0.0, prec, Kn, Kb)   # dict of sumrule1 ... sumrule5
print(float(sr["sumrule1"]), float(sr["sumrule2"]))        # number / energy monitors
\end{lstlisting}

% ===========================================================================
\section{The coupled decoupling solver}
\label{sec:solver}
% ===========================================================================

This is the headline of the port: evolve the plasma temperature $\Tpl$ together
with the per-mode neutrino spectrum $f[3,2,N]$ through $e^\pm$ annihilation, so
$\Neff$ emerges self-consistently from the spectral distortions.  The state is
the PyTree $\{\Tpl,\;f\}$; the integration variable is $\ln a$, with
$\Tcm=\Tcm^{\rm init}e^{-\ln a}$.  The module is \code{neutt/decoupling.py}.

\begin{fmap}
The temperature ODE is faithful to \code{bbn_yes_evolve_v2.f90}; the collision
RHS to \code{trans_evolve.f90} (through \code{transport.assemble}).
\end{fmap}

\subsection{The system of equations}
\label{sec:thesystem}

It is worth writing the evolved system out in full, since it is rarely seen in
one place.  The state is $\{\Tpl,\ f_i(\eps_a)\}$, one plasma temperature plus
the occupation of every neutrino mode $i\in\{$flavour$\}\times\{$cp$\}$ (6 modes)
on every grid bin $\eps_a$, and the independent variable is $s\equiv\ln a$,
with the comoving clock $\Tcm(s)=\Tcm^{\rm init}e^{-s}$ fixed analytically.

\paragraph{(1) The neutrino transport}
For each mode $i$ and bin $\eps_a$,
\begin{equation}
  \frac{\dd f_i(\eps_a)}{\dd s}=\frac{1}{H}\,C_i(\eps_a),
  \qquad
  C_i(\eps_a)=\Tcm^5\;\mathrm{net}_i(\eps_a)\;\Theta_{\rm gate},
  \label{eq:sysnu}
\end{equation}
where $C_i$ is the collision rate [s$^{-1}$] summed over the five weak processes
(Section~\ref{sec:collint}).  After the analytic angular reduction each process
contributes (illustrated for one process)
\begin{equation}
  \mathrm{net}_i(\eps_a)=\mathrm{prefac}\sum_b\mathrm{kinem}(a,b)
  \sum_{p_3}w(a,b,p_3)\,F\!\big(f_1,f_2,f_3,f_4\big),
  \label{eq:sysnet}
\end{equation}
with $\mathrm{prefac}$ and $\mathrm{kinem}$ as in Eq.~\eqref{eq:net},
$w(a,b,p_3)=\mathrm{glw}\cdot\mathrm{w_3}$ the combined Gauss and leg-3 quadrature weight,
$f_1=f_i(\eps_a)$, $f_2$ the partner occupation at $\eps_b$, $f_3=f(p_3)$,
$f_4=f(\eps_a{+}\eps_b{-}p_3)$ (the off-grid leg), and the Pauli-blocked
gain$-$loss factor
\begin{equation}
  F=\mathrm{flafact}\Big[(1{-}f_1)(1{-}f_2)f_3 f_4-f_1 f_2(1{-}f_3)(1{-}f_4)\Big]
  \equiv\mathrm{forward}-\mathrm{reverse}.
  \label{eq:sysF}
\end{equation}
Here $\mathrm{flafact}$ is the per-process flavour/spin statistical weight (the
combinatorial factor counting the contributing $\nu$ and $e^\pm$ states).
$\Theta_{\rm gate}$ is the precision gate (Section~\ref{sec:gatenum}); it exists
because near equilibrium forward and reverse nearly cancel, so a bin's net rate
can be dominated by round-off.  It is $1$ where
$|\sum\mathrm{net}|/(\sum\mathrm{frs}\cdot\mathrm{prec})\ge\mathrm{nuratiotol}=3$
and $0$ otherwise, with $\mathrm{frs}=\mathrm{forward}+\mathrm{reverse}$ the
cancellation-free throughput and $\mathrm{prec}$ the per-bin equilibrium-noise
floor.  Equation~\eqref{eq:sysnu} is exactly the gated collision RHS, and it can
be written in the affine form $\dd f/\dd s=A-Bf$ used by the hybrid head
(Section~\ref{sec:hybridnum}).

\paragraph{(2) The plasma temperature}
The decoupled $\mu=0$ $e^\pm/\gamma$ plasma conserves comoving entropy except for
the energy it leaks to the neutrinos, which gives
\begin{equation}
  \frac{\dd\Tpl}{\dd s}=\frac{-3\,(\rho_{\rm pl}+P_{\rm pl})
       -\dfrac{1}{H}\dfrac{\dd\rho_\nu}{\dd t}\Big|_{\rm exch}}
       {\dd\rho_{\rm pl}/\dd\Tpl},
  \qquad
  \frac{\dd\rho_\nu}{\dd t}\Big|_{\rm exch}
   =\frac{\Tcm^9}{2\pi^2}\sum_{i,a}w_a\,\eps_a^3\,
     \big(\text{gated }\mathrm{leg4}\big)_i,
  \label{eq:sysT}
\end{equation}
where the energy-exchange rate $\dd\rho_\nu/\dd t|_{\rm exch}$ collects only the
three $\nu$--$e^\pm$ channels (\code{leg4}, BURST's $e^\pm$ partner slot), the ones that move energy across the
$\nu$/plasma boundary.  The $-3(\rho_{\rm pl}+P_{\rm pl})$ term is adiabatic
dilution; the second term is the extra cooling because the plasma is feeding the
distorting neutrinos.  $\dd\rho_{\rm pl}/\dd\Tpl\equiv C_{\rm pl}$ is the plasma
heat capacity.

\paragraph{(3) Closure: the Friedmann equation and the constitutive relations}
\begin{align}
  H&=\frac{1}{\hbar}\sqrt{\frac{8\pi}{3M_{\rm Pl}^2}\,
        \big(\rho_{\rm pl}+\rho_\nu+\rho_{\rm dr}\big)}, &
  \rho_\nu&=\frac{\Tcm^4}{2\pi^2}\sum_{i}\int_0^{\eps_{\max}}\!\eps^3 f_i(\eps)\,\dd\eps,
  \notag\\
  \rho_{\rm pl}&=\rho_\gamma+\rho_{e^\pm}+\delta\rho_{\rm QED}, &
  \rho_\gamma&=\frac{\pi^2}{15}\Tpl^4,
  \qquad
  \rho_{\rm dr}=\dNeff\,\big(\tfrac{4}{11}\big)^{4/3}\frac{7}{8}\frac{\pi^2}{15}\Tcm^4 .
  \label{eq:sysclose}
\end{align}
$M_{\rm Pl}$ is the Planck mass (\code{MPL}), $\rho_{e^\pm}$ the Gauss--Laguerre free-gas
integral (Section~\ref{sec:thermo}), $\delta\rho_{\rm QED}$ the finite-$T$ QED correction
(Section~\ref{sec:qedeos}), and $\rho_{\rm dr}$ the optional inert dark radiation, with
$\dNeff$ here the input \code{delta_neff_init} (Section~\ref{sec:standalonedemos}).  The headline observable is
the late-time $\Neff=\frac{8}{7}\big(\frac{11}{4}\big)^{4/3}\rho_\nu/\rho_\gamma$.

\paragraph{The QKE generalization}
The flavour-mixing solver (Section~\ref{sec:qke}) replaces the diagonal $f_i(\eps)$
with the $3\times3$ density matrix $\rho(\eps)$ per cp and evolves
\begin{equation}
  \frac{\dd\rho(\eps)}{\dd s}=\frac{1}{H}\Big(-i\,[\,H_{\rm flav}(\eps),\rho(\eps)\,]
        +C[\rho]\Big),
  \label{eq:sysqke}
\end{equation}
with the flavour Hamiltonian $H_{\rm flav}$ of Section~\ref{sec:qke} and the
collision $C[\rho]$ the matrix generalization of the gated $C_i$ of~\eqref{eq:sysnu} (its
kernels the matrix form of~\eqref{eq:sysnet}), to which it reduces on a diagonal $\rho$.

\subsection{\texttt{make\_decoupling\_solver}}
\label{sec:evolve}

The reference transport solve: one adaptive \code{diffrax} solve of the whole
window.  Like every solver it is a factory that builds the kernels, the noise
floor, and the equilibrium initial condition once and returns a jitted \code{solve}
that closes over them (Section~\ref{sec:knobs}).

\paragraph{Signature}
\begin{lstlisting}
make_decoupling_solver(grid=None, *, n_save=64, max_steps=50_000,
                       pcoeff=0.4, icoeff=0.3, dcoeff=0.0,
                       solver=None, implicit=False,
                       nuratiotol=3.0, use_renorm=True, use_sched=True,
                       renorm_pdep=True, qed_nlo=True, tcm_cutoff=0.015, phie=0.0,
                       eta0=None, q_baryon=0.5, xi=(0.0,0.0,0.0), sumrules=False,
                       throw=True, lag_tail_fix=True) -> (solve, grid)
# solve(T_init=10.0, lna_max=7.0, rtol=1e-6, atol=1e-30, dtmax=None,
#       delta_neff_init=0.0, me=None, xi=None) -> dict
\end{lstlisting}
The builder's \code{xi} is the build-time degeneracy (it sets the baked noise
floor and initial spectrum), so a scan over $\xi$ uses \code{solve(xi=...)}
instead.  The omitted build-time arguments (\code{dark}, \code{chi_sector},
\code{decay_tail}, \code{decay_to_sm} and their satellites) and the per-call
\code{f_x0}, \code{f_chi0} drive the dark sectors of Part~II
(Sections~\ref{sec:darktransport}--\ref{sec:chinu}); with them off the solve is
the one described here.

\paragraph{Build-time arguments}
\begin{itemize}[nosep]
\item \code{n_save}: number of saved trajectory states.
\item \code{implicit}: \code{True} uses L-stable Kvaerno5 (an alternative for
  a stiff start); \code{False} (default) uses Tsit5, which with the
  $\code{atol}=10^{-30}$ default already completes the $10$~MeV start (it is
  stability-limited but does not blow up; the hybrid is the faster cure).
\item \code{use_renorm}/\code{renorm_pdep}: include the finite-$T$ QED EoS
  correction (default \code{True}$=$BURST-faithful; \code{False} reproduces the
  decoupling-only $\Neff=3.034029$).
\item \code{use_sched}: lepton-kernel path (Section~\ref{sec:sched}; default
  \code{True}).
\item \code{tcm_cutoff}: transport-window low edge ($0.015$ MeV; below it the
  $\propto\Tcm^5$ collision rates are negligible, so they are masked to zero and
  the spectrum is frozen).
\item \code{sumrules}: also return the conservation monitors.
\item \code{grid}: the $\eps$-grid (default \code{make_eps_grid()}); \code{max_steps},
  \code{pcoeff}/\code{icoeff}/\code{dcoeff}: the diffrax step budget and PID gains
  (Section~\ref{sec:odenum}); \code{solver}: an explicit diffrax solver, overriding
  \code{implicit}; \code{throw}: raise on a failed solve.
\item \code{nuratiotol}, \code{qed_nlo}, \code{lag_tail_fix}: the gate threshold
  (Eq.~\eqref{eq:gate}), the NLO QED term (Section~\ref{sec:nloqed}), and the lepton-kernel
  tail quadrature (Section~\ref{sec:lepkernels}).
\item \code{phie}, \code{eta0}, \code{q_baryon}: the electron degeneracy, a fixed float or
  \code{"auto"} for charge neutrality at the present baryon-to-photon ratio \code{eta0}
  (\code{None} is \code{thermo.ETA0_DEFAULT}) and net charge per baryon \code{q_baryon}
  (Section~\ref{sec:thermo}).
\end{itemize}

\paragraph{Per-call arguments (traced)}
\begin{itemize}[nosep]
\item \code{T_init}: the start temperature $\Tpl=\Tcm$ [MeV].  The default
  $10$ is the top of BURST's transport window; neutrinos are tightly coupled
  there ($\Tpl/\Tcm=1$).
\item \code{lna_max}: integrate $\ln a\in[0,\code{lna_max}]$; $\sim7$ lets
  $e^\pm$ annihilation finish so the $(11/4)^{4/3}$ formula is valid.
\item \code{rtol}, \code{atol}, \code{dtmax}: PID tolerances and step cap
  ($10^{-6}$, $10^{-30}$; \code{dtmax=None} is $0.08$ with \code{use_sched},
  $0.15$ on the reference path; Section~\ref{sec:knobs}).
\item \code{delta_neff_init}, \code{me}, \code{xi}: the inert dark radiation, the electron
  mass [MeV] (\code{None} is \code{M_ELECTRON} as a compile-time constant), and a per-call
  degeneracy (\code{None} keeps the build-time floor and initial spectrum); see
  Sections~\ref{sec:standalonedemos} and~\ref{sec:knobs}.
\end{itemize}

\paragraph{Returned dict}
\begin{lstlisting}
out["lna"], out["T_pl"], out["T_cm"], out["ratio"]    # background trajectory
out["f_final"]      # the distorted final spectrum  (3,2,N)
out["f_traj"]       # full trajectory               (n_save,3,2,N)
out["neff_final"]   # N_eff from the distorted final spectrum
out["ratio_final"]  # T_pl/T_cm at the end
out["sol"]          # the raw diffrax solution
out["sumrules"]     # only if sumrules=True
\end{lstlisting}
\code{lna} holds the \code{n_save} saved $\ln a$ nodes, uniform on $[0,\code{lna_max}]$, and
\code{ratio}$=\Tpl/\Tcm$ at each.

\subsection{The RHS builders: \texttt{make\_rhs} and \texttt{make\_split\_rates}}
\label{sec:makerhs}

\code{make_decoupling_solver} assembles its vector field from \code{make_rhs}; the
hybrid head uses the affine \code{make_split_rates}.  Both close over the grid,
the kernels, and the noise floor, and take the same physics knobs.

\begin{lstlisting}
make_rhs(grid, prec_trans, nunu_kernel, nunubar_kernel, T_cm_init,
         nuratiotol=3.0, use_renorm=True, use_sched=False, renorm_pdep=True,
         qed_nlo=True, tcm_cutoff=0.015, phie=0.0, ..., delta_neff_init=0.0,
         me=M_ELECTRON, lag_tail_fix=True, dark=None, chi_sector=None, ...)
    -> rhs(lna, state, args) -> {"T_pl": dT_pl, "f": df}     # per d(ln a)

make_split_rates(grid, prec_trans, nunu_kernel, nunubar_kernel, T_cm_init, ...)
    -> split(lna, T_pl, occ) -> (A, B, dT_pl, H)             # A - B*occ == make_rhs's df
\end{lstlisting}
Here \code{T_cm_init} is $\Tcm$ at $\ln a=0$ (the solvers pass \code{T_init}).  The
temperature derivative is
$\dd\Tpl/\dd\ln a=\big[-3(\rho_{\rm pl}+P_{\rm pl})-\dot\rho_\nu/H\big]/C_{\rm pl}$, with
$\dot\rho_\nu\equiv\dd\rho_\nu/\dd t|_{\rm exch}$ (\code{drhonudt}, Eq.~\eqref{eq:sysT}) and the
plasma heat capacity $C_{\rm pl}=\dd\rho_{\rm pl}/\dd T$ (distinct from
the collision integral $C_{\nu_i}$ of Section~\ref{sec:collint}), and the Hubble rate
$H=\mathrm{hubble}(\rho_{\rm pl}+\rhonu+\rho_{\rm dr})$ includes the inert
dark-radiation term $\rho_{\rm dr}=\code{delta_neff_init}\cdot(4/11)^{4/3}\cdot\frac78\frac{\pi^2}{15}\Tcm^4$
(LINX's \code{Delt_Neff_init} normalization; the $(4/11)^{4/3}$ factor lives in
\code{decoupling._DR_DNEFF_FACTOR}).  The plasma EoS triple
$(\rho_{\rm pl},P_{\rm pl},C_{\rm pl})$ carries the QED $\delta$-corrections when
\code{use_renorm} is on.

\begin{fmap}
The QED correction is BURST's additive $\delta\rho/\delta P$, not a
free-gas mass shift: the lepton-kernel \code{meps} stays the bare
$\me/\Tcm$ (\code{meps = xmelec/tcmev}, snapshot \code{trans_evolve.f90:97}; no thermal
$m_{\rm eff}$ correction is applied to the kernel mass).  A varied \code{me}
varies that bare mass coherently throughout the RHS
(kernels, free-gas EoS, QED $\delta$-terms, and the \code{phie="auto"} charge-neutrality solve).
\end{fmap}

\subsection{The solver menu}
\label{sec:menu}

\begin{center}\small
\begin{tabular}{@{}R{0.27\textwidth}l R{0.50\textwidth}@{}}
\toprule
\textbf{path} & \textbf{warm/call} & \textbf{when} \\
\midrule
hybrid (default) & $\sim$2.8 s & production standalone transport \\
plain Tsit5 (\code{hybrid=False}) & $\sim$6.5 s & reference / debugging; transport starts below $\sim$6 MeV (where the stiff head the hybrid feeds on is mostly gone) \\
implicit Kvaerno5 (\code{implicit=True}) & slower & an L-stable alternative for a stiff start (the explicit default already completes the 10 MeV start, just slowly) \\
\bottomrule
\end{tabular}
\end{center}
The \code{hybrid}/\code{implicit} selectors are the \code{solver_kwargs} keys on
the coupled-BBN path (Section~\ref{sec:linx}); standalone, the hybrid is built with
\code{make_hybrid_solver} (Section~\ref{sec:hybrid}) while \code{make_decoupling_solver}
is the plain Tsit5 path (Kvaerno5 with \code{implicit=True}).  For why the hybrid is the production default and
what its one blind spot is, see Section~\ref{sec:hybrid}.

\subsection{Standalone demos}
\label{sec:standalonedemos}

\paragraph{The pin values (NLO / LO / decoupling-only)}
\begin{lstlisting}
from neutt.decoupling import make_decoupling_solver
print(float(make_decoupling_solver()[0]()["neff_final"]))                  # 3.043322  (NLO QED, default)
print(float(make_decoupling_solver(qed_nlo=False)[0]()["neff_final"]))     # 3.044276  (LO; 3.044225 at SIN2THW_BURST = BURST-faithful)
print(float(make_decoupling_solver(use_renorm=False)[0]()["neff_final"]))  # 3.034029  (no finite-T QED)
\end{lstlisting}
Each line builds and compiles its own solver; in a scan, build once and reuse
the \code{solve}.

\paragraph{The spectral distortion $\delta f(\eps)$}
Distortions are $f/f^{\rm eq}-1$ on the frozen comoving grid (the equilibrium
reference is the same FD array at every epoch):
\begin{lstlisting}
import numpy as np
from neutt.grid import make_eps_grid
from neutt.transport.assemble import equilibrium_occ
from neutt.hybrid import make_hybrid_solver
solve, grid = make_hybrid_solver()
out  = solve()
eps  = np.asarray(grid.abscissas)
dist = np.asarray(out["f_final"]) / np.asarray(equilibrium_occ(grid)) - 1.0
# dist[0,0] = nu_e fractional heating vs eps (positive, grows with eps,
#             percent-level by eps ~ 10); dist[1],dist[2] = the smaller mu/tau distortions
np.savez("spectra.npz", eps=eps, w=np.asarray(grid.weights),
         lna=out["lna"], T_cm=out["T_cm"], T_pl=out["T_pl"], f_traj=out["f_traj"])
\end{lstlisting}

\begin{fmap}
\code{dist[0,0]} is the comoving-grid version of BURST's \code{nu_72.dat}
$\nu_e$ distortion.  The trajectory \code{out["f_traj"]} (shape
$(\code{n_save},3,2,N)$) at the nodes \code{out["lna"]} reproduces, in memory,
what BURST streams to its \code{nu_7x} units over the run.
\end{fmap}

\paragraph{Inert dark radiation ($\dNeff$)}
In the standalone solver the knob is \code{delta_neff_init}; it is traced, so no recompile:
\begin{lstlisting}
solve, grid = make_hybrid_solver()
for dn in (0.0, 0.1, 0.2):
    print(dn, float(solve(delta_neff_init=dn)["neff_final"]))
\end{lstlisting}
\begin{diffbox}
\textbf{One convention for extra radiation.}  \code{delta_neff_init} carries
LINX's \code{Delt_Neff_init} normalization both standalone and coupled:
$\rho_{\rm dr}=\code{delta_neff_init}\cdot(4/11)^{4/3}\cdot\frac78\frac{\pi^2}{15}\Tcm^4$,
so one unit is one fermionic degree of freedom (dof) at the transport start
($\dNeff^{\rm init}=1$).
BURST's bare \code{rhors} knob \emph{omits} the $(4/11)^{4/3}$ factor, so the same number
there is $(11/4)^{4/3}\approx3.85\times$ more radiation; neutt folds that factor into
\code{delta_neff_init} (\code{decoupling._DR_DNEFF_FACTOR}), so standalone and coupled use
the identical convention and the LINX coupling passes \code{delta_neff_init} straight
through with no conversion (Section~\ref{sec:linx}).
\end{diffbox}

\paragraph{A varied electron mass}
\code{me} is traced and reaches the kernels' \code{meps}, the free-gas EoS, and
the QED $\delta$-terms coherently:
\begin{lstlisting}
solve, grid = make_hybrid_solver()
print(float(solve()["neff_final"]),
      float(solve(me=0.6)["neff_final"]))     # dN_eff/dme ~ +0.124 / MeV
\end{lstlisting}

\paragraph{A lepton asymmetry ($\xi$)}
$\xi=(\xi_e,\xi_\mu,\xi_\tau)$ sets both the degenerate FD initial condition and
the matching noise floor.  The builder's \code{xi} bakes them in; \code{solve(xi=...)}
recomputes both inside the compiled solve, so a scan over $\xi$ does not recompile:
\begin{lstlisting}
solve, grid = make_hybrid_solver()
print(float(solve(xi=(0.05, 0.0, 0.0))["neff_final"]))   # degenerate-FD shift; watch out["qualified"]
\end{lstlisting}
\begin{diffbox}
$\xi\ne0$ is validated end to end at one point, $\xi=(0.05,0,0)$ at LO QED: the hybrid
matches an \code{rtol=1e-8} arbiter to $7\times10^{-7}$ in $\Neff$, and $\dNeff$ matches the
degenerate-FD prediction to $0.8\%$.  The test suite pins no $\xi\ne0$ solve (its only
$\nu$/$\bar\nu$-asymmetric case is one spectrum in the \code{nuepma} kernel oracle), so run an
\code{rtol=1e-8} arbiter comparison at any other $\xi$ before production use (and in coupled
runs use full mode; the semi path's plasma side cannot see the
asymmetry).
\end{diffbox}

\paragraph{Conservation telemetry}
\begin{lstlisting}
solve, grid = make_decoupling_solver(sumrules=True)
out = solve()
# out["sumrules"] is a dict of sumrule1 ... sumrule5, each an array along the trajectory:
print(np.asarray(out["sumrules"]["sumrule1"]).shape)   # (n_save,) number monitor
\end{lstlisting}

\subsection{Outputs and I/O: getting at the spectra}
\label{sec:io}

The port writes no files: every data product is returned in memory, unlike
BURST, which streams to numbered Fortran units.  Knowing where each distribution
lives is therefore an I/O question in itself.

\paragraph{What the transport solvers return}
The \code{solve} closures from \code{make_decoupling_solver}
(Section~\ref{sec:evolve}) and \code{make_hybrid_solver} (Section~\ref{sec:hybrid})
return dicts that share
the same user-facing fields (the extras differ: the decoupling \code{solve} adds the
raw \code{sol}; the hybrid \code{solve} adds the qualification diagnostics ---
\code{qualified}, \code{qss_dist}, \code{head_*}, \code{tail_*}; Section~\ref{sec:hybrid}).  The fields
a user reads:
\begin{itemize}[nosep]
\item \code{f_final} $[3,2,N]$: the distorted final occupation
  $f_{\nu_i}(\eps)$ (flavour $\times$ cp $\times$ bin) on the frozen comoving grid;
\item \code{f_traj} $[\code{n_save},3,2,N]$: the trajectory at the saved $\ln a$
  nodes \code{out["lna"]}, i.e.\ the in-memory analogue of BURST's
  \code{nu_72.dat}\,\ldots\,\code{nu_75.dat} streams;
\item \code{T_cm}, \code{T_pl}, \code{ratio}: the background clocks at each
  node; \code{neff_final}, \code{ratio_final}, the end-point scalars.
\end{itemize}
The grid abscissas and Boole weights come from the returned \code{grid}
(\code{grid.abscissas}, \code{grid.weights}); they are what turn an occupation
into a moment, a density, or a distortion.

\paragraph{Recovering the spectral distortion}
The $\eps$-grid is frozen comoving, so the equilibrium reference is the
same Fermi--Dirac array at every epoch and $\delta f=f/f^{\rm eq}-1$ is a
single elementwise division; the demo in Section~\ref{sec:standalonedemos} shows
this and saves the result to an \code{.npz} with \code{numpy.savez}.  Densities
and $\Neff$ from any saved \code{f} go through the \code{grid} moments
(Section~\ref{sec:moments}) and \code{neff_from_spectrum}
(Section~\ref{sec:neff}).

\paragraph{Coupled-run outputs}
In the LINX coupling the products are different: \code{run} returns
\code{(abundance_tuple, N_eff)}, with $\Yp=\code{res[5]}\cdot4$ and
$\DoverH=\code{res[2]}/\code{res[1]}$ (Section~\ref{sec:linx}); the distorted
$\nu_e$ spectrum is reachable through \code{f_e_occupation}
(Section~\ref{sec:adapter}).  As above, persist with \code{numpy.savez}; there
is no built-in file output to configure.

% ===========================================================================
\section{The exponential-midpoint hybrid solver}
\label{sec:hybrid}
% ===========================================================================

This is the production transport solver, in \code{neutt/hybrid.py}.  The
numerical idea is sketched in Section~\ref{sec:hybridnum}; here is the API and
the operating envelope.

\subsection{\texttt{make\_hybrid\_solver}}

\begin{lstlisting}
make_hybrid_solver(grid=None, *, lna_sw=1.5, n_head=1600, head_power=1.0,
                   n_save_head=128, n_save_tail=128,
                   use_sched=True, use_renorm=True, renorm_pdep=True, qed_nlo=True,
                   nuratiotol=None, tcm_cutoff=0.015, phie=0.0,
                   eta0=None, q_baryon=0.5, xi=(0.0,0.0,0.0), lag_tail_fix=True,
                   max_steps=50_000, pcoeff=0.4, icoeff=0.3, dcoeff=0.0,
                   throw=True)  -> (solve, grid)
\end{lstlisting}
\code{throw=False} returns a failed tail solve's partial result with \code{tail_ok} False
instead of raising.  The kernels, the noise floor, and the equilibrium IC are built once at
construction; \code{solve} closes over them.  Arguments:
\begin{itemize}[nosep]
\item \code{lna_sw}: the head$\to$tail handoff in $\ln a$ (static).  Default
  $1.5$ means ``hand off at $\Tcm=0.22\,\code{T_init}$'', tuned for a $10$~MeV
  start (handoff at $2.2$~MeV).  Keep $\lesssim1.5$; a later handoff puts the
  state too far off the QSS attractor and stalls the tail.
\item \code{n_head}: number of fixed exponential-midpoint head steps (static).
  $1600$ uniform is the validated configuration.
\item \code{head_power}: grading exponent for the head step edges; $1=$
  uniform.  Do not use $\ne1$: graded heads have measured failure pockets
  and are parameter-fragile (the qualification catches them, but each failed call
  pays full-solve price).
\item \code{use_sched}, \code{use_renorm}, \code{renorm_pdep}, \code{qed_nlo},
  \code{nuratiotol} (\code{None} is BURST's $3.0$), \code{tcm_cutoff}, \code{phie},
  \code{eta0}, \code{q_baryon}, \code{lag_tail_fix} (Section~\ref{sec:evolve}): forwarded identically to the head
  (\code{make_split_rates}) and the tail (\code{make_rhs}), so both halves are
  the same physics.
\item \code{xi}: the build-time default degeneracy; it sets the
  degenerate-FD initial occupation and the matching \code{prec_trans} noise floor
  once at construction (it is not a \code{make_split_rates}/\code{make_rhs}
  argument, but head and tail both close over \code{prec_trans} and start from the
  same IC, so it reaches them consistently).  \code{solve} also takes a traced
  per-call \code{xi} that overrides this with no recompile (\code{xi=None} reuses
  the baked value bit-for-bit); $\xi\ne0$ shifts the QSS attractor (watch
  \code{qualified}).
\item \code{n_save_head}, \code{n_save_tail} (default 128 each): saved
  trajectory nodes in the head $[0,\code{lna_sw})$ and tail
  $[\code{lna_sw},\code{lna_max}]$; the stitched \code{f_traj} has
  \code{n_save_head}$+$\code{n_save_tail} rows.
\item \code{max_steps}, \code{pcoeff}, \code{icoeff}, \code{dcoeff}: the tail's diffrax
  step budget and PID gains (the production PI($0.4,0.3$)).
\end{itemize}
The head's building blocks are public (the QKE collision substep uses \code{exp_step} too):
\begin{lstlisting}
phi1(z)                                # (e^z - 1)/z, series 1 + z/2 for |z| <= 1e-8
exp_step(f, a, b, h)                   # f + h*phi1(-b*h)*(a - b*f) = Eq. (expstep), exact as b -> 0
head_edges(lna_sw, n_head, power=1.0)  # (n_head+1,) static numpy step edges on [0, lna_sw]
\end{lstlisting}
The midpoint head calls \code{exp_step} twice per step: a half step to the midpoint, then a full
step with $(A,B)$ evaluated there.

\subsection{The returned \texttt{solve}}

\begin{lstlisting}
solve(T_init=10.0, lna_max=7.5, rtol=1e-6, atol=1e-30, dtmax=0.08,
      delta_neff_init=0.0, qual_tol=1e-4, me=None, xi=None) -> dict
\end{lstlisting}
\code{me=None} compiles the bare mass ($\code{me}=\code{M_ELECTRON}$) in as a constant,
the graph the canonical pins use; a passed \code{me} is traced.  \code{delta_neff_init},
\code{qual_tol}, and a per-call \code{xi} are traced (no recompile).  The dict has the base keys \code{lna}, \code{T_pl},
\code{T_cm}, \code{ratio}, \code{f_traj}, \code{f_final}, \code{ratio_final},
\code{neff_final} (it does not carry the decoupling \code{solve}'s raw
\code{sol}) plus the hybrid diagnostics:
\begin{lstlisting}
out["qualified"]   # bool: False means the full-window Tsit5 fallback ran
out["qss_dist"]    # gated max|f - A/B| at the handoff
out["head_T_ok"], out["head_f_finite"], out["head_n_neg"], out["head_n_super"]
out["tail_steps"], out["tail_rejected"], out["tail_ok"]
\end{lstlisting}
The \code{head_*} entries check the handoff state ($\Tpl$ finite and positive, every $f$
finite, and the numbers of bins with $f<0$ and $f>1$); \code{tail_steps},
\code{tail_rejected} and \code{tail_ok} are the tail's accepted and rejected step counts and
its diffrax success flag.

\subsection{Handoff qualification and silent fallback}

At the handoff $\ln a=\code{lna_sw}$ the solve checks that the spectrum is on the
quasi-static-equilibrium attractor and otherwise sane ($\Tpl$ finite and positive, every
$f$ finite and in $[0,1]$):
\[
  \mathrm{qualified}=\mathrm{sane}\ \wedge\ \big(\mathrm{qss\_dist}\le\code{qual_tol}\big),
  \qquad \mathrm{qss\_dist}=\max_{\rm gated\ bins}\big|f_{\rm sw}-A_{\rm sw}/B_{\rm sw}\big|.
\]
where the subscript sw marks values at the handoff.  If the check fails, the tail
\code{diffeqsolve} (the diffrax integration call) is fed a traced start time and state
$t_0,y_0$
that restarts it from $\ln a=0$ at the equilibrium IC, so qualified and
fallback share one compiled executable and the safe path is always available
without a recompile.  The final state and $\Neff$ are exact in both
modes; a mis-qualified configuration costs only the extra head work.  The default
\code{qual_tol}$=10^{-4}$ gives $\sim2\times$ margin at the
$\code{lna_sw}=1.5$ handoff (gated $\max|f-A/B|=5.6\times10^{-5}$ there) while
rejecting every handoff later than $\ln a\approx1.7$.

\subsection{Operating envelope, which solver, when}
\label{sec:whichsolver}

The hybrid's win is the stiff $10\to2.2$~MeV head.  Its one blind spot is
that the head is order-2 and its qualification polices the spectrum, not
accumulated $\Tpl$/energy-budget error, so accuracy-critical physics deposited
in the head can slip through.  The asymmetry that should drive the choice:
picking the hybrid wrongly costs speed (or, for exotic high-$T$ physics, a bias
a one-time pin would catch); but the genuinely dangerous wrong choice is the
semi coupling (Section~\ref{sec:linx}), which silently drops the
$\nu\leftrightarrow$plasma energy exchange.

\begin{center}\small
\begin{tabular}{@{}p{0.50\textwidth}p{0.44\textwidth}@{}}
\toprule
\textbf{scenario} & \textbf{recommendation} \\
\midrule
SM baseline / BBN-only params (baryon-to-photon ratio $\eta$, neutron lifetime $\tau_n$, nuclear rates) & cached transport history (\code{cached_transport_history}, Section~\ref{sec:adapter}); transport is constant, solve once \\
SM $+$ inert dark radiation ($\dNeff$ only) & full default (hybrid) \\
BSM with all dynamics below $\sim2$ MeV & full default; the head is SM-equilibrium there \\
energy injection above $\sim2$ MeV, or a freeze-out/resonance/decay at 3--10 MeV & \code{hybrid=False} (the head's blind spot) \\
transport start below $\sim6$ MeV & \code{hybrid=False}; the stiff head is mostly gone \\
any new model family, first time & run the family pin: solve at default and at \code{qual_tol=-1, rtol=1e-8} (\code{qual_tol=-1} forces the full adaptive fallback), require $|\dNeff|\le2\times10^{-6}$, then scan with whichever passed \\
debugging / bit-level cross-checks & \code{hybrid=False, use_sched=False, rtol=1e-8} (the bit-faithful reference) \\
\bottomrule
\end{tabular}
\end{center}

Watch per-call wall time in scans: a hybrid call that suddenly slows to the
plain-solve time (standalone $\sim$2.8 s $\to$ $\sim$6.5 s; coupled $\sim$4.2 s
$\to$ $\sim$8 s) is the fallback firing (the answer is still right).  If a whole
parameter region does it, switch that scan to \code{hybrid=False}.

\subsection{Validation headline}

Default ($1600$ uniform head steps, \code{lna_sw=1.5}): the hybrid's $\Neff$
sits $-6.9\times10^{-8}$ from the \code{rtol=1e-8} arbiter, comfortably inside
the pre-registered bar ($|\dNeff|\le2\times10^{-6}$ and
$\max|\Delta f_{\rm final}|\le1\times10^{-5}$), at $2.36\times$ the speed of the
plain transport solve (2.78 s vs.\ 6.54 s warm).  The \code{neutt} test suite passes; a
harness that rebuilds the executable for each run finds an executable-to-executable $\Neff$
spread of $7.4\times10^{-7}$; and all $5$ LINX coupling pins are reproduced.

\begin{diffbox}
The hybrid is a port-only construction with no BURST counterpart (BURST uses one
adaptive RK integrator throughout).  It is mathematically exact for the linear
collision diagonal and falls back to the faithful adaptive solve whenever its
premises do not hold, so it adds speed without changing the physics.
\end{diffbox}

% ===========================================================================
\section{Coupling to LINX for abundances}
\label{sec:linx}
% ===========================================================================

The transport solver produces $\Neff$ and the distorted spectra; to get the
light-element abundances it is coupled to \textbf{LINX}, the differentiable BBN
code.  The integration lives in the LINX source tree, in three files:
\code{linx/coupled_bbn.py} (the one-call driver \code{make_coupled_bbn}) and
\code{linx/transport_adapter.py} (the transport$\leftrightarrow$LINX seam,
imported as \code{ta}), plus \code{linx/weak_rates_transport.py} (\code{TransportWeakRates}, the $n\leftrightarrow p$ rates on the distorted $f_e$).

\begin{diffbox}
\textbf{This entire coupling is a port-only addition.}  Fortran BURST has its own
internal nuclear network and computes $\Yp$, $\DoverH$ itself; the JAX port instead
ports only the neutrino transport and hands the distorted spectrum to LINX's
abundance model.  Consequently the JAX$+$LINX path also has no
recombination / sound-horizon (\code{sdratio}) stage; its $\Neff$ is the
late-time frozen value from transport, not a recombination-era observable.  See
Appendix~\ref{sec:fortranmap}.
\end{diffbox}

\subsection{\texttt{make\_coupled\_bbn}}

\begin{lstlisting}
make_coupled_bbn(nuclear_net="key_recommended", use_fe=True, n_save=256,
                 use_sched=None, gate_width=0.0, solver_kwargs=None,
                 device_split=True, cpu=None, gpu=None,
                 transport="full", xi=(0.0,0.0,0.0), grid=None,
                 qke_kwargs=None, solver=None, **bg_kwargs)  -> (run, info)
\end{lstlisting}
Key arguments:
\begin{itemize}[nosep]
\item \code{nuclear_net}: LINX reaction network (default
  \code{"key_recommended"}: the key rates of the 2023 compilation of the PRIMAT
  BBN code, with an updated $n+p\to d+\gamma$ rate).  The choice affects the abundances ($\Yp$, $\DoverH$)
  but not $\Neff$ (which is a transport result, network-independent).
\item \code{use_fe} (default \code{True}): use the distorted-$f_e$ weak rates
  (\code{TransportWeakRates}) rather than equilibrium ones.  This is what threads
  the transport-distorted $\nu_e$ spectrum into the $n\leftrightarrow p$ rates.
\item \code{transport} $\in\{\code{"full"},\code{"semi"},\code{"qke"}\}$: the
  coupling mode (Section~\ref{sec:fullsemi}; \code{"qke"} is the flavour-mixing
  solver of Section~\ref{sec:qke}; its builder options go in \code{qke_kwargs}, whose \code{lna_max}, default $7.5$, goes to the solve).
\item \code{solver}: a convenience front-end that selects the transport solver
  per mode without reaching into \code{solver_kwargs}: \code{"full"} accepts
  \code{hybrid} (default) / \code{tsit5} / \code{implicit}; \code{"semi"} only
  \code{tsit5}; \code{"qke"} \code{hybrid} (default) /
  \code{expm} / \code{strang}.  Passing both \code{solver} and the matching low-level key is an
  error.
\item \code{use_sched}: transport kernel path (Section~\ref{sec:sched}).
  Default \code{None} resolves to each mode's validated default (\code{select}
  for full/semi, \code{"union"} for qke).
\item \code{solver_kwargs}: forwarded to the transport builder in full and semi modes (in full
  mode a key that is neither a build key nor \code{rtol}/\code{atol}/\code{dtmax}/\code{lna_max}
  raises; qke mode takes its options in \code{qke_kwargs} and rejects \code{solver_kwargs}).  Notably
  \code{solver_kwargs={"hybrid": False}} selects plain Tsit5 transport
  (neutt's \code{make_decoupling_solver}, $\sim8$~s), and
  \code{{"lna_max": 6.6}} matches the end temperature for a
  lower transport start.
\item \code{device_split} (default \code{True}): run LINX on CPU and the
  transport solve on GPU, in one process.  (LINX's \code{BackgroundModel}
  compiles prohibitively slowly on GPU.)  \code{False} runs everything on one
  device; with no visible GPU the split switches off by itself.
\item \code{grid}: the transport $\eps$-grid (every mode), fixed at
  build time (never vary per call; it would retrace).
\item \code{**bg_kwargs}: forwarded to LINX \code{BackgroundModel}.
  \textbf{Pass \code{max_steps=8192} here}: the LINX background default (512) is
  marginal for the NLO plasma EoS and can fail the coupled run with a ``maximum
  number of steps'' error.  This is the background ceiling; the
  transport solver's own \code{max_steps} is set separately via
  \code{solver_kwargs}.
\end{itemize}
\code{info} is \code{dict(cpu, gpu, grid, device_split)}.  Everything
parameter-independent (kernels, noise floor, LINX models) is built once;
\code{run} is serialized-MCMC safe (physics scalars are traced, nothing
recompiles per call).

\subsection{The \texttt{run} closure}

\begin{lstlisting}
run(delta_neff_init=0.0, T_anchor=4.0, T_init=10.0, xi=None, return_diag=False,
    **am_kwargs)  -> (abundance_tuple, N_eff)
\end{lstlisting}
\begin{itemize}[nosep]
\item \code{delta_neff_init}: the public inert-dark-radiation input
  (LINX's \code{Delt_Neff_init} convention; see below).
\item \code{T_init}: where the full and qke solves start ($10$ MeV).
\item \code{T_anchor}: where the semi spectrum solve starts ($4$ MeV).
\item \code{xi}: optional per-call neutrino-degeneracy override (Section~\ref{sec:knobs});
  \code{None} reuses the build-time value, a passed triple is traced (no recompile).
\item \code{**am_kwargs}: forwarded to LINX's \code{AbundanceModel}: the
  traced physics scalars \code{eta_fac} (baryon-to-photon rescale),
  \code{tau_n_fac} (neutron-lifetime rescale), \code{nuclear_rates_q} (per-reaction
  log-normal draws), \code{me} (electron mass, threaded coherently to the LINX
  background, the weak rates, the transport and the swapped \code{Neff_vec}; LINX's
  \code{AbundanceModel} builds its own $t(T)$ and $a(T)$ from the EM energy density at
  \code{const.me}, so a varied \code{me} does not reach the network's clock), plus non-traced controls like \code{rtol},
  \code{sampling_nTOp}.
\item \code{return_diag=True} adds a third return value, a CPU dict with the
  transport solve's own \code{neff_final}.
\end{itemize}
The abundance tuple is LINX's \code{AbundanceModel} output; the standard
extractions are $\Yp=\code{res[5]}\cdot4$ and $\DoverH=\code{res[2]}/\code{res[1]}$.
\code{N_eff} is the late-time value at the lowest (frozen) $\Tg$ of the swapped
background, LINX's background with its neutrino sector replaced by the transport
history (\code{ta.swap_neutrino_sector}).

\subsection{Full vs.\ semi coupling}
\label{sec:fullsemi}

\begin{center}\small
\begin{tabular}{@{}R{0.09\textwidth}R{0.49\textwidth}R{0.09\textwidth}R{0.18\textwidth}@{}}
\toprule
\textbf{mode} & \textbf{what it does} & \textbf{warm} & \textbf{$\nu$--plasma feedback?} \\
\midrule
\code{"full"} (default) & per-call full coupled transport: $\Tpl$ co-evolved with the spectrum, finite-$T$ QED (with the momentum-dependent \code{renorm_pdep} term) in the RHS, $10$~MeV start, BURST's own topology.  $\nu\leftrightarrow$plasma energy exchange feeds $\Tpl$/$H$.  The only mode correct for energy-exchanging physics. & $\sim$4.2 s & yes \\
\code{"semi"} & SM-grade fast path: spectrum-only solve driven by LINX's plasma history from a $4$~MeV anchor.  \textbf{Discards} the $\nu\leftrightarrow$plasma exchange ($\dot\rho_\nu$).  (Still produces the SM distortion and accepts inert \code{delta_neff_init} through LINX's background.) & $\sim$2.4 s & no \\
\bottomrule
\end{tabular}
\end{center}

\begin{diffbox}
\textbf{Semi is a silent physics error if misused.}  Full mode reproduces
BURST's coupling topology; semi is a port-only convenience that drops
$\dot\rho_\nu$.  It is admissible only for SM-grade work (or SM $+$ inert dark
radiation) at $\sim10^{-4}$ $\Neff$ tolerance.  Never use it with any sector that
exchanges energy with the neutrinos.  When in doubt, stay on full; nothing
warns you if semi is wrong.
\end{diffbox}

\subsection{The $\dNeff$ convention}
\label{sec:dneffconv}

The single input \code{delta_neff_init} uses LINX's \code{Delt_Neff_init} meaning:
$\rho_{\rm DR}=\code{delta_neff_init}\cdot\frac78(\frac{4}{11})^{4/3}\rho_\gamma$ at the
start, redshifting as $a^{-4}$.  The transport carries the \emph{same}
normalization natively: neutt folds the $(4/11)^{4/3}$ factor into its
dark-radiation term (\code{decoupling._DR_DNEFF_FACTOR}), so
$\rho_{\rm dr}=\code{delta_neff_init}\cdot(4/11)^{4/3}\cdot\frac78\frac{\pi^2}{15}\Tcm^4$
equals LINX's $\rho_{\rm DR}$ and the glue passes \code{delta_neff_init} straight
through with no conversion at the seam.  (BURST's own \code{rhors} omits the
$(4/11)^{4/3}$ factor, so the same number there is $\approx3.85\times$ more
radiation.)  The hybrid stays qualified at least to $\code{delta_neff_init}=0.3$.

\subsection{The transport$\leftrightarrow$LINX seam (\texttt{transport\_adapter})}
\label{sec:adapter}

These are the functions that map the comoving-grid spectrum onto LINX's $\Tg$
clock and feed the distorted $f_e$ into the weak rates.  The most important for a
user:
\begin{lstlisting}
make_transport_solver(grid=None, n_save=256, use_sched=True, use_renorm=True,
                      qed_nlo=True, implicit=False, ..., hybrid=None,   # None: not implicit
                      ...) -> (solve, grid)                 # build-once transport
cached_transport_history(grid=None, T_init=10.0, lna_max=7.5, ...) -> history
make_spectrum_solver(grid=None, n_save=256, use_sched=True, ..., gate_width=0.0,
                     xi=(0.0,0.0,0.0)) -> (solve_spectrum, grid)  # the semi-mode solver
build_history_from_full_coupled(full_out, grid) -> history     # the qke-mode history
build_history_from_transport(transport_out, grid) -> history
f_e_occupation(E_nu_MeV, T_g, history, cp=0)      # distorted nu_e occupation
T_cm_of_Tg(T_g_query, history)                    # comoving temp on the T_g clock
r_nu_of_Tg(T_g_query, history)                    # rho_nu_tot/rho_gamma on the T_g clock
\end{lstlisting}
\begin{itemize}[nosep]
\item \code{make_transport_solver} is the production build-once transport for
  serial MCMC; with \code{hybrid=True} (default) it delegates to
  \code{make_hybrid_solver}, with \code{hybrid=False} to
  \code{make_decoupling_solver}.
\item \code{cached_transport_history} is the SM entry point: transport is
  constant across an SM MCMC, so solve once and reuse the history for every
  abundance call (zero per-call transport cost).
\item \code{f_e_occupation} is the convenience wrapper (slice-once over $\Tg$ then
  evaluate over $\eps$); \code{TransportWeakRates} actually calls the two split
  primitives it composes (\code{f_e_slice_at_Tg} once per $\Tg$ and
  \code{f_e_eval_slice} per momentum point) for compile efficiency.  In the
  undistorted limit it returns exactly $1/(e^{E_\nu/\Tg}+1)$, LINX's thermal
  factor, which is the correctness anchor.  \code{cp}: 0$=\nu_e$, 1$=\bar\nu_e$.
\item \code{r_nu_of_Tg} / \code{T_cm_of_Tg} interpolate clamped into the
  transport $\Tg$ range, so above the transport start ($\nu$ in equilibrium, $r_\nu$
  flat) and below $e^\pm$ annihilation ($r_\nu$ frozen) the fills are physically
  exact, no spurious slope through the BBN epoch.
\end{itemize}
The rest of the public seam, which \code{make_coupled_bbn} composes, and the weak-rate class:
\begin{lstlisting}
make_history(T_g_hist, r_nu_hist, *, ratio_hist=None, f_traj=None, eps=None,
             weights=None) -> history                  # the history dict, sorted ascending in T_g
linx_plasma_nodes(bg_output, T_anchor=4.0, me=None)
    -> (lna_nodes, Tg_nodes, H_nodes, lna0, lna1, Tcm_anchor)   # semi-mode drive, from LINX's background
swap_neutrino_sector(bg_output, history, me=const.me) -> bg_output  # LINX 7-tuple, transport rho_nu, N_eff
rho_nu_vec_on_grid(T_g_vec, rho_g_vec, history)        # per-species rho_nu = r_nu*rho_g/3 [MeV^4]
neff_asymptotic(history)                               # (8/7)(11/4)^{4/3} r_nu at the lowest T_g
neff_vec_on_grid(T_g_vec, history)                     # N_eff(T_g) along LINX's grid
f_e_slice_at_Tg(T_g, history, cp=0) -> (N,)            # f_e(eps) at one T_g, cubic in log T_g
f_e_eval_slice(E_nu_MeV, T_cm, f_slice, eps_grid)      # the slice at E_nu, cubic in eps, 0 past the grid

# linx/weak_rates_transport.py
TransportWeakRates(nu_history, RC_corr=True, FM_corr=True, weak_mag_corr=True, thermal_corr=True)
\end{lstlisting}
\code{make_history} is what the two \code{build_history_*} functions call, and what a stored
trajectory (a saved file, a test fixture) goes through.  \code{bg_output} is the 7-tuple returned
by LINX's \code{BackgroundModel}; \code{swap_neutrino_sector} keeps its electromagnetic sector and
replaces \code{rho_nu_vec} and \code{Neff_vec} from the history.  \code{TransportWeakRates}
subclasses LINX's \code{WeakRates} and replaces only the Born neutrino factor by the transport
$f_{\nu_e}$ or $f_{\bar\nu_e}$ (the channel's neutral lepton is a $\nu$ when
\code{sgnq}$\cdot e_\nu>0$ and a $\bar\nu$ otherwise), which the radiative correction multiplies
as in LINX; the finite-mass correction and LINX's tabulated thermal corrections keep the thermal
neutrino.  \code{make_coupled_bbn(use_fe=True)} builds it from each
history.

\subsection{Coupled-BBN demos}

\paragraph{SM baseline}
\begin{lstlisting}
import neutt
from linx.coupled_bbn import make_coupled_bbn
run, info = make_coupled_bbn(max_steps=8192)
res, neff = run(delta_neff_init=0.0)
print(neff, float(res[5]*4.0), float(res[2]/res[1]))   # N_eff 3.043322; Yp/D-H depend on nuclear_net
\end{lstlisting}

\paragraph{A dark-radiation point}
\begin{lstlisting}
res, neff = run(delta_neff_init=0.2)        # traced -> no recompile
print(neff, float(res[5]*4.0))
\end{lstlisting}

\paragraph{Varying BBN-only parameters}
\code{eta_fac}, \code{tau_n_fac} and \code{nuclear_rates_q} are traced, so a chain over
them never recompiles; for SM-grade work the semi transport is the faster choice:
\begin{lstlisting}
run, info = make_coupled_bbn(transport="semi", max_steps=8192)     # SM-grade, fast
for ef in (0.98, 1.0, 1.02):
    res, neff = run(delta_neff_init=0.0, eta_fac=ef)
    print(ef, float(res[5]*4.0))
\end{lstlisting}

\paragraph{Plain-Tsit5 transport (reference path)}
\begin{lstlisting}
run, info = make_coupled_bbn(max_steps=8192, solver_kwargs={"hybrid": False})   # ~8 s warm
res, neff = run(delta_neff_init=0.0)
\end{lstlisting}

\paragraph{A varied electron mass, end to end}
\begin{lstlisting}
res0, n0 = run(delta_neff_init=0.0)
res1, n1 = run(delta_neff_init=0.0, me=0.6)   # reaches bg + weak rates + transport coherently
print(n0, n1, float(res0[5]*4.0), float(res1[5]*4.0))   # dYp/dme ~ +0.227 / MeV
\end{lstlisting}

% ===========================================================================
\section{Parameters, tolerances, and knobs}
\label{sec:knobs}
% ===========================================================================

\subsection{Traced vs.\ build-time}

A parameter is \textbf{traced} if it can change per call with zero recompile
(MCMC-safe), or \textbf{build-time} if changing it rebuilds and recompiles the
solver ($\sim$40--80 s).  The dividing line is whether the parameter changes the
shape of the computation (build-time) or only its numeric inputs (traced).

Every solver is a factory, \code{make_*_solver(...) -> (solve, grid)}, whose
arguments are build-time and whose \code{solve} is an ordinary \code{jax.jit}
function: build the solver once and call \code{solve} repeatedly.  \code{solve}'s
arguments are traced, so a new value reuses the compiled executable.  \code{jax.jit}
compiles one variant per call form, so the first call that passes its arguments
differently (a keyword not passed before, a positional instead of a keyword, a
Python number instead of a float64 array) compiles once more, and later calls of
that form reuse it.  In particular \code{me=None} and \code{xi=None} select
build-time constants (\code{M_ELECTRON}; the builder's degeneracy), and the first
call with a value compiles one traced variant.  The neutt banner
(\code{neutt.compiling()}, called first in each jitted \code{solve}) prints on every
compile; a call that prints nothing ran a cached executable.

\paragraph{Coupled-BBN \texttt{run(\dots)} arguments (traced, except the abundance-ODE controls)}
\begin{center}\small
\begin{longtable}{@{}R{0.16\textwidth}p{0.30\textwidth}p{0.40\textwidth}@{}}
\toprule
\textbf{parameter} & \textbf{meaning} & \textbf{notes} \\
\midrule
\endhead
\code{delta_neff_init} & extra inert dark radiation & LINX's \code{Delt_Neff_init} normalization; same convention standalone and coupled --- no glue conversion (Section~\ref{sec:dneffconv}) \\
\code{eta_fac} & baryon-to-photon rescale & $1=$ fiducial \\
\code{tau_n_fac} & neutron-lifetime rescale & $1=$ fiducial \\
\code{nuclear_rates_q} & nuclear-rate draws & $q\sim\mathcal N(0,1)$ per reaction \\
\code{me} & electron mass [MeV] & reaches bg $+$ weak rates $+$ transport coherently; unset keeps transport on \code{M_ELECTRON} \\
\code{T_init} & transport start [MeV] & full and qke modes \\
\code{T_anchor} & spectrum anchor [MeV] & semi mode only (default 4) \\
\code{xi} & per-flavour $\nu$ degeneracy $(\xi_e,\xi_\mu,\xi_\tau)$ & \code{run(xi=...)}/\code{solve(xi=...)} is traced: \code{None} (default) reuses the build-time value bit-for-bit, a passed triple recomputes the initial Fermi--Dirac spectrum and the noise floor inside the compiled solve with no recompile.  Only \code{full}/\code{qke} carry the asymmetry into $H$ and the $n\leftrightarrow p$ rates (the \code{semi} plasma is symmetric).  Validated end to end at $\xi_e=0.05$ (LO QED; matches an \code{rtol=1e-8} arbiter to $7\times10^{-7}$ and the degenerate-FD prediction to $0.8\%$), with no pinned test (Section~\ref{sec:standalonedemos}); run a tight-tolerance arbiter for any other or production $\xi$ \\
\code{rtol}, \code{atol}, \code{sampling_nTOp}, \ldots & abundance-ODE controls & forwarded to LINX \code{AbundanceModel} (its own defaults \code{rtol}$=10^{-6}$, \code{atol}$=10^{-9}$); these are the abundance-ODE tolerances, not the transport \code{atol}$=10^{-30}$ below \\
\bottomrule
\end{longtable}
\end{center}

\paragraph{Build-time arguments}
\begin{center}\small
\begin{longtable}{@{}R{0.17\textwidth}p{0.24\textwidth}p{0.40\textwidth}@{}}
\toprule
\textbf{parameter} & \textbf{where} & \textbf{meaning} \\
\midrule
\endhead
\code{nuclear_net} & \code{make_coupled_bbn} & LINX reaction network (default \code{key_recommended}: PRIMAT-2023 $+$ updated $n+p\to d+\gamma$) \\
\code{use_fe} & \code{make_coupled_bbn} & distorted $f_e$ in the $n\leftrightarrow p$ rates (default True) \\
\code{transport} & \code{make_coupled_bbn} & \code{"full"}/\code{"semi"}/\code{"qke"} (Section~\ref{sec:fullsemi}) \\
\code{solver_kwargs} & \code{make_coupled_bbn} & dict forwarded to the transport builder: \code{max_steps}/\code{implicit}/\code{pcoeff}/\code{icoeff}/\code{dcoeff}/\code{hybrid}/\code{hybrid_kwargs}/\code{use_renorm}/\code{renorm_pdep}/\code{tcm_cutoff}/\code{nuratiotol}/\code{phie}/\code{eta0}/\code{q_baryon}/\code{qed_nlo} (build-time; in full mode any other key raises; qke mode takes \code{qke_kwargs} instead) and \code{rtol}/\code{atol}/\code{dtmax}/\code{lna_max} (per-solve).  The only way to set the transport tolerances on the coupled path \\
\code{gate_width} & \code{make_coupled_bbn} & precision-gate smoothing width on the semi path (default 0 $=$ BURST hard gate) \\
\code{n_save} & \code{make_coupled_bbn} & saved output-trajectory nodes (default 256) \\
\code{use_sched} & all builders & lepton-kernel path (Section~\ref{sec:sched}) \\
\code{qed_nlo} & transport builders / \code{make_coupled_bbn} & NLO $\mathcal O(e^3)$ QED (Section~\ref{sec:nloqed}); default on in every builder (the low-level \code{thermo} QED functions default to \code{nlo=False}), \code{False} recovers the LO BURST value; in \code{make_coupled_bbn} tied to LINX's \code{NLO} \\
\code{use_renorm}, \code{renorm_pdep} & transport builders & finite-$T$ LO QED in the EoS (default on $=$ BURST-faithful) \\
\code{lag_tail_fix} & lepton kernels, \code{make_rhs}/\code{make_split_rates}/\code{make_decoupling_solver}/\code{make_hybrid_solver} & corrected wide-regime quadrature (default True; Section~\ref{sec:lepkernels}); False recovers the buggy bit-faithful BURST form \\
\code{phie} & transport builders & electron degeneracy $\mu_e/\Tpl$ (default 0 $=$ symmetric; a fixed float, or \code{"auto"} for the charge-neutrality $\phi(\Tpl)$, Section~\ref{sec:knownissues}) \\
\code{tcm_cutoff} & transport builders & transport-window low edge (BURST 0.015 MeV) \\
\code{nuratiotol} & transport builders & precision-gate threshold (BURST 3.0) \\
\code{lna_sw}, \code{n_head}, \code{head_power} & \code{make_hybrid_solver} & hybrid head geometry \\
\code{xi} & all builders & per-flavour $\nu$ degeneracy $(\xi_e,\xi_\mu,\xi_\tau)$; this sets the default, and the per-call \code{run(xi=...)}/\code{solve(xi=...)} override is traced (no recompile, see the traced table above) \\
\code{nbins}, \code{eps_range} & \code{make_eps_grid} & grid resolution/extent (101, 25 $=$ BURST's) \\
\bottomrule
\end{longtable}
\end{center}

\paragraph{Electron degeneracy \texttt{phie} vs.\ neutrino degeneracy \texttt{xi}: independent}
The electron degeneracy $\phi_e=\mu_e/\Tpl$ and the neutrino degeneracy $\xi$ are
different physical quantities and are not derived from one another.  $\xi=(\xi_e,\xi_\mu,\xi_\tau)$
is the neutrino lepton asymmetry, a free input you choose.  $\phi_e$ is instead
fixed by charge neutrality with the baryons, $n_{e^-}-n_{e^+}=q_{\rm baryon}\,n_b$,
i.e.\ it is a function of the baryon-to-photon ratio $\eta$, of $\Tpl$, and of
$\me$, not of $\xi$; it is
never inferred from the neutrino sector.  Their overrides are therefore separate:
\begin{itemize}[nosep]
\item \code{xi}: default symmetric $\xi=0$; set the factory default at build
  time, or override per call with \code{run(xi=...)} / \code{solve(xi=...)} (traced,
  no recompile).
\item \code{phie}: \code{0.0} (symmetric, default), a fixed float, or \code{"auto"}
  for the time-dependent charge-neutrality $\phi(\Tpl)$ solved inside the transport
  RHS (Section~\ref{sec:knownissues}); its $\Neff$ effect is $\ll10^{-7}$.
\end{itemize}
Both reach the Hubble rate and the $n\leftrightarrow p$ rates only in \code{full}/\code{qke}
mode (the \code{semi} path's external plasma is symmetric $e^\pm$).

\subsection{Tolerances and step control (change with care)}

These are the transport-solve tolerances.  The standalone solvers
(the \code{solve} of \code{make_decoupling_solver} and of \code{make_hybrid_solver}) take them directly; on the
coupled path they are set through \code{make_coupled_bbn(solver_kwargs={...})}.
The defaults are the tuned, validated state.  Each non-obvious choice:
\begin{itemize}[nosep]
\item \textbf{\code{rtol}$=10^{-6}$} (transport).  Counterintuitively $\sim2\times$
  faster than $10^{-4}$, which runs rejection-bound at the stability edge.  Its
  $\Neff$ differs from the $10^{-8}$ arbiter by $+5.7\times10^{-7}$.  Tightening to
  $10^{-8}$ costs only $\sim20\%$ (arbiter-grade); loosening buys nothing.
\item \textbf{\code{atol}$=10^{-30}$} $=$ pure-relative error norm, BURST-faithful.
  Never raise it.  At $10^{-7}$ the stiff $\bar\nu$-tail bins go
  error-uncontrolled $\to$ executable-dependent stalls and a silent
  $+1.5\times10^{-3}$ $\Neff$ bias (Section~\ref{sec:knownissues}).
\item \textbf{\code{dtmax}} $=0.08$ with \code{use_sched} (default), $0.15$ on the
  reference path (\code{make_decoupling_solver} only; the hybrid and coupled-path solves default to $0.08$).  The scheduled kernels reorder the gate sum at round-off and
  cascade Tsit5 at $0.15$.  Do not raise.
\item \textbf{\code{pcoeff,icoeff,dcoeff}}$=(0.4,0.3,0.0)$: a damped PI law.
  A pure-I controller oscillates at the explicit stability limit; the PI damping
  changes only the step-adaptation, never the enforced tolerance.
\item \textbf{\code{use_renorm}$+$\code{renorm_pdep}}$=$ True: finite-$T$ QED in
  the EoS (matches the collaborator oracle to $5.5\times10^{-7}$).  Turn it off
  only for the decoupling-only configuration ($\Neff=3.034$, Section~\ref{sec:pins}).
\item \textbf{No QKE coherence tolerance.}  The above are the diagonal
  transport-ODE tolerances.  The QKE coherent (off-diagonal) substep is an exact
  matrix exponential (Section~\ref{sec:qkesplit}) with no per-component
  \code{rtol}/\code{atol}; there is nothing trace-scaled to set on the real or
  imaginary coherences.  The only QKE-specific numerical control is \code{n_steps},
  the number of operator-split macro-steps, and (orthogonally) the grid size $N$; the
  result is converged in both (Section~\ref{sec:qke}).  A step too coarse for the plasma-temperature
  update ($\dd\ln a\gtrsim0.03$ over the full window, i.e.\ \code{n_steps}$\lesssim250$) drives
  $\Tpl$ unstable; the QKE solver then raises (it checks that $\Tpl$ never exceeds
  \code{T_init} and the final $\Tpl/\Tcm$ lies in $[0.99,1.45]$) instead of returning a
  meaningless $\Neff$.
\end{itemize}

\subsection{Starting transport at a lower temperature}

The transport start is decoupled from the BBN start: above the transport window
the history is clamped to equilibrium ($r_\nu=\rho_{\nu,\rm tot}/\rho_\gamma$ constant,
$\Tpl/\Tcm=1$), which is
exactly the physics there, so LINX's abundance evolution still begins at its own
high $T_{\rm start}$.  To start transport at, e.g., 4 MeV: \code{run(T_init=4.0)}
(traced, no recompile).  Rules of thumb:
\begin{itemize}[nosep]
\item \textbf{Floor $\approx3$--$4$ MeV.}  The equilibrium initial condition is only valid while
  weak interactions are still fast ($\nu$ decoupling $\sim2$--$3$ MeV); below that
  you erase real distortion physics.
\item \textbf{Shrink \code{lna_max} to keep the same end temperature}:
  $T_{\rm end}=\code{T_init}\,e^{-\code{lna_max}}$.  Unlike the traced
  \code{T_init}, \code{lna_max} is a build-time solver knob: pass it via
  \code{make_coupled_bbn(solver_kwargs={"lna_max":6.6})} (one recompile)
  together with \code{run(T_init=4.0)}.  The production end of
  $5.5\times10^{-3}$ MeV corresponds to the hybrid/coupled default
  $\code{lna_max}=7.5$ ($10\,e^{-7.5}$); note \code{make_decoupling_solver}'s
  \code{solve} defaults to $\code{lna_max}=7.0$ (ending slightly warmer, $\sim9\times10^{-3}$ MeV),
  so matching $5.5\times10^{-3}$ from $4$ MeV needs $\code{lna_max}\approx6.6$.
\item \textbf{Hybrid caveat:} \code{lna_sw} is $e$-folds from the start, so its
  $1.5$ default means ``hand off at $\Tcm=0.22\,\code{T_init}$''.  Below
  $\code{T_init}\approx6$~MeV just use \code{hybrid=False}; the stiff head the
  hybrid feeds on is mostly gone, and a mis-set \code{lna_sw} only costs the
  fallback price anyway.
\end{itemize}

% ===========================================================================
\section{Validation and testing}
\label{sec:validation}
% ===========================================================================

\subsection{The validation philosophy}

The port is gated at three levels (the BURST paper's ``checksums'' raised to a
first-class target):
\begin{enumerate}[nosep]
\item \textbf{Built-in sum rules (free).}  At equilibrium every net rate
  vanishes to the noise floor and the number/energy moments
  ($\int\eps^2\,\mathrm{net}$, $\int\eps^3\,\mathrm{net}$) are zero.  These
  validate each kernel standalone with no Fortran run
  (Section~\ref{sec:sumrules}).
\item \textbf{Module-level diffs against the Fortran oracle.}  The
  \code{fortran_reference/} snapshot compiles and runs, so each JAX module is
  checked against its Fortran counterpart at the intermediate level (a
  single $\dd f/\dd t$ at a fixed background, $\delta m_e^2$ from \code{renorm},
  $e^\pm$ $\rho$/$P$ at a given $T$), localizing any discrepancy to the module
  that introduced it.
\item \textbf{End-to-end pins.}  $\Neff^{\rm SM}=3.0433222$ (NLO, the default),
  $3.0442761$ (LO, \code{qed_nlo=False}; $3.0442252$ at \code{SIN2THW_BURST}), and $3.034029$
  (decoupling-only); the coupled $\Neff=3.0433222$ (network-independent, equal to the
  transport's); $\Yp\approx0.2470$,
  $\DoverH\approx2.44\times10^{-5}$ depend on \code{nuclear_net}.
\end{enumerate}

\subsection{Canonical numbers}
\label{sec:pins}

\begin{center}\small
\begin{tabular}{@{}R{0.30\textwidth}R{0.32\textwidth}R{0.30\textwidth}@{}}
\toprule
\textbf{quantity} & \textbf{value} & \textbf{configuration} \\
\midrule
$\Neff$ (full, NLO, default) & $3.0433222$ & NLO QED, \code{use_renorm}$+$\code{renorm_pdep}, 10 MeV \\
$\Neff$ (full, LO) & $3.0442761$ & \code{qed_nlo=False} \\
$\Neff$ (full, LO, BURST-faithful) & $3.0442252$ & \code{qed_nlo=False} at \code{SIN2THW_BURST}$=0.23$ \\
$\Neff$ (decoupling-only) & $3.034029^{\dagger}$ & \code{use_renorm=False} \\
$\Neff$ (momentum-indep.\ QED, LO) & $3.044313^{\dagger}$ & \code{renorm_pdep=False}, \code{qed_nlo=False} \\
$\Tpl/\Tcm$ asymptote (background only) & $1.401019$ & \code{evolve_background}, vs.\ $(11/4)^{1/3}=1.401020$ \\
$\Neff$ (background only) & $3.000007$ & \code{evolve_background} \\
coupled $\Neff$ (NLO) & $3.0433222$ & \code{make_coupled_bbn(solver_kwargs={"hybrid": False, "rtol": 1e-8})} \\
coupled $\Yp,\DoverH$ & $\approx0.2470,\ \approx2.44\!\times\!10^{-5}$ & network-dependent (\code{key_recommended}) \\
\bottomrule
\end{tabular}

\smallskip
\noindent{\footnotesize All rows use $\sin^2\theta_W=0.23122$ (\code{SIN2THW}) except the
BURST-faithful one, and come from the \code{rtol=1e-8} arbiter
(\code{scripts_qke/regen_pins_weinberg.py}).  The production angle raises every $\Neff$ here by
$+5.1\times10^{-5}$ relative to BURST's $0.23$ (\code{SIN2THW_BURST}).  The script sets XLA's
deterministic GPU flags, under which every pin is bitwise reproducible; with default autotuning
the LO pin can land up to $3\times10^{-7}$ higher.  $^{\dagger}$At $\sin^2\theta_W=0.23$ these two
configurations give $3.033977$ and $3.044262$.}
\end{center}

\subsection{Running the test suite}

\begin{lstlisting}[style=sh]
cd $BURST_ROOT && PYTHONPATH=$PWD \
  srun -n1 -c32 --gpus-per-task=1 \
  python -m pytest tests/ -q          # default deselects slow-marked tests (pyproject addopts)
\end{lstlisting}
The suite includes the foundation tests (grid/quadrature/moments vs.\ analytic)
and the per-kernel oracle comparisons (\code{test_nunu.py}, \code{test_nunubar.py},
\code{test_nuer1.py}, \code{test_nuer2.py}, \code{test_nuepma.py}, each diffed
against the Fortran goldens in the matching \code{tests/oracle_*} directory), plus
the equilibrium sum rules.  Three caveats: (i) the default invocation deselects
\code{slow}-marked tests, including the lepton-kernel \code{meps}-sweep
robustness test, which you run with \code{pytest tests/ -m slow}; (ii)
\code{tests/} also contains the QKE suites that a bare
\code{pytest tests/} will collect, so the total count is not a fixed SM number;
(iii) the end-to-end $\Neff$ pins ($3.0433222$ NLO / $3.0442761$ LO, coupled
$3.0433222$; Section~\ref{sec:pins}) are reference values from a tight-tolerance arbiter run (the
\code{rtol=1e-8} reference of Section~\ref{sec:knobs}), not part of the default
pytest suite.

\begin{center}\small
\fbox{\parbox{0.92\textwidth}{The first \code{pytest} invocation triggers a long
JAX compile ($\sim$10 min) before any test runs; this is not a hang.}}
\end{center}

% ===========================================================================
\section{Caveats and safeguards}
\label{sec:knownissues}
% ===========================================================================

This section is self-contained; each item is described in full here.  Each
caveat below is a numerical failure mode with its symptom and the safeguard the
default configuration applies, so a reader who sees the symptom, or turns the
safeguard off, knows what is happening.

\paragraph{Lepton-kernel quadrature at a regime boundary (\texttt{lag\_tail\_fix}, default on)}
Fortran BURST mis-resolves bins where an energy cut diverges and a fixed 5-point
Gauss--Legendre rule is applied over a near-infinite interval: the wide-regime
bins of \code{nuer1} and \code{nuepma} near a kinematic regime boundary (the
per-bin rate wrong by $2$--$4\times$ in a narrow window, the kernel discontinuous
at the boundary), and the lowest \code{nuer2} bins once $\Tcm\lesssim0.1$~MeV
(off by $20$--$100\%$).  By default (\code{lag_tail_fix=True},
Section~\ref{sec:lepkernels}) the port integrates those bins with a Gauss--Laguerre
tail rule instead; the end-to-end $\Neff$ shift is $\le8.4\times10^{-7}$ (below
$10^{-9}$ from \code{nuer2}).  \code{lag_tail_fix=False} gives BURST's (inaccurate)
bit-faithful form, used only by the Fortran-parity oracle tests.  The one case the tail
rule does not cover is the measure-zero exact boundary coincidence,
where the Fortran crashes and the port instead returns a finite zero-weight
tile, never sampled by the continuous solve.

\paragraph{Constant-folding of the precision-gate floor}
The equilibrium noise floor is input-independent, so under \code{jit} XLA would
evaluate it through its constant-folding path, which uses different
floating-point arithmetic and loses the $10^{-9}$ cancellation that defines the
floor; the gate would then invert and $\Neff$ collapse by $\sim0.03$ (in the
decoupling-only configuration, $3.034\to3.004$).  The port prevents this with an
\code{optimization_barrier} on the equilibrium occupation (Section~\ref{sec:gatenum}),
a runtime identity.

\paragraph{The \texttt{atol} floor}
Raising \code{atol} toward $10^{-7}$ leaves the stiff $\bar\nu$-tail bins
error-uncontrolled, giving executable-dependent \code{max_steps} stalls and a
silent $+1.5\times10^{-3}$ $\Neff$ bias.  The default $\code{atol}=10^{-30}$ (a
pure-relative norm) keeps every bin error-controlled.  Do not raise it
(Section~\ref{sec:knobs}).

\paragraph{NaN error estimates in the step controller}
A NaN error estimate slips past diffrax's controller guards and produces
$\dd t=\mathrm{NaN}$, a permanent rejection cascade (a quasi-random
\code{max_steps} stall).  The port's error norm, \code{nanrobust_rms_norm}, maps the
error norm's NaN to $+\infty$ so the library's designed reject-and-shrink path
handles it (bit-identical on every finite step).  It is the controller norm of every adaptive
transport solve (\code{make_decoupling_solver}, the hybrid tail, and the semi-mode
\code{make_spectrum_solver}); the reference helper \code{thermo.evolve_background} and LINX's
own background and abundance solves use diffrax's default norm.

\subsection{Features in BURST not present in the port}

For the SM/$\dNeff$ scope:
\begin{itemize}[nosep]
\item \textbf{Sterile neutrinos} ($n_\nu>3$).  The SM transport is hardwired to
  3 flavours $\times$ 2 cp.
\item \textbf{The nuclear reaction network} and \textbf{recombination /
  \code{sdratio}}.  Not ported; abundances come from LINX, and there is no
  recombination-era $\Neff$ (Section~\ref{sec:linx}).
\item \textbf{The $\phi_e$ degeneracy: trajectory ported, closure not.}  The
  default is $\phi_e=0$ (symmetric $e^\pm$).  The time-dependent charge-neutrality
  electron degeneracy is available as \code{phie="auto"}: it solves
  $\phi(\Tpl)$ from charge neutrality $n_{e^-}-n_{e^+}=q_{\rm baryon}\,n_b$ inside
  the transport RHS each step and feeds it to the lepton kernels and the QED EoS
  (full/qke modes).  It is valid while $\Tpl>\me/300\approx1.7$~keV ($\ln a\lesssim8.8$
  from 10~MeV): below that BURST's $E/T>300$ occupation cutoff removes every quadrature node,
  so \code{solve_phie} returns NaN and the solve fails rather than continuing on a wrong $\phi_e$.  Its $\Neff$ effect is below the $10^{-7}$
  reproducibility floor ($\phi\sim6\times10^{-10}$ at 10 MeV, rising to
  $\mathcal O(0.1)$ near the window bottom), so it is a completeness knob, not an
  $\Neff$ lever.  What is not ported is BURST's way of closing its
  own $\phi_e$ ODE, the \code{glag6} ($\dd\rho_e/\dd\phi_e$, BURST's
  \code{drhoepdphi}) EoS-derivative term; the port gets $\phi$ from exact charge
  neutrality directly instead.
\item \textbf{MPI parallelization.}  Replaced by the on-GPU tensor contraction.
\item \textbf{Integration in time $t$.}  Replaced by $\ln a$ with
  $\Tcm=\Tcm^{\rm init}e^{-\ln a}$ analytic.
\end{itemize}

\subsection{Troubleshooting quick reference}
\label{sec:troubleshoot}

The operational pitfalls scattered through the manual, collected for lookup; each
row points to the section with the full explanation.

\begin{center}\small
\begin{longtable}{@{}p{0.40\textwidth}p{0.54\textwidth}@{}}
\toprule
\textbf{symptom} & \textbf{cause / fix} \\
\midrule\endhead
x64 silently off; sum rules fail; results look float32-grade & \code{import neutt} before creating any JAX array; x64 is enabled at import (Section~\ref{sec:float64}). \\
\code{pytest} appears to hang for $\sim$10 min before any test runs & the one-time JAX compile, not a hang (Section~\ref{sec:validation}). \\
Every loop iteration recompiles ($\sim$40--80 s each; the neutt compile banner prints every time) & the solver is rebuilt inside the loop (each \code{make_*_solver} call returns a new \code{solve} that compiles afresh), or a build-time argument changes per call; build once outside the loop and vary only traced \code{solve} arguments (Section~\ref{sec:knobs}). \\
A scan call jumps from $\sim$4 s to $\sim$8 s & the hybrid qualification failed and the full-window fallback ran (the answer is still correct); switch a whole bad region to \code{hybrid=False} (Section~\ref{sec:whichsolver}). \\
A single retrace fires mid-scan & the first call that passes \code{solve}'s arguments in a new form, e.g.\ a value for an argument whose default is \code{None} (\code{me}, \code{xi}) or a keyword not passed before: one more variant is compiled once and reused (Section~\ref{sec:knobs}). \\
$\Neff$ biased by $\sim+1.5\times10^{-3}$, or executable-dependent stalls & \code{atol} set above $10^{-30}$; keep it at $10^{-30}$ (Sections~\ref{sec:odenum}, \ref{sec:knownissues}). \\
Energy-exchanging BSM gives the wrong coupled $\Neff$ & \code{transport="semi"} silently drops $\dot\rho_\nu$; use \code{transport="full"} for anything that exchanges energy with the neutrinos (Section~\ref{sec:fullsemi}). \\
Coupled run dies with a \code{diffeqsolve} ``maximum number of steps'' error & the LINX background solver's \code{max_steps} (default 512) is marginal for the NLO plasma EoS; pass \code{make_coupled_bbn(..., max_steps=8192)}.  This is the background ceiling; the transport solver's own \code{max_steps} is separate, via \code{solver_kwargs} (Section~\ref{sec:linx}). \\
\bottomrule
\end{longtable}
\end{center}

% ===========================================================================
\section{The QKE flavour-mixing sector}
\label{sec:qke}
% ===========================================================================

This is the Standard Model's own flavour-mixing refinement of the $\Neff$
calculation, not a BSM extension.  Three-neutrino oscillations redistribute the spectral
distortions across flavours as the neutrinos decouple, so the $\nu_e$ that drives
the $n\leftrightarrow p$ rates carries a flavour-redistributed
distortion.  It is the most elaborate piece of the SM transport: it replaces
the diagonal occupation $f[3,2,N]$ with the full flavour density matrix and
evolves the off-diagonal coherences.  The lineage is the \code{qburst} Fortran
code, but its end-to-end solve is not a trustworthy oracle (its operator-split
scheme is reported to be step-size unstable), so the port is validated instead by
exactness-by-construction (a matrix exponential that conserves neutrino number and
positivity by design), by unit oracles (the verbatim-\code{gfortran} weak-potential
and coherent-commutator checks), and, as the central check, by the
\textbf{mixing-off reduction}: with the coherences set to zero the whole
machinery reduces bit-identically to the validated classical transport of the
preceding sections.  The modules are under \code{neutt/transport/qke/}:
\code{coherent} (the density-matrix state, PMNS mixing, the Hamiltonian and
$-i[H,\rho]$), \code{nonadiabatic} (the coherent-step propagators and the moving
eigenframe), \code{collide_matrix}, \code{collide}, \code{collide_nue},
\code{collide_nunu} (the density-matrix collision), \code{full_coupled} (the
solver), and \code{solve} (a standalone Strang driver over a given background).

\subsection{The density-matrix state}

The classical state is the real occupation $f[\text{flavour}=3,\text{cp}=2,
\text{bin}=N]$, diagonal only.  The QKE generalizes it to a $3\times3$ complex
Hermitian flavour density matrix per (cp, bin):
\begin{equation}
  \rho[\text{cp}=2,\ \text{bin}=N,\ a,\ b]\in\mathbb C^{2\times N\times3\times3},
\end{equation}
with the diagonal real, $\rho[\text{cp},i,m,m]\equiv f[m,\text{cp},i]$ (the
classical occupations embedded exactly), and the off-diagonal complex
$\rho[\text{cp},i,a,b]$ ($a\ne b$) the \textbf{flavour coherences}, the new
degrees of freedom, zero at the initial condition (decohered flavour
states).  We write $\rho$ for the neutrino block (cp$=0$) and $\bar\rho$ for the
antineutrino block (cp$=1$); $\rho$ with a species subscript ($\rho_\gamma$,
$\rho_{e^\pm}$) remains an energy density.  The state evolves under
$i\,\dd\rho/\dd t=[H,\rho]+iC[\rho]$, with $H$ the flavour Hamiltonian ($H_{\rm flav}$
of Eq.~\eqref{eq:sysqke}; in this section $H$ denotes it, not the Hubble rate) and
$C[\rho]$ the collision term: the coherent commutator is integrated by an exact,
number- and positivity-preserving matrix exponential (no per-component ODE tolerance;
Section~\ref{sec:qkesplit}), the collision term by the classical kernels.

\begin{diffbox}
The batch axes (cp, bin) come first and the matrix (flavour) axes
last, so the commutator and \code{matmul} broadcast over (cp, bin) with no
transposes, a deliberate layout change from Fortran's \code{dens_mat(a,b,bin)}
(same content, reordered for JAX vectorization).  Helpers: \code{diag_to_rho},
\code{rho_to_diag}, \code{hermitize}, \code{commutator}; with the off-diagonal set
to zero, \code{rho_to_diag} recovers the classical $f[3,2,N]$ exactly.
\end{diffbox}

\subsection{The Hamiltonian and what is on by default}

The coherent RHS is $\dd\rho/\dd t|_{\rm coh}=-i[H,\rho]$ with, per (cp, bin) at
momentum $p_1=\eps_i\Tcm$ and $\mathrm{parfact}=\pm1$ for $\nu$/$\bar\nu$,
\begin{equation}
  H = H_{\rm vac} + H_{\rm mat}^{(1)} + H_{\rm mat}^{(2)}
    + H_{\rm si}^{(1)} + H_{\rm si}^{(2)}.
\end{equation}
\begin{center}\small
\begin{tabular}{@{}R{0.13\textwidth}R{0.44\textwidth}R{0.30\textwidth}@{}}
\toprule
\textbf{term} & \textbf{physics} & \textbf{production default} \\
\midrule
$H_{\rm vac}$ & vacuum oscillation $\propto m^2_{\rm weak}/2p_1$ & \textbf{ON} \\
$H_{\rm mat}^{(1)}$ & 1st-order MSW (Mikheyev--Smirnov--Wolfenstein) matter potential $\propto(n_{e^-}-n_{e^+})$ & OFF (\code{include_matter_asym=False}; suppressed by the baryon-to-photon ratio $\eta$) \\
$H_{\rm mat}^{(2)}$ & 2nd-order thermal MSW (N\"otzold--Raffelt) $\propto\rho_{e^\pm}$ & ON iff \code{include_matter=True} \\
$H_{\rm si}^{(1)}$ & 1st-order $\nu$--$\nu$ self-int $\propto(\rho-\bar\rho)$ & not included (zero) \\
$H_{\rm si}^{(2)}$ & 2nd-order $\nu$--$\nu$ self-int $\propto(\rho+\bar\rho)$ & not included (zero) \\
\bottomrule
\end{tabular}
\end{center}
The vacuum numerator $m^2_{\rm weak}=U M^2 U^\top$ uses the real PMNS
(Pontecorvo--Maki--Nakagawa--Sakata) mixing matrix
$U=R_{23}R_{13}R_{12}$ (no CP phase) with $\sin^2\theta_{12}=0.307$,
$\sin^2\theta_{13}=0.0212$, $\sin^2\theta_{23}=0.417$, and
$M^2=\mathrm{diag}(0,\Delta m^2_{21},\Delta m^2_{31})$ with $\Delta
m^2_{21}=7.53\times10^{-17}$ MeV$^2$ and $\Delta m^2_{31}=2.59\times10^{-15}$
MeV$^2$ (normal hierarchy); the trace$/3$ is removed (an $\propto$identity shift
drops out of the commutator but is removed explicitly to match Fortran
bit-for-bit) and $m^2_{\rm weak}(\bar\nu)=-m^2_{\rm weak}(\nu)$.  The two MSW and
two self-interaction terms carry the standard prefactors $\sqrt2\,G_{\rm
F}(n_{e^-}-n_{e^+})$ and $-\tfrac{8\sqrt2}{3}G_{\rm F}E_\nu\rho_{e^\pm}/m_W^2$
(matter), $\sqrt2\,G_{\rm F}(\rho-\bar\rho)$-type and
$-\tfrac{8\sqrt2}{3}(G_{\rm F}/m_Z^2)(\rho+\bar\rho)$-type (self-interaction).
Production $\Neff$ / coherence solves activate $H_{\rm vac}$, plus the thermal MSW
term with \code{include_matter=True} and the asymmetric one with
\code{include_matter_asym=True}; the two $\nu$--$\nu$ terms are not included.
\code{coherent.hamiltonian} is oracle-validated (full commutator to rel $10^{-15}$).

\paragraph{The $\nu$--$\nu$ terms are not included}
The solver carries no $\nu$--$\nu$ refractive term: $H_{\rm si}^{(1)}=H_{\rm
si}^{(2)}=0$.  Both are nonzero in the symmetric universe.  $H_{\rm
si}^{(2)}\propto\Sigma w\eps^3(\rho+\bar\rho)$ ($w$ the grid quadrature weights) is
sourced only by the momentum-dependent flavour misalignment with the mean field (the
commutator of the mean field with $\rho(p)$ vanishes when every momentum shares one
flavour structure).  $H_{\rm si}^{(1)}\propto\Sigma w\eps^2(\rho-\bar\rho)$ does not vanish
at $\xi=0$ either: the solve's symmetric state is $\bar\rho=\rho^\ast$, not
$\bar\rho=\rho$, so $H_{\rm si}^{(1)}\propto\Sigma w\eps^2\,{\rm Im}\,\rho$, the
mean-field term of the flavour-pendulum system (the ATAO, adiabatic transfer of
averaged oscillations, treatment of Froustey
et al., in which $\rho$ is real, drops it).  Each moves $\Neff$ by
$|\delta\Neff|\lesssim3\times10^{-7}$ (measured with the vacuum and with the
thermal-MSW Hamiltonian), below the precision of every result in this manual; the
coherence frames are those of the model without $\nu$--$\nu$ refraction.

\paragraph{The induced $\nu$--$\bar\nu$ asymmetry}
$H_{\rm mat}^{(1)}\propto\sqrt2\,G_{\rm F}n_p$ (the ``solar'', linear-$G_{\rm F}$
MSW term, with $n_p=q_{\rm baryon}n_b\propto\eta$ the baryon charge the $e^\pm$ gas
balances) is the one $\xi=0$ term that breaks $\nu\leftrightarrow\bar\nu$ in
the diagonal.  With it off, $m^2_{\rm weak}\!\to\!-m^2_{\rm weak}$ (vacuum) and the
$\bar\nu$ sign of the thermal term give $H_{\bar\nu}=-H_\nu$, hence
$\bar\rho=\rho^\ast$ and equal diagonals; $H_{\rm mat}^{(1)}$ alone carries
no $\bar\nu$ flip, the complex-conjugate (Sigl--Raffelt) equivalent of the lab
frame where the asymmetric potential flips and the thermal one does not
(N\"otzold--Raffelt), so it seeds a genuine flavour asymmetry
$\rho_{aa}(\nu)-\rho_{aa}(\bar\nu)$ by acting on the oscillation coherence.  It is
switched on by \code{include_matter_asym} (default off, which leaves the solve
bit-identical), with a per-call traced \code{matter_asym_scale} on $n_p$ ($0$ null,
$1$ physical $\eta$, ${>}1$ source sweep, all on one compile).  Measured: the
decoupling-era ($\Tcm\!\sim\!2$ MeV), null-subtracted asymmetry is linear
in $n_p$ over scale $1\to100$ (saturating by $1000$), with physical-$\eta$
magnitude $A_e/N_e\sim2\times10^{-6}$ ($A_e$ the grid-integrated $\nu_e\!-\!\bar\nu_e$
number difference, $N_e$ the $\nu_e$ number) and
$|\delta f_e|/f_e\sim\text{few}\times10^{-4}$ (local), $\dNeff=-1.4\times10^{-7}$,
negligible for $\Neff$ but a nonzero, $\eta$-controlled induced asymmetry.
Raw peak-over-time measures are contaminated at the physical scale by the
late-time lepton-kernel tail amplification of Section~\ref{sec:lepkernels};
null-subtract / restrict to the decoupling era for the clean signal.

\paragraph{Lower-level API (\texttt{coherent})}
\begin{lstlisting}
diag_to_rho(f) -> rho[2,N,3,3]          # f[3,2,N] embedded; coherences 0
rho_to_diag(rho) -> f[3,2,N]            # the real diagonal; coherences dropped
hermitize(rho)                          # (rho + rho^dagger)/2 over the flavour axes
commutator(A, B)                        # A@B - B@A, batched over leading axes
pmns_matrix(s12=sqrt(SIN2_THETA12), s13=sqrt(SIN2_THETA13), s23=sqrt(SIN2_THETA23),
            nnu=3) -> U (3,3)           # real R23 R13 R12, from the sines of the angles
m2_mass_matrix(dm2_21=DELTA_M2_21, dm2_31=DELTA_M2_31, normal_hierarchy=True,
               nnu=3) -> M2 (3,3)       # diag(0, dm2_21, +-dm2_31) [MeV^2]
m2_weak(s12=..., s13=..., s23=..., dm2_21=..., dm2_31=..., normal_hierarchy=True,
        nnu=3) -> (3,3,2)               # U M2 U^T - Tr/3; [...,1] = -[...,0] (nubar)
matter_terms(therm_corr, npro, nflavor=3) -> (mata, mats)   # the two MSW potentials (ee entry)
self_interaction_terms(tcm, nudiff, nusum) -> (sia, sis)    # the nu-nu potentials
hamiltonian(eps, tcm, m2_weak, mata_mat, mats_mat, sia_mat, sis_mat) -> H[2,N,3,3]   # [s^-1]
coherent_rhs(rho, eps, tcm, m2_weak, *, therm_corr=0.0, npro=0.0,
             nudiff=None, nusum=None, hermitian_project=True)   # -i[H, rho]; 0 at eps = 0
\end{lstlisting}
\code{tcm} is $\Tcm$ [MeV], \code{therm_corr} the $e^\pm$ energy density [MeV$^4$] of
$H_{\rm mat}^{(2)}$ and \code{npro} the proton density $n_p$ [MeV$^3$] of $H_{\rm mat}^{(1)}$;
\code{nudiff}, \code{nusum} are the $\nu-\bar\nu$ and $\nu+\bar\nu$ moments of
\code{solve.nudiff_nusum} (below); the solvers pass zero $\nu$--$\nu$ potentials.

\subsection{The operator split: \texttt{hybrid} by default}
\label{sec:qkesplit}

The solve uses operator splitting: rather than integrating the coherent
(Hamiltonian) and collision parts of $\dd\rho/\dd s$ ($s\equiv\ln a$) together, each
of the $n_{\rm steps}$ (\code{n_steps}) fixed $\ln a$ macro-steps, of duration $\dd t$
in seconds, advances them in alternation as two substeps, which lets each use
the integrator best suited to it (an exact propagator for the stiff coherent
oscillation, the L-stable exponential-midpoint for the stiff collision).
\begin{itemize}[nosep]
\item \textbf{Collision substep.}  The production collision advances the
  diagonal with the validated classical kernels in affine form
  (\code{make_split_rates} $\to A-Bf$ over the five kernels $+$ the $\dd\Tpl$
  ODE), with the L-stable exponential-midpoint step.  How the collision acts on
  the coherences is the \code{collision} option of
  Section~\ref{sec:qkecollision}: by default a scalar damping, or the
  density-matrix (Sigl--Raffelt) operator of \code{collide.py},
  \code{collide_nue.py} and \code{collide_nunu.py}.
\item \textbf{Coherent substep.}  $\rho\leftarrow
  \mathrm{expm}(L_{\rm coh}\,\dd t)\rho$ per bin, with $L_{\rm coh}=-i[H,\cdot]-
  \Lambda_{\rm off}$ the $9\times9$ Liouvillian ($9\times9$ because the $3\times3$
  density matrix is flattened to a 9-vector, turning the commutator $[H,\cdot]$
  into an ordinary matrix$\times$vector product) and $\Lambda_{\rm off}$ an
  off-diagonal-only dephasing: it exponentially damps the off-diagonal
  coherences at $\Gamma_{ab}=\tfrac12(\Gamma_a+\Gamma_b)$ ($\Gamma_a$ the total
  collision rate of flavour $a$, Section~\ref{sec:qkecollision}) while, being trace- and
  positivity-preserving, never pushing an occupation negative.  Two
  choices matter: the diagonal is undamped here (its collisional loss
  already lives in the collision substep; damping it again double-counts), and
  $L_{\rm coh}$ is a Lindblad pure-dephasing $\oplus$ unitary channel, so
  $\mathrm{expm}$ \textbf{conserves $\sum_a\rho_{aa}$ (number) and keeps the
  diagonal $\ge0$ (positivity) by construction}, which is why it cannot
  trip the collision kernel's occupation-floor cliff.  (\code{max_squarings=40},
  not the JAX default 16: at late times the rotation $\omega\,\dd t$ ($\omega$ the
  largest eigenvalue splitting of $H$) reaches
  $\sim10^{9}$ rad/step, which the default silently clamps $\to$ NaN.  At that
  phase the squaring is not exactly unitary in floating point; see
  \code{hybrid} below.)
\end{itemize}
The \code{expm} integrator is thus the dense $9\times9$ matrix exponential
$\mathrm{expm}(L_{\rm coh}\,\dd t)$ of the dissipative Liouvillian, BURST's
corotating-frame Magnus idea realized as a single exact interaction-picture
propagator, with no tuned subtraction and no per-bin frame-switch logic, exact
for any stiffness up to a round-off that grows with the phase per step (below).
It is also tolerance-free: there is no per-component
\code{rtol}/\code{atol} on the coherences; the stiff phase $\omega\,\dd t$ is
exponentiated exactly rather than error-controlled, so a small imaginary
off-diagonal cannot drift against a loose tolerance the way it does in an
adaptive operator-split integrator (the failure mode of a corotating scheme such
as \code{qburst}'s; see ``Watch the magnitude'' in Section~\ref{sec:qkeresults}).
(The opt-in Strang split below instead realizes the
unitary part as a $3\times3$ Hermitian eigendecomposition
$U=Ve^{-i\Lambda\dd t}V^\dagger$ ($V$, $\Lambda$ the eigenvectors and eigenvalues of
$H$; this $U$ is the step unitary, not the PMNS matrix; \code{nonadiabatic.propagate_coherent_H}), plus a
Hadamard dephase, dropping the $9\times9$ \code{expm} and its
\code{max_squarings=40}.)

\paragraph{Accuracy of the coherent integrators; \texttt{hybrid}}
Scaling-and-squaring is not exactly unitary at large phase: per step the dense
\code{expm} moves occupations, per-bin number included, by
$\approx10^{-17}\,\omega\,\dd t$, with $\omega$ the largest eigenvalue splitting of
$H$.  At low $\eps$ and late times $\omega\,\dd t$ reaches $10^{8}$--$10^{10}$, and the
accumulated round-off makes the mixing-on solve ill-conditioned: a $10^{-12}$
change of \code{T_init} moves $\Neff$ by $\sim10^{-6}$ and the low-$\eps$
($\eps\le1$) flavour-differential distortion $\delta_e-\delta_\mu$ ($\delta_a$ the
fractional distortion $\delta f$ of flavour $a$) by $\sim40\%$, and
$\dNeff$ carries a round-off bias of $\approx+3\times10^{-6}$ (damping, $N=101$,
$n_{\rm steps}=2000$).  The Strang split (\code{coherent_method="strang"}:
flavour-basis dephasing as half-step kicks around an exact unitary-plus-source
step) has no round-off, but applies each kick between rotations of
$\omega\,\dd t\gg1$ rad.  That split error costs $1.5\times10^{-6}$ in $\dNeff$ at
$n_{\rm steps}=2000$, spreads the low-$\eps$ $\delta_e-\delta_\mu$ by 5\% across step
counts and grids, and leaves a spurious rotating low-$\eps$ coherence in the mass basis of
magnitude $9\times10^{-8}$, where the dense \code{expm} gives $1.6\times10^{-10}$ for the same
coherence.

\code{coherent_method="hybrid"} (\code{nonadiabatic.coherent_step_sourced_hybrid})
has neither error.  Per bin it keeps the dense \code{expm} where
$\omega\,\dd t\le10^{3}$ (round-off $\lesssim10^{-14}$).  Elsewhere, if
$\Gamma_{\max}\dd t\le1$ ($\Gamma_{\max}$ the bin's largest decoherence rate), it takes a
secular step in the eigenbasis of $H$: each
component's rotation and own decay are exact phases (unitary at any
$\omega\,\dd t$), the population block of the transformed dephasing is exponentiated
exactly, and the oscillating (non-secular) couplings enter averaged over the step
(first Magnus term in the interaction picture).  Its errors are
$\mathcal O(\Gamma\dd t\,(\Gamma/\omega)^2)$ on populations and
$\mathcal O(\Gamma\dd t\,\Gamma/\omega)$ on coherences, and both branches are
accurate at the switch (moving it to $10^{4}$ changes $\dNeff$ by $\sim10^{-11}$).
Measured (damping, $N=101$/$201$, $n_{\rm steps}=2000$--$32000$): the $10^{-12}$
nudge moves $\Neff$ by $3\times10^{-13}$, $\delta_e-\delta_\mu$ at $\eps\le1$ agrees
to 0.3\% across all six runs (\code{expm}: 41\%), per-bin number is conserved to
$10^{-13}$, and the rotating coherence equals \code{expm}'s.  Under
\code{collision="matrix_full"} the nudge moves $\Neff$ by $3\times10^{-14}$
(\code{expm}: $1.1\times10^{-6}$) and the low-$\eps$ spread is 0.8\%
(\code{expm}: 22\%).  It costs $\sim$17\% more wall time in damping mode and
$\sim$8\% in \code{matrix_full}.  \code{hybrid} is the default.

\paragraph{The moving eigenframe of the step}
The coherent step freezes $H$ at the step midpoint.  Where the thermal potential
rotates the eigenbasis of $H$ as $T$ falls (most at high $\eps$, where
$V_{\rm th}/\omega_{\rm vac}\propto\eps^2T^6$, with $V_{\rm th}$ the thermal MSW
potential and $\omega_{\rm vac}$ the vacuum splitting), a frozen $H$ changes direction at
every step boundary and leaves the coherence behind by the rotation angle, at an
effectively random phase once $\omega\,\dd t\gg1$.  Those kicks random-walk into
bin-to-bin structure in the high-$\eps$ coherences that falls only as
$1/n_{\rm steps}$: in the non-adiabaticity measure
$\Theta=\|[H_{\rm vac},\rho]\|^2/(\|H_{\rm vac}\|^2\|\rho\|^2)$ (Grohs et al.;
Frobenius norms, unit-trace $\rho$), $\Theta\approx10^{-13.5}$ at $\eps=20$--$25$ for
$n_{\rm steps}=6000$, against $<10^{-16}$ converged.  Each step therefore runs in
the frame that carries the eigenbasis of $H(t_{\rm start})$ onto that of
$H(t_{\rm end})$.  With the frame generator $K$ (Hermitian) defined by
$e^{-iK\dd t}=V_{\rm end}V_{\rm start}^\dagger$ (eigenvectors
in matching order, phase-aligned; \code{nonadiabatic.eigenframe_generator}) and
$\tilde R=e^{-iK\dd t/2}$,
\begin{equation}
  \rho(\dd t)=\tilde R\;\mathrm{step}\big(H_{\rm mid}-K,\ \tilde R\rho\tilde R^\dagger\big)\;\tilde R^\dagger ,
\end{equation}
where $H_{\rm mid}$ is $H$ at the step midpoint, $\mathrm{step}$ is the coherent
substep above, and $-K$ is the geometric term of the rotating frame; in that frame the
eigenbasis is nearly fixed and the coherence is transported adiabatically.  The
frame coincides with the flavour basis at the step midpoint, where $\Gamma$ and the
collision source are evaluated.  $K=0$ when the eigenvectors do not move (vacuum),
which recovers the plain step bit-for-bit, and, as a guard, where the frame would
turn by more than $0.5$~rad in one step (an eigenvalue crossing reorders the
eigenvectors).  Measured with the thermal potential: at $n_{\rm steps}=6000$ no
high-$\eps$ bin has $\Theta>10^{-15}$ (with the frozen step all twenty at
$\eps=20$--$25$ do, and the frozen step is still worse at $n_{\rm steps}=150\,000$);
between 0.5 and
5~MeV the coherences match the frozen step at $n_{\rm steps}=150\,000$ to
$|\Delta\log_{10}\Theta|\le0.008$; and $\Neff$ moves by $+5.0\times10^{-8}$ ($n_{\rm steps}=2000$) and
$+1.0\times10^{-8}$ (6000), the frozen step's error shrinking with the step.  The unit test is the
adiabatic theorem: an eigenstate of a slowly rotating $H$ leaks
$3.8\times10^{-9}\approx(|K|/\omega)^2$ with the frame and $7.7\times10^{-6}$ with
the frozen step.

\paragraph{Lower-level API (\texttt{nonadiabatic}, \texttt{solve})}
Matrices are $(\ldots,3,3)$ complex, batched over leading axes such as (cp, bin); $H$ and the
rates are in s$^{-1}$, \code{dt} in s, \code{gamma_diag} is $\Gamma_a$ $(\ldots,3)$ and \code{src}
the sourced step's $s_{\rm coll}$.
\begin{lstlisting}
# the coherent substep (full_coupled's coherent_method) and the moving frame
coherent_liouvillian(H, gamma_diag) -> L_coh (...,9,9)            # -i[H,.] - Lambda_off, row-major vec
coherent_step(H, gamma_diag, rho, dt, max_squarings=40)            # expm(L_coh dt) rho
coherent_step_sourced(H, gamma_diag, rho, dt, src, max_squarings=40)   # ... + dt phi1(L_coh dt) src
coherent_step_strang(H, gamma_diag, rho, dt)                       # D(dt/2) U(dt) D(dt/2)
coherent_step_sourced_strang(H, gamma_diag, rho, dt, src)          # "strang"
coherent_step_sourced_secular(H, gamma_diag, rho, dt, src)         # secular step, omega dt >> 1
coherent_step_sourced_hybrid(H, gamma_diag, rho, dt, src, wdt_switch=None)   # "hybrid"
HYBRID_WDT_SWITCH = 1.0e3                                          # its default omega_max dt switch
eigenframe_generator(H_start, H_end, dt, max_angle=0.5) -> K       # the moving-eigenframe generator
unitary_exp(K, t)                                                  # e^{-iKt}
# exact unitary, frozen-diagonal coherence, collision step, diagnostics
eigendecompose(H) -> (evals, V)            # of the hermitized H; evals ascending, V's columns
magnus_propagator(evals, V, dt) -> U       # V e^{-i Lambda dt} V^dagger
propagate_coherent_H(rho, H, dt)           # U rho U^dagger, U = e^{-iH dt}; U = I where H = 0
liouvillian(H, gamma_diag)                 # -i[H,.] - {diag(Gamma/2), .}: damps the diagonal too
offdiag_blocks(H, gamma_diag) -> (L_oo, L_od)        # of liouvillian, (...,6,6) and (...,6,3)
qss_offdiag(L_oo, L_od, rho_diag) -> (...,6)          # -L_oo^{-1} L_od rho_diag (a diagnostic)
evolve_offdiag(L_oo, s, rho_off, dt, max_squarings=40)    # exact step of rho_off' = L_oo rho_off + s
rho_from_offdiag(rho_off, rho_diag) -> rho            # the Hermitian (...,3,3)
off_diag_amps(rho, V=None, H=None)                    # mass-basis amplitude envelope (qburst's)
lstable_collision_step(rho, A_rate, B_rate, dt)       # exact step of drho/dt = A - (B rho + rho B)
strang_step(rho, H_fn, collision_fn, t, dt)           # half collision, exact unitary, half collision
# solve.py: a standalone Strang driver over a given background
make_coherent_H_fn(grid, m2_weak, bg_fn) -> H_fn(t, rho=None)   # bg_fn(t) -> {"T_cm", "therm_corr", "npro"}
evolve(rho0, t_grid, H_fn, collision_fn=None) -> (M,2,N,3,3)     # strang_step over t_grid [s]
number_density(grid, rho, T_cm), energy_density(grid, rho, T_cm)   # moments of rho [MeV^3], [MeV^4]
nudiff_nusum(grid, rho, diff_pow=2, sum_pow=3) -> (nudiff, nusum)  # the nu-nu source moments
coherence_norm(rho)                                                # sum over a != b of |rho_ab|
\end{lstlisting}
\code{strang_step} and \code{solve.evolve} take explicit half steps of the collision, so they suit
only a resolved (non-stiff) collision; the production solve is \code{full_coupled}.

\subsection{The collision acting on the coherences}
\label{sec:qkecollision}

Two options set how the collision substep reaches the density matrix:
\code{collision_coupling} (how the classical diagonal update is handed to the
coherent step) and \code{collision} (which operator acts on the coherences).

\paragraph{\texttt{collision\_coupling}: \texttt{"sourced"} (default) or \texttt{"jump"}}
Under \code{"sourced"} the collision's change of the diagonal enters the coherent
step as a constant rate $s_{\rm coll}=(f_{\rm coll}-f)/\dd t$ ($f_{\rm coll}$ the
diagonal after the collision substep) inside the exponential
(\code{nonadiabatic.coherent_step_sourced} and its \code{hybrid}/\code{strang}
counterparts), so the diagonal ramps during the step and the off-diagonal's
quasi-steady target $-L_{oo}^{-1}L_{od}\rho_{\rm diag}$ moves continuously
($\rho_{\rm diag}$ and $\rho_{\rm off}$ the diagonal and off-diagonal entries of $\rho$,
$L_{oo}$ and $L_{od}$ the off-diagonal$\to$off-diagonal and diagonal$\to$off-diagonal
blocks of $L_{\rm coh}$); for a constant $(L_{\rm coh},s_{\rm coll})$ the step's fixed
point is exactly $-L_{\rm coh}^{-1}s_{\rm coll}$, so a solve sitting on
the quasi-steady state stays on it.  \code{"jump"} applies the collision first and
runs the coherent step on the updated state: the target then jumps once per step
and $\rho_{\rm off}$ trails it by $\mathcal O(\dd t)$, a lag that is generated
across the transition epoch and freezes in once $\Gamma\to0$.  It puts
${\rm Im}\,\rho_{e\tau}$ at $2.5\times10^{-7}$ ($n_{\rm steps}=10^4$) where the
quasi-steady value is $\sim4\times10^{-15}$, converging only at first order;
\code{"sourced"} gives $5.5\times10^{-11}$, $7.9\times10^{-12}$, $2.4\times10^{-12}$ at
$n_{\rm steps}=10^4$, $5\times10^4$, $10^5$, with $\Neff$, \code{num_drift} (the
per-step number-conservation residual) and \code{min_f} (the minimum occupation)
unchanged.  $\Neff$ barely sees the difference (it rides on the
diagonal, which the collision dominates).

\paragraph{\texttt{collision}: the operator on the coherences}
\begin{itemize}[nosep]
\item \code{"damping"} (default): the scalar model
  $\dd\rho_{ab}/\dd t=-\tfrac12(\Gamma_a+\Gamma_b)\rho_{ab}$, with $\Gamma_a$ the
  ungated total collision throughput of flavour $a$ (\code{assemble.decoherence_rate};
  \code{gamma_ungated=False} substitutes the gated transport rate, a comparison
  mode).  It is exact as a loss but drops the coherence gain, the
  momentum-nonlocal repopulation of $\rho_{ab}(p)$ by the other momenta.
\item \code{"matrix"}: all five SM channels act through the density-matrix
  (Sigl--Raffelt) operator $\dd\rho/\dd t|_{\rm coll}=A_{\rm gain}-\{B_{\rm
  loss},\rho\}$: $\nu\bar\nu\to e^+e^-$ (\code{collide.py}), $\nu$--$e^\pm$
  scattering (\code{collide_nue.py}) and $\nu$--$\nu$, $\nu$--$\bar\nu$
  (\code{collide_nunu.py}).  They share the classical kernels' energy and angular
  quadrature, with the blocking factors $(1-f)$, $f$ promoted to $(I-\rho)$, $\rho$
  inside anticommutators, so $A_{\rm gain}$ and $B_{\rm loss}$ are Hermitian, and
  $B_{\rm loss}\succeq0$ on a physical state ($0\preceq\rho\preceq I$) except at
  \code{nuepma}'s lowest bin below $\Tcm\approx0.1$~MeV, where BURST's quadrature leaves a
  small negative eigenvalue (Section~\ref{sec:lepkernels}).  On a diagonal $\rho$ each reduces to its classical
  kernel, and $2\,{\rm diag}(B_{\rm loss})$ equals the ungated throughput, so each
  channel's scalar damping is replaced by $2B_{\rm loss}$ and the remainder of the
  operator enters the sourced step as a source: the off-diagonal right-hand side is
  $\Tcm^5\,\mathrm{net}_{ab}$ with no truncation in the coherence order.  The
  diagonal is still advanced by the classical kernels.
\item \code{"matrix_full"}: as \code{"matrix"}, and the diagonal is taken from the
  same operator, $A_{{\rm gain},aa}-\{B_{\rm loss},\rho\}_{aa}$: the classical
  affine $A-Bf$ plus the $\mathcal O(|\rho_{\rm off}|^2)$ coherence feedback the
  classical kernels cannot see.  Its affine halves go to the same
  exponential-midpoint step, so the two modes differ in the operator only.  The
  precision gate acts on the diagonal alone, with the classical gate decision (the
  same open bins); gating the coherences would switch their damping off at
  freeze-out.
\item \code{"matrix_nue"}: the three $\nu$--$e$ channels on the matrix operator,
  $\nu$--$\nu$ and $\nu$--$\bar\nu$ on the damping model (a comparison mode).
\end{itemize}
The matrix modes need \code{collision_coupling="sourced"} (the gain is a source,
and the jump path has nowhere to put it) and \code{gamma_ungated=True}; the
factory rejects other combinations.  The matrix kernels take the same off-grid
tail closure as the classical ones (Section~\ref{sec:qkecap}).  Measured (vacuum
mixing, $N=101$, $n_{\rm steps}=2000$; $\Neff^{\rm damping}=3.043945$):
\begin{center}\small
\begin{tabular}{@{}lrl@{}}
\toprule
\textbf{\code{collision}} & \textbf{$\Neff-\Neff^{\rm damping}$} & \textbf{notes} \\
\midrule
\code{"matrix_nue"} & $-4.5\times10^{-7}$ & the $\nu$--$e$ gain alone \\
\code{"matrix"} & $-9.2\times10^{-5}$ & $\nu$--$\nu$ gain dominates; $|\rho_{\mu\tau}|$ $+35\%$ \\
\code{"matrix_full"} & \code{"matrix"} $\pm5\times10^{-7}$ & feedback $\lesssim10^{-5}$ of the throughput \\
\bottomrule
\end{tabular}
\end{center}
The $\nu$--$\nu$ gain is 5--15 times the $\nu$--$e$ one, and the
\code{"matrix"} shift is converged to $1.5\%$ in step count ($-9.4\times10^{-5}$ at
\code{n_steps=6000}) and $0.7\%$ in grid ($N=101\to201$).  The \code{"matrix_full"} shift is not resolved by the
fixed-step solve: it scatters at the $\pm3\times10^{-7}$ level between step counts,
so it is quoted as the bound $|\delta\Neff|\lesssim5\times10^{-7}$.  Number
conservation is unchanged (largest \code{num_drift} $6\times10^{-14}$ in every mode).

\subsection{The off-grid tail closure}
\label{sec:qkecap}

The QKE solve needs an off-grid tail closure, and getting it right
is the most important numerical ``why'' in this sector.

\paragraph{The root cause: the lepton kernels amplify an off-grid tail}
The $\nu$--$\nu$ / $\nu$--$e$ kernels read an off-grid source occupation
$f(p_4)$, $p_4=p_1+E_2-E_3$, which reaches the Gauss--Laguerre tail past the grid
edge ($25<p_4<50$).  There the order-5 log-occupation polynomial
extrapolates tens of grid-steps past the data, so a tiny spectral
perturbation is amplified into an unphysical occupation.  In the QKE
solve the coherent back-reaction seeds exactly such a tail perturbation
$\to$ an $f(p_4)$ spike $\to$ pure gain $\to f$ pumped toward 1 $\to$ fake energy
$\to\Tpl$ corrupts $\to\Neff\to0/$NaN.  The culprit is the SM lepton kernels
(\code{nuer2}/$\nu\nu$) amplifying a coherence-induced super-thermal tail
perturbation, the lepton-kernel boundary fragility
(Section~\ref{sec:knownissues}), H-theorem-violating on non-smooth spectra.  The
coherence engine is sound; the collision amplifies a seed that should have been
negligible, so a well-posed solve needs a closure that bounds the overshoot
without erasing real signal.

\paragraph{The default: the parameter-free \texttt{srcfix} closure}
The QKE solver defaults to \code{tail_closure="srcfix"}, a parameter-free
pair of closures.  \code{fd_decay} carries the live grid-edge $f/f^{\rm eq}$
ratio outward along the FD shape (bounding the polynomial overshoot with no
tunable constant while preserving a genuine super-FD edge), and
\code{monotone_read} clamps the on-grid interpolant to its two-node bracket,
killing the odd--even Nyquist mode the kernels self-amplify.  With it the collision
damps the spurious spike rather than amplifying it.  Because the same closure
acts in both mixing branches it cancels in the mixing-on$-$mixing-off $\dNeff$:
it reproduces that differential (vacuum $+7.6\times10^{-4}$, matter
$+6.7\times10^{-4}$) and keeps the classical (mixing-off) reduction bitwise
unchanged.

\paragraph{The opt-in FD-envelope cap}
The alternative, \code{tail_closure="cap"}, caps the off-grid source to a
Fermi--Dirac envelope $f(p)\le C\cdot f^{\rm eq}(p)$ with $C=10$ (i.e.\ $\log
f\le\log C-\mathrm{logaddexp}(p,0)$, applied with \code{jnp.minimum}) on both
\code{full_coupled} kernel edges, including the non-cancelling
\code{decoherence_rate}, which must be capped explicitly or it silently
contaminates $\dNeff$, and the three scheduled lepton kernels.  The mixing-on $\Neff$
is flat in $C$ across $\sim[5,40]$.  The cap is the numeric \code{fd_cap_logC}
(default $+\infty$: $\mathrm{minimum}(x,+\infty)=x$, a bitwise no-op,
jit-folded at zero hot-path cost), so the classical SM transport is untouched and
BURST-faithful.  \code{srcfix} is preferred as the default because it is strictly
more correct at a degenerate spectrum: at $\xi=3>\ln C$ the absolute
$C\!\cdot\!f^{\rm eq}$ cap clips the genuine degenerate tail ($\Neff=13.1$,
wrong) whereas the edge-anchored closure preserves it ($\Neff=20.2$).  A naked,
no-closure solve diverges to NaN, so a closure is always load-bearing.

\begin{diffbox}
The off-grid closure is a documented improvement over BURST's bare
extrapolation, applied only where it is physically valid (the near-thermal QKE
tail) and held off where it would move
the SM collaborator oracle.  Tail-accuracy ordering: \code{srcfix} $>$ FD-envelope
$>$ BURST-bare.  Approaches that do not work, and why: a
global FD cap shifts the SM oracle by $-1.2\times10^{-6}$; one-edge gating
in \code{full_coupled} leaves the \code{decoherence_rate} edge uncapped; and
a projection onto positive-semidefinite $\rho$ or a number renormalization,
interp-extrapolation caps, coherence tail-insulation, and a $\Gamma$ gate are
diagnostic-only or ineffective; nor is the tail pump an operator-split aliasing
artefact (the behaviour is node-converged to $n_{\rm steps}=192000$).
\end{diffbox}

\subsection{Headline results}
\label{sec:qkeresults}

\begin{center}\small
\begin{tabular}{@{}R{0.22\textwidth}R{0.22\textwidth}R{0.22\textwidth}@{}}
\toprule
\textbf{medium} & \textbf{$\Neff$ (mixing on)} & \textbf{mixing $\dNeff$} \\
\midrule
vacuum & 3.044071 & $+7.64\times10^{-4}$ \\
matter (MSW) & 3.043976 & $+6.69\times10^{-4}$ \\
\bottomrule
\end{tabular}
\end{center}
The absolute $\Neff$ are at the NLO QED default and $n_{\rm steps}=6000$, with the default
damping collision on the coherences; $\dNeff$ is taken against the mixing-off solve at the same
$n_{\rm steps}$ ($\Neff=3.043307$, the classical reduction).  The density-matrix collision
(\code{collision="matrix_full"}) lowers the vacuum value by $9.4\times10^{-5}$
(Section~\ref{sec:qkecollision}).  For comparison, Bennett et al.\ (2021,
\href{https://arxiv.org/abs/2012.02726}{arXiv:2012.02726}, Table~3) obtain $\Neff=3.0440$ in the
Standard Model with flavour oscillations and NLO QED; there oscillations raise $\Neff$ by
$+5.7\times10^{-4}$ with the full collision integrals, and damping terms on the coherences add a
further $1\times10^{-4}$.  The mixing $\dNeff$ (the deliverable) is QED-baseline-independent: QED
acts identically on the mixing-on and mixing-off runs and cancels in the difference.  MSW lowers
$\dNeff$ by $\sim10^{-4}$ (no resonance).  $\dNeff$ changes by $\le3\times10^{-6}$ between
\code{n_steps} 2000 and 6000 (the order-2 absolute-baseline error cancels in the difference
because the off-grid tail closure, Section~\ref{sec:qkecap}, is applied identically in both
branches), so \code{n_steps=6000} is the default for absolute observables and 2000 a validated
$3\times$ faster delta mode.

\paragraph{The mixing $\dNeff$ is adiabatic}
During decoupling the collision rate is far below the oscillation frequencies
($\Gamma/\omega_{31}\sim10^{-3}$ at $T\sim3$~MeV, $\omega_{31}=\Delta m^2_{31}/2p$), so $\rho$ stays
diagonal in the instantaneous eigenbasis of $H$ up to $\mathcal O(\Gamma/\omega)$.  With large
vacuum mixing that state is not flavour-diagonal: its flavour off-diagonals are the image of
unequal eigenbasis populations.  Projecting the solver's collision $A-\tfrac12\{B,\rho\}$ onto it
gives, per ($\nu/\bar\nu$, bin), a closed equation for the eigenbasis populations $p_i$,
\begin{equation*}
\frac{\dd p_i}{\dd\ln a}=\sum_a|U_{ai}|^2A_a-\Big(\sum_a|U_{ai}|^2B_a\Big)p_i,\qquad
f_a=\sum_i|U_{ai}|^2p_i,
\end{equation*}
with $U$ the eigenvectors of $H$ and $A$, $B$ the classical affine-split rates.  Integrated with
the same exponential-midpoint step, it reproduces the full solve's mixing $\dNeff$ to $0.2\%$
(vacuum) and $0.13\%$ (matter) at $n_{\rm steps}=6000$, and its $\Neff$ to
$1.5\times10^{-6}$.  The mixing $\dNeff$ is therefore carried by the adiabatically transported
eigenbasis populations, and the non-adiabatic remainder is $\lesssim2\times10^{-6}$ in $\Neff$.
The pairwise flavour-conversion rate $R_{ab}=2|H_{ab}|^2\Gamma_{ab}/(\Delta_{ab}^2+\Gamma_{ab}^2)$
(the two-flavour quasi-steady result, $\Delta_{ab}=H_{aa}-H_{bb}$) is not this limit: it assumes the
flavour-basis coherence is small (small mixing, or $\Gamma\gg\omega$), and as a closure on the
flavour diagonal it gives a mixing $\dNeff$ of $+3.0\times10^{-4}$ (vacuum), $2.5\times$ too small.

\paragraph{The off-diagonal coherences are small}
The largest flavour-basis coherence (the maximum of $|\rho_{ab}|$ over bins) rises to
$1.8\times10^{-5}$ and stays there to the end of the run (vacuum, $n_{\rm steps}=6000$).  This is
the expected size: the off-diagonal is the eigenbasis image of the flavour differences $\delta f$
(the $f_0\,\mathbb 1$ part cancels under the eigenbasis projectors), and the largest $|f_a-f_b|$ is
$3.5\times10^{-5}$.  The mass-basis coherence $|\rho^{\rm mass}_{ij}|$ (maximum over bins) peaks at
$1.6\times10^{-7}$ near $\Tcm=6$~MeV and freezes at $2.9\times10^{-10}$ once $\Gamma\to0$.
Corner-cutting computations give coherences of order $0.1$--$1$ instead, an artefact of either an
aliased explicit operator-split dump (whose coherences grow non-convergently and violate
positivity) or of the frozen-diagonal quasi-steady coherence $-L_{oo}^{-1}L_{od}\rho_{\rm diag}$
used as a closure.  The matrix $|U_{ai}|^2$ is nearly singular for the PMNS angles (smallest
singular value $1.1\times10^{-3}$, near $\mu$--$\tau$ symmetry), so once
$\Gamma\lesssim10^{-4}\omega_{31}$ a diagonal that is not of the form $|U|^2p$ maps to $\mathcal
O(1)$ coherences.  The full solve does neither.  $\nu$ and $\bar\nu$ coincide in the symmetric
case ($\bar\rho=\rho^\ast$).

\paragraph{Watch the magnitude, not the lab-frame phase}
The instantaneous lab-frame real and imaginary parts of $\rho_{ab}$ are not
the right object to plot over time.  The off-diagonal carries the fast vacuum
phase $e^{-i\Delta t}$ ($\Delta=\Delta m^2_{ij}/2p$, the vacuum oscillation
frequency, winds $\sim10^{9}$--$10^{10}$ rad
per macro-step at late times), so its real and imaginary components rotate;
the value sampled at frame boundaries is therefore an aliased phase, and it changes
with the step count or the grid, an artifact of displaying a fast
rotation, not an instability in the solve.  The physically meaningful quantities
are phase-/aliasing-immune and do converge: the magnitude $|\rho_{ab}|$, the
mass-basis (eigenbasis-of-$H$) coherence, and $\Neff$ itself.  Doubling the binning
($N=101\to201$ at fixed \code{n_steps}) leaves $\Neff$ unchanged to $9\times10^{-7}$, the mixing
$\dNeff$ to $3\times10^{-7}$, and $|\rho_{ab}|$ identical to four figures.  After decoupling
$\Gamma\to0$, so the dephasing $\Lambda_{\rm off}\to0$ and the late evolution is the undamped free
precession $\rho\to e^{-iH_{\rm vac}t}\,\rho\,e^{+iH_{\rm vac}t}$, with $|\rho_{ab}|$ flat to the
end of the run.  The flavour off-diagonal is then a constant, essentially real offset (the image,
through the real $\delta_{CP}{=}0$ PMNS rotation, of the unequal mass-basis occupations) plus the
frozen $2.9\times10^{-10}$ mass-basis coherence, which precesses.  Neither moves $\Neff$: the
trace is invariant under the free precession, and positivity ($\min f\ge0$) and per-bin number
(drift $\lesssim10^{-13}$) are conserved at every step by the completely-positive
trace-preserving (CPTP) coherent step, so the coherence cannot ``leak into'' and corrupt the
diagonal.  The exact CPTP exponential propagator has no per-component tolerance for the small
imaginary off-diagonal to drift against, which is why the solve is free of the
corotating-scheme instability (Section~\ref{sec:qkesplit}).

\subsection{The public API and status}

\begin{lstlisting}
make_full_coupled_solver(grid=None, *, mixing=True, include_matter=False,
        include_matter_asym=False, use_renorm=True, renorm_pdep=True, qed_nlo=True,
        use_sched="union", tcm_cutoff=0.015, n_steps=6000,
        phie=0.0, eta0=None, q_baryon=0.5, xi=(0.0,0.0,0.0), normal_hierarchy=True,
        gamma_ungated=True, coh_eps_cut=None, n_frames=600,
        tail_closure="srcfix", fd_cap_C=None,
        coherent_method="hybrid", collision_coupling="sourced",
        collision="damping") -> (solve, grid)
# solve(T_init=10.0, lna_max=7.0, delta_neff_init=0.0, me=None, xi=None,
#       matter_asym_scale=1.0) -> dict
\end{lstlisting}
The builder arguments are build-time and the \code{solve} arguments traced
(Section~\ref{sec:knobs}).  The returned
\code{solve} gives \code{neff_final}, the background scalars, \code{f_final}
(3,2,N), \code{rho_final}, the per-step histories \code{T_pl_hist}/\code{T_cm_hist},
the coherence frame arrays (\code{off_fr}, the off-diagonal; \code{ss_fr}, the
frozen-diagonal quasi-steady coherence on the same diagonal, a diagnostic with the conditioning
of Section~\ref{sec:qkeresults}; \code{diag_fr}; at \code{lna_fr},
\code{T_cm_fr}, and the end-of-step \code{T_cm_end_fr}), and the conservation
diagnostics \code{num_drift}/\code{min_f}/\code{any_nan} (plus \code{lna_left}, \code{frame_idx} and the closure provenance \code{fd_cap_C}/\code{fd_cap_logC}).  Key knobs:
\code{mixing=False} gives the exact classical reduction; \code{include_matter}
adds the thermal MSW potential and \code{include_matter_asym} the asymmetric
``solar'' term, scaled per call by \code{matter_asym_scale} (Section~\ref{sec:qke}
Hamiltonian table; \code{eta0}, \code{q_baryon} set its $n_p$);
\code{coherent_method} selects \code{hybrid} (default), \code{expm} or
\code{strang} (Section~\ref{sec:qkesplit}); \code{collision} and
\code{collision_coupling} set the collision on the coherences
(Section~\ref{sec:qkecollision}); \code{tail_closure} (\code{"srcfix"}, default, or
\code{"cap"}) selects the off-grid tail
closure (Section~\ref{sec:qkecap}), and \code{fd_cap_C} is the FD-envelope
multiplier $C$ that takes effect only under \code{tail_closure="cap"} (default
\code{None}: $C=$ \code{interp.FD_ENVELOPE_C} $=10$), threaded to the kernels as $\code{fd_cap_logC}=\log C$.
The remaining arguments mirror the SM transport builders: \code{use_renorm} /
\code{renorm_pdep} (finite-$T$ QED), \code{use_sched} (kernel path, default
\code{"union"} here), \code{tcm_cutoff} (window low edge), \code{phie} / \code{xi}
(electron / neutrino degeneracies), \code{normal_hierarchy} (mass ordering),
\code{n_steps} (operator-split macro-steps; Section~\ref{sec:extend}), plus
three coherence knobs best left at their defaults: \code{n_frames} (how
many coherence snapshots are saved), \code{gamma_ungated} (the off-diagonal damping rate,
Section~\ref{sec:qkecollision}) and \code{coh_eps_cut} (an optional $\eps$ above which $H$ is
zeroed, a diagnostic).
It is reached from BBN by
\code{make_coupled_bbn(transport="qke", qke_kwargs={...})}, which builds the
QKE solve on GPU and runs LINX on CPU (the device split).  Warm $\sim$25 s
($\sim6\times$ the SM \code{transport="full"}), cold $\sim$2 min (A100,
\code{n_steps=6000}).

\paragraph{Demo: the mixing $\dNeff$}
The headline number is the difference between the mixing-on and mixing-off
solves (the off-grid tail closure of Section~\ref{sec:qkecap} cancels in the difference):
\begin{lstlisting}
import neutt                                   # FIRST: enables x64
from neutt.transport.qke.full_coupled import make_full_coupled_solver

solve_on,  _ = make_full_coupled_solver(mixing=True)    # build ONCE (~76 s cold)
solve_off, _ = make_full_coupled_solver(mixing=False)   # the exact classical reduction
neff_on  = float(solve_on()["neff_final"])
neff_off = float(solve_off()["neff_final"])
print(neff_on - neff_off)                        # -> Delta N_eff = +7.6e-4 (vacuum)
\end{lstlisting}
Pass \code{include_matter=True} for the thermal-MSW variant ($+6.7\times10^{-4}$).
From BBN the same solve is driven by
\code{make_coupled_bbn(transport="qke", qke_kwargs={...})}.

\begin{diffbox}
\textbf{Status.}  The mixing $\dNeff$ is the validated headline result.  Through
\code{make_coupled_bbn(transport="qke")} at the NLO default and \code{key_recommended}, the
coupled run gives $\Neff=3.044067$ and $\Yp=0.247061$: the mixing raises $\Yp$ by
$5.4\times10^{-5}$, and the mixing-off run (\code{qke_kwargs=\{"mixing": False\}}) reproduces
the classical \code{transport="full"} $\Yp$ to $3\times10^{-7}$.  This manual gives no
convergence-pinned $\Yp$/$\DoverH$ table for the QKE path.
\end{diffbox}

% ===========================================================================
\section{Extending the code}
\label{sec:extend}
% ===========================================================================

This section is the practical guide to changing what the code computes:
tightening the $\Neff$ precision, taking gradients, scanning the parameters of the
existing model, and adding a collision channel, a background term or an option,
with the invariants each change must preserve and the tests that cover it.  It
closes with the structure that is fixed throughout the code.

\subsection{Increasing the precision of $\Neff$}

The default tolerances and resolutions are tuned, not maximal; for an
arbiter-grade number you can tighten them, at a (usually modest) cost.
\begin{itemize}[nosep]
\item \textbf{Tolerance.}  Lower \code{rtol} from $10^{-6}$ to $10^{-8}$ (it costs
  only $\sim$20\%; the result moves by $\sim$$5\times10^{-7}$).  Do not
  touch \code{atol}: it is already at the pure-relative floor $10^{-30}$;
  raising it silently biases $\Neff$ (Section~\ref{sec:knownissues}).
\item \textbf{Solver.}  For a reference solve, use the plain
  adaptive path \code{hybrid=False} (and \code{use_sched=False}, the bit-faithful
  kernels) with \code{rtol=1e-8}; the implicit \code{Kvaerno5} path
  (\code{implicit=True}) is the L-stable cross-check.  For the hybrid, increase
  \code{n_head} (more fixed exponential-midpoint steps), though its blind spot
  is the order-2 head, so the plain path is the cleaner arbiter.  A fifth-order
  method is required, not optional: the discontinuous precision gate fools the
  lower-order embedded error estimators (a third-order solver mis-estimates $\Neff$
  by $\sim$$9\times10^{-3}$), so the default \code{Tsit5} (order 5) is the floor.
\item \textbf{Grid and quadrature resolution.}  Raise \code{nbins} (must stay
  $4n{+}1$) and/or \code{eps_range} by passing \code{grid=make_eps_grid(...)} to the
  builder.  The Gauss orders are not builder arguments: \code{n_gl}/\code{n_lag} are
  arguments of the kernel functions (and \code{n_gl} of \code{build_nunu_kernel} /
  \code{build_nunubar_kernel}), which \code{assemble._kernel_rate_groups} and the
  builders call at their defaults, and the Gauss--Laguerre node counts of
  \code{thermo} are module constants (100 for the free gas, 25 for the $\phi_e$ and
  QED terms), so raising them is an edit at those call sites.  All of these are
  build-time and trade memory/compile for accuracy.
\item \textbf{QKE node count.}  For the QKE solve, the absolute
  $\Neff$ converges only at order 2, so raise \code{n_steps} (default 6000)
  toward $\sim$18000 for an arbiter absolute number; the difference
  $\dNeff$ already converges by \code{n_steps=2000}.
\item \textbf{Physics completeness.}  Turn on the pieces that are implemented but
  off by default: the density-matrix collision on the coherences
  (\code{collision="matrix"} or \code{"matrix_full"}, $\delta\Neff=-9.2\times10^{-5}$,
  Section~\ref{sec:qkecollision}), the thermal-MSW matter
  potential (\code{include_matter=True}), and the charge-neutrality electron
  degeneracy (\code{phie="auto"}; its $\Neff$ effect is $\ll10^{-7}$).  The only
  un-modelled pieces are the QKE $\nu$--$\nu$ refractive terms
  ($|\delta\Neff|\lesssim3\times10^{-7}$, Section~\ref{sec:qke}), BURST's $\phi_e$
  EoS-derivative closure, and any QED beyond
  the ported NLO $\mathcal O(e^3)$ (the next term of the QED equation of state,
  $\mathcal O(e^4)$, is estimated at $\delta\Neff\simeq3.5\times10^{-6}$ by Bennett et al.,
  \href{https://arxiv.org/abs/1911.04504}{arXiv:1911.04504}, Eq.~4.35, and is not included).
\end{itemize}

\subsection{Gradients and gradient-based inference}

\texttt{neutt} is differentiable and \code{jax.grad} runs through the whole
solve, but one caveat must be respected before trusting a gradient.  The precision
gate (Section~\ref{sec:gatenum}) is a hard threshold, flat on either side, so
its derivative is zero almost everywhere, and the reverse-mode adjoint silently
drops the gate's contribution to $\dd\Neff/\dd\theta$.  The forward $\Neff$
is unaffected (and bit-faithful), but a gradient taken through the full coupled
solve is not validated against finite differences and can be biased.
\textbf{Before any gradient-based inference (HMC / NUTS, a Fisher
forecast), check \code{jax.grad} against a finite-difference $\Neff(\theta)$ for
the parameters you will differentiate.}  The smooth gate \code{gate_width}$>0$
replaces the hard threshold with a sigmoid.  It is an argument of
\code{assemble_rhs} and reaches a solve only through
\code{make_coupled_bbn(gate_width=...)} on the semi path; the decoupling, hybrid and
QKE builders (and the affine \code{assemble_split_rhs}) use the hard gate.  Keep the
forward pass at \code{gate_width=0} for the bit-faithful number.

\subsection{Scanning the parameters of an existing model}

The hot path is built to be swept with no recompile.  Build the solver once and
vary the traced inputs per call:
\begin{itemize}[nosep]
\item \textbf{SM / cosmology:} \code{delta_neff_init}, \code{eta_fac}, \code{tau_n_fac},
  \code{nuclear_rates_q}, the electron mass \code{me}, and the transport start
  \code{T_init} are all traced (Section~\ref{sec:knobs}), and an MCMC varies them
  with zero recompile.  For a batch, \code{vmap} the built \code{solve} over a
  parameter array.
\item \textbf{Lepton asymmetry:} the neutrino degeneracy \code{xi} is traced per
  call: \code{run(xi=...)} (coupled) or \code{solve(xi=...)} (standalone)
  recomputes the initial spectrum and the noise floor inside the compiled
  solve, so a $\xi$-scan reuses one build with no recompile (Section~\ref{sec:knobs}).
\end{itemize}

\paragraph{Parallelism: batch, do not shard}  A single coupled solve is tiny
(the state is a few kilobytes) and launch/bandwidth-bound, so sharding one
solve across several GPUs buys nothing ($\lesssim1.2\times$, often $<1\times$).  The
useful parallelism is across independent solves: \code{vmap} a batch of
cosmologies onto one GPU (the per-cosmology cost drops sharply once the launch
overhead is amortized), or map an MCMC's walkers onto separate GPUs, one whole
solve each.

\subsection{Adding a channel, a background term, or an option}

\begin{itemize}[nosep]
\item \textbf{A classical collision channel.}  A kernel is a function returning
  \code{(net, frs)}, each of shape $(3,2,N)$ per $\Tcm^5$ (the caller multiplies by
  $\Tcm^5$).  It enters through \code{assemble._kernel_rate_groups}, whose per-kernel
  tuple \code{_all_kernel_rates} sums into the gate numerator and denominator (and
  into \code{leg4} if the channel exchanges energy with the plasma), and which
  \code{compute_sumrules} and \code{decoherence_rate_split} unpack.  Everything
  downstream reads those sums: \code{assemble_rhs}, the affine
  \code{assemble_split_rhs}, the noise floor \code{compute_prec_trans}, the QKE
  damping \code{decoherence_rate}, and through \code{make_rhs} /
  \code{make_split_rates} every solver.
\item \textbf{A background energy density.}  $H$ and $\dd\Tpl/\dd\ln a$ are formed
  in \code{decoupling._background} (the inert $\rho_{\rm dr}$ is the template), which
  both \code{make_rhs} and \code{make_split_rates} call, so a term added there
  reaches the plain, hybrid and QKE solves alike.  A plasma-EoS correction enters
  \code{decoupling._plasma_eos} as a $(\delta\rho,\delta P,\delta C_{\rm pl})$
  triple, the pattern of \code{thermo.qed_plasma_deltas}.
\item \textbf{A density-matrix (QKE) channel.}  The matrix kernels in
  \code{transport/qke/} reuse their classical module's tiles, nodes and measure and
  change only the coupling and the blocking: $(1-f)$, $f$ become $(I-\rho)$, $\rho$
  inside anticommutators (e.g.\ \code{collide_matrix.assemble_from_phi}).  A new
  channel follows that pattern and is wired into \code{full_coupled}'s
  \code{collision} option (Section~\ref{sec:qkecollision}).
\item \textbf{An option.}  A switch that changes the shape of the computation is a
  builder argument (build-time); a number to be scanned is a \code{solve} argument
  passed to the RHS as data, never closed over as a Python constant
  (Section~\ref{sec:knobs}).
\end{itemize}

\paragraph{Invariants to preserve}
\begin{itemize}[nosep]
\item \textbf{Detailed balance and the affine form.}  \code{net} must vanish at the
  equilibrium spectrum (the noise floor \code{prec_trans} is that residual,
  recomputed from the same kernels), and every forward rate must carry the Pauli
  factor $(1-f_1)$ and every reverse rate $f_1$, so that the affine split
  (Eq.~\eqref{eq:affine}) is exact.
\item \textbf{No data-dependent control flow.}  Kinematic branches become fixed
  tiles with runtime regime masks (Section~\ref{sec:kerneltensor}), and masked
  square roots and denominators carry double-\code{where} guards so the kernel stays
  \code{grad}-safe (Section~\ref{sec:float64}).
\item \textbf{Faithful defaults.}  A port-only change is default-off, a runtime
  identity, or reversible with one flag (Appendix~\ref{sec:fortranmap}), so the
  Fortran-parity oracle tests (\code{lag_tail_fix=False}, \code{qed_nlo=False},
  \code{SIN2THW_BURST}) and the canonical pins are unchanged.
\item \textbf{The mixing-off reduction.}  A QKE change keeps \code{mixing=False}
  bit-identical to the classical solve, and a matrix kernel reduces to its
  classical kernel on a diagonal $\rho$ and sources no coherence from it.
\end{itemize}

\paragraph{Tests to add}  A new kernel gets an oracle diff (a Fortran golden in
\code{tests/oracle_*/} with its \code{test_<kernel>.py}) and the equilibrium sum
rule (net $\to0$ at $f^{\rm eq}$; Section~\ref{sec:sumrules}); a scheduled form gets
a parity test against its reference (\code{test_sched_parity.py}); a matrix kernel
gets the diagonal-reduction test (the pattern of \code{tests/test_qke_collide*.py}).
An $\Neff$-affecting change is checked against the canonical pins with the
\code{rtol=1e-8} arbiter (Section~\ref{sec:validation}).

\subsection{Fixed structure}

\begin{itemize}[nosep]
\item \textbf{Three flavours $\times$ two cp.}  The flavour axis is three flavours
  by two cp states throughout the kernels, the noise floor and the QKE $3\times3$
  density matrix; a sterile species ($n_\nu>3$) means generalizing that axis in every
  kernel.
\item \textbf{Abundances and recombination live outside.}  The nuclear network is
  LINX's; there is no recombination / sound-horizon stage.  A new observable on
  the BBN side is a LINX change, not a transport one.
\item \textbf{The transport window.}  Transport runs over comoving $T\in[0.015,10]$
  MeV (the defaults of \code{tcm_cutoff} and \code{T_init}); above it the spectrum is
  the equilibrium initial condition, below it the rates are masked and the spectrum is
  frozen, so a model with relevant dynamics outside that window needs the window and
  the equilibrium initial condition changed.
\end{itemize}

% ###########################################################################
\part{The interacting dark sectors}
% ###########################################################################

% ###########################################################################
\section{The dark sectors: overview and scope}
\label{sec:dark}
% ###########################################################################

The research purpose of BURST is dark-sector physics in BBN, and the port now
carries several dark extensions of the neutrino transport.  This part documents
them, with emphasis on why the (often subtle) design choices were made.

\subsection{What is covered, and what is not}

\textbf{Covered here:}
\begin{itemize}[nosep]
\item the \textbf{dark boson decay sectors}: a massive boson $X$ that couples
  to neutrinos through $X\leftrightarrow\nu\bar\nu$: a \textbf{vector $Z'$} and a
  \textbf{majoron / pseudoscalar $J$}
  (Sections~\ref{sec:darkamp}, \ref{sec:darktransport});
\item the \textbf{decay-tail solvers} and the dark hybrid/implicit heads
  (Section~\ref{sec:darksolvers});
\item the \textbf{thermal-majoron $\dNeff=+2.43$ case study}
  (Section~\ref{sec:majoron});
\item the out-of-equilibrium (OOE) scalar decay (folded into
  Sections~\ref{sec:darksolvers}--\ref{sec:majoron});
\item the \textbf{$\chi\leftrightarrow\nu$ contact-interaction sector}: a
  dark-matter fermion scattering on neutrinos, including its strong-coupling
  exact-shape solver, documented up to the in-progress lepton-kernel source
  fix (Section~\ref{sec:chinu}).
\end{itemize}
\textbf{Not covered:} the legacy \code{chi}/\code{dv} kernels these scenarios
superseded, and an end-to-end external dark BBN oracle; the dark sectors
are validated against brute-force collision integrals and conservation laws rather
than a full reference code (a deliberately different validation contract from the
SM sector).

The LINX-side entry points quoted in this part (\code{make_coupled_bbn} with
\code{dark=}, \code{chi_sector=} or \code{solver="exact"}, and the scalar-decay
\code{mcmc/} driver built on it) belong to the dark-capable LINX coupling the
sectors were developed with; the coupling of Section~\ref{sec:linx} does not take
them.  Everything on the \texttt{neutt} side (the sectors, kernels and solvers) is
in this package.

\subsection{The one physical picture behind every dark design choice}
\label{sec:darkpicture}

Almost every subtlety below follows from a single fact about the coordinate
system.  Transport evolves the spectrum in comoving energy $\eps=p/\Tcm$
with $\Tcm=\Tcm^{\rm init}e^{-\ln a}$.  A neutrino that has decoupled
(below $\Tcm\approx1$ MeV the SM collision rates $\propto\Tcm^5$ are negligible)
then sits at \textbf{fixed $\eps$}: its physical momentum redshifts but its
comoving-grid label does not move.

Now let a dark boson $X$ of mass $m_X$ decay $X\to\nu\bar\nu$.  At rest it emits
two neutrinos each at physical energy $\simeq m_X/2$, i.e.\ at comoving
\begin{equation}
  \eps_{\rm feature}=\frac{m_X/2}{\Tcm},
  \label{eq:feature}
\end{equation}
a quasi-monochromatic line.  Because $\Tcm$ is falling, the same physical
decay line climbs to ever-higher comoving $\eps$ as time goes on, and once
deposited it freezes there (the decoupled neutrinos do not move it).  Two
failure modes recur, and the dark machinery exists to defeat them:
\begin{enumerate}[nosep]
\item \textbf{deposition into near-empty high-$\eps$ bins.}  The SM Fermi--Dirac
  occupation at $\eps\gtrsim14$ is $\lesssim10^{-6}$ (a few $\times10^{-6}$ at
  $\eps\simeq13$).  A decay puts a feature
  there; a pure-relative error tolerance sees a large relative
  change in those bins and drives the step size to zero, the ``$\approx$28
  keV step collapse.''
\item \textbf{the climbing edge runs off the grid.}  Once $\eps_{\rm feature}$
  exceeds $\eps_{\max}$ the injection piles up in the boundary bin, and because
  $\rhonu=\int\eps^3 f\,\dd\eps$ the $\eps^3$ weight amplifies that
  mis-resolved pile-up into a spurious $\Neff$.
\end{enumerate}
Both are step-control / discretization pathologies, not physics bugs; the
kernels themselves are oracle-validated.  Keep
Eq.~\eqref{eq:feature} and these two modes in mind; they explain the decay-tail
leg structure (Section~\ref{sec:darksolvers}), the down-weighted error norm, the
energy-fraction monitor, and the thermal-majoron off-grid fix
(Section~\ref{sec:majoron}).

\subsection{A different validation contract}
\label{sec:darkvalidation}

\begin{diffbox}
The SM sector is gated bit-for-bit against the Fortran BURST snapshot.  The dark
sectors cannot be: BURST has no QKE/dark end-to-end oracle one can trust;
the Fortran \code{qburst} end-to-end solve is reported to diverge
(its operator-split corotating scheme appears step-size unstable).  So the dark
work is validated against \textbf{unit-level oracles and exact invariants}
instead: a brute-force direct evaluation of the collision integral, the
conservation sum rules (number, energy, and the cross-species ledger), the
matrix exponential and analytic two-flavour oscillation for the coherent sector,
and, importantly, the \textbf{mixing-off / dark-off bit-reduction} to the
already-validated SM transport.  This is why the dark design leans on
exactness-by-construction (positivity-preserving exponential steps,
trace-preserving channels) rather than on tolerance tuning.
\end{diffbox}

% ===========================================================================
\section{The dark amplitudes and data types}
\label{sec:darkamp}
% ===========================================================================

\subsection{The scenarios and their matrix elements}

A massive ($\sim$MeV) dark boson $X$ couples to neutrinos via
$X\leftrightarrow\nu\bar\nu$; it is populated and depleted by (inverse) decay
against the neutrino bath, and its energy density enters the Hubble rate.  The
decay sectors implement two scenarios; couplings are flavour-diagonal 3-vectors
$g_\alpha$ (default universal), neutrinos massless and left-handed.

\begin{itemize}[nosep]
\item \textbf{Vector $Z'$} (\code{spin_dof = 3}):
  $\mathcal L\supset g_\nu Z'_\mu\,\bar\nu_\alpha\gamma^\mu P_L\nu_\alpha$.
  Process $Z'\to\nu\bar\nu$, partial width $\Gamma=g_\nu^2 m_{Z'}/(24\pi)$ per
  flavour.
\item \textbf{Majoron / pseudoscalar $J$} (\code{spin_dof = 1}):
  $\mathcal L\supset-\tfrac{i}{2}g_{\rm CGJPS} J\,\bar\nu_\alpha\gamma^5\nu_\alpha$
  (Majorana $\nu$; the $\tfrac12$ is the Majorana factor), with $g_{\rm CGJPS}$
  the universal coupling of the cosmological majoron of Chang--Giovanetti et al.\
  (\href{https://arxiv.org/abs/2401.00687}{arXiv:2401.00687}).  Process
  $J\to\nu\bar\nu$, partial width $\Gamma=g_{\rm CGJPS}^2 m_J/(16\pi)$ per flavour
  ($\tau_J=16\pi/3 g_{\rm CGJPS}^2 m_J$); cross-checked against Heeck
  (\href{https://arxiv.org/abs/1709.07670}{arXiv:1709.07670}).  The kernel
  coupling is $g=g_{\rm CGJPS}/\sqrt2$ (the Majorana--Dirac pair-counting factor),
  so $g^2=g_{\rm CGJPS}^2/2$ below.
\end{itemize}

Both reduce to a single kernel form,
\begin{equation}
  \overline{\sum}|M|^2 = 2\,g^2\,m_X^2 ,
\end{equation}
($\sum$ over all final spins, no averaging, the degeneracy division lives in
the kernel prefactor), with $g^2=g_\nu^2$ (vector) and $g^2=g_{\rm CGJPS}^2/2$ (majoron)
(the Majorana--Dirac pair-counting factor).  All eight tree-level $|M|^2$ for
the three scenarios (the two decays plus the $\mathcal O(g^4)$ $2\!\to\!2$
partners) were derived with FeynCalc and gated by 16 internal $+$ 6 literature
checks plus an independent adversarial re-derivation; the $\mathcal O(g^4)$
amplitudes are kept for completeness but are negligible at BBN-viable couplings.

\begin{diffbox}
The dark amplitudes are new physics, not a port of a Fortran routine;
they were derived from scratch (the headless-FeynCalc workflow in
\code{dark_amplitudes/}, archived on the BURST repository's \code{dark} branch).  The $\chi$ amplitudes were validated the same way
(and shown to be the SM lepton-kernel family at $C_V=1,C_A=0,m_e\to m_\chi$,
$G_{\rm eff}^2=g_\chi^2 g_\nu^2/2m_{Z'}^4$); the $\chi$ transport is documented in
Section~\ref{sec:chinu}.
\end{diffbox}

\subsection{The dark data types (\texttt{dark.py})}

\code{BosonDecaySector} is an \code{equinox.Module} whose coupling/mass fields
are dynamic (so a mass/coupling scan compiles once and re-runs); only
\code{spin_dof} is static.

\begin{center}\small
\begin{tabular}{@{}R{0.16\textwidth}R{0.74\textwidth}@{}}
\toprule
\textbf{field} & \textbf{meaning} \\
\midrule
\code{m_x} & boson mass [MeV] \\
\code{gsq_diag} & per-flavour $g^2$, shape (3,); $\overline{\sum}|M|^2=2\,g^2 m_X^2$ \\
\code{tratio} & $\Tcm/T_X$, the (constant) grid-unit ratio for the $X$ grid ($\eps=p/T_X$) \\
\code{grid} & the $X$ $\eps$-grid (default \code{make_dv_grid()}: 101 bins, range 25, offset 0.0025) \\
\code{spin_dof} & static; $X$ degrees of freedom: \textbf{3} (vector) or \textbf{1} (scalar) \\
\bottomrule
\end{tabular}
\end{center}

Constructors \code{make_vector_decay_sector(m_zp, g_nu, ...)} (spin\_dof 3,
\code{gsq_diag} $=g_\nu^2$) and \code{make_majoron_sector(m_j, g, ...)}
(spin\_dof 1, \code{gsq_diag} $=g_{\rm CGJPS}^2/2$ for the input coupling $g=g_{\rm CGJPS}$).
Supporting functions:
\begin{lstlisting}
boson_moments(grid, f_x, m_over_T, T_x, spin_dof) -> (n, rho, P)   # massive-boson moments
dark_be_equil(eps, m_over_T, xi=0.0)   # Bose-Einstein, the default initial X spectrum
dark_fd_equil(eps, m_over_T, xi=0.0)   # Fermi-Dirac (used by the chi sector)
\end{lstlisting}
With $e=\sqrt{\eps^2+(m/T)^2}$ and $\mathrm{pref}=\code{spin_dof}/2\pi^2$,
$n=\mathrm{pref}\,T_x^3\!\int\eps^2 f$, $\rho=\mathrm{pref}\,T_x^4\!\int\eps^2 e f$,
$P=\tfrac{\mathrm{pref}}{3}T_x^4\!\int\eps^4 f/e$.  \code{boson_moments} is
statistics-agnostic (it just integrates $f$), so it provides $\rho_X$ for the
Hubble rate in every dark RHS.  The Bose--Einstein occupation diverges as
$(\eps,m)\to0$, which is why the $X$ grid carries a small offset.

% ===========================================================================
\section{Dark boson decay transport}
\label{sec:darktransport}
% ===========================================================================

The decay sector adds an \code{'f_x'} component to the state and evolves the two
coupled processes $X\to\nu\bar\nu$ (decay; depletes $X$, injects $\nu$) and
$\nu\bar\nu\to X$ (inverse decay; depletes $\nu$, repopulates $X$).  The module
is \code{neutt/transport/dv_nu.py}.

\subsection{The 1$\leftrightarrow$2 (inverse-)decay kernels}

A $1\!\leftrightarrow\!2$ kernel differs from the SM $2\!\to\!2$ kernels in
three physically important ways: the $X$ final/initial state carries a Bose
enhancement/absorption factor $(1+f_X)$/$f_X$, the neutrino legs carry Pauli
blocking $(1-f_\nu)$, and the remaining collision integral is one-dimensional
(the two-body kinematics fix the rest), over the parent-$X$ energy
(Gauss--Laguerre) for the $\nu$-tracked kernel, or over the daughter-$\nu$
momentum (Gauss--Legendre) for the $X$-tracked kernel.  The kernels return $(\mathrm{net},
\mathrm{frs})$ directly in s$^{-1}$ (a $1\!\leftrightarrow\!2$ rate scales
as $g^2 m^2/E$, not $G_{\rm F}^2 T^5$, so, unlike the SM kernels, there is no
$\Tcm^5$ factor to apply later).

\begin{lstlisting}
# d f_X / dt  (X is the tracked species; daughter-nu momenta span a FINITE range -> Gauss-Legendre)
x1_nu_decay_rates(nu_grid, x_grid, f_nu, f_x, meps, tratio, tcmev, gsq_diag,
                  spin_dof=3, nleg=20, off_grid_mode="default") -> (net, frs)   # (3, Nx)

# d f_nu / dt (nu tracked; parent-X energy runs a SEMI-INFINITE tail -> Gauss-Laguerre)
nu1_x_decay_rates(nu_grid, x_grid, f_nu, f_x, meps, tratio, tcmev, gsq_diag,
                  nlag=20, off_grid_mode="default") -> (net, frs)               # (3, 2, Nnu)

# affine forms  net = A - B*f  (for the exponential-midpoint stiff head)
x1_nu_decay_affine(nu_grid, x_grid, f_nu, meps, tratio, tcmev, gsq_diag,
                   spin_dof=3, nleg=20, range_limit=False,
                   off_grid_mode="default", forward_only_tcm=0.0) -> (A_x, B_x)
nu1_x_decay_affine(nu_grid, x_grid, f_nu, f_x, meps, tratio, tcmev, gsq_diag,
                   nlag=20, range_limit=False,
                   off_grid_mode="default", forward_only_tcm=0.0) -> (A_nu, B_nu)
\end{lstlisting}
Here $\mathrm{meps}=m_X/\Tcm$, $\mathrm{tratio}=\Tcm/T_X$, $\mathrm{tcmev}=\Tcm$
[MeV].  Both kernels are exactly linear in the local tracked occupation
(it factors out of the inner quadrature), which is what makes the affine
$A-Bf$ split exact and lets the exponential-midpoint head integrate the stiff
decay loss $-B_X f_X$ analytically (Section~\ref{sec:darksolvers}).  The $X$
(boson) spectrum is interpolated without the fermionic high-occupation
clamp (BE occupations exceed 1 at small momentum); the neutrino legs use the
standard fermionic log-interp.  Convenience aliases
\code{nu1_dv_decay_rates}/\code{nu1_j_decay_rates} (the $\nu$-side is identical
for a vector or scalar parent) and \code{dv1_nu_decay_rates}/\code{j1_nu_decay_rates}
fix \code{spin_dof}.

\subsection{Why the dark legs are ungated: the two gate traps}
\label{sec:darkgate}

This is the central correctness choice and deserves the space.  Recall
(Section~\ref{sec:gatenum}) that the SM kernels carry a per-step
quadrature-cancellation noise: $\mathrm{net}$ is a small difference of
large forward/reverse integrals on a coarse grid, so near equilibrium they need
the precision gate to suppress sign-flickering noise that would otherwise corrupt
the adaptive error estimate.  The $1\!\leftrightarrow\!2$ decay kernels are
fundamentally different and \textbf{must not be gated}:

\begin{quote}
The decay integrands are smooth log-interpolated products over a clean
Gauss--Legendre/Laguerre quadrature; there is no catastrophic
cancellation.  Their detailed-balance residual is a smooth $\sim10^{-16}\cdot
\mathrm{frs}$ drift (negligible, $\lesssim10^{-13}$ s$^{-1}$ even at strong
coupling).  Integrating that residual is harmless; gating it is what hurts.
\end{quote}

Two concrete traps were diagnosed (and both are why the legs are added raw,
after the SM gate):
\begin{itemize}[nosep]
\item \textbf{Trap 1: folding the dark legs into the SM gate sums.}  The gate
  keeps bin $i$ open when $|\sum\mathrm{net}_i|/\sum\mathrm{frs}_i/\mathrm{prec}_i
  \ge\mathrm{nuratiotol}$, with $\mathrm{prec}_i$ the tiny per-bin
  equilibrium-noise floor (Section~\ref{sec:gatenum}).  The dark net is
  real signal far above that floor, so folding it into the numerator flips
  $\sim$600 near-equilibrium SM bins open (the dark $\mathrm{frs}$ would only
  enter the denominator, if anything lowering $\mathrm{test\_ratio}$); the
  per-step flicker of those newly-open gates breaks the explicit error estimate
  $\to$ systematic tail drift $\to$ step collapse.
\item \textbf{Trap 2: giving the dark legs their own hard relative-floor
  gate.}  A hard $|\mathrm{net}|/\mathrm{frs}\ge$ floor reintroduces threshold
  flicker (bins crossing the floor between steps) and a per-executable
  step-count ``roulette'' (the same nondeterministic-stall family as the SM
  \code{atol} stall).
\end{itemize}
The resolution: discard the dark $\mathrm{frs}$, add the raw $\mathrm{net}$ legs
to the SM RHS after \code{assemble_rhs} and after the SM gate, and add
$\rho_X$ to $H$.  Because the residual is smooth and tiny, the ungated
integration is accuracy-neutral (validated to $\sim10^{-6}$ in $\Neff$ against
the no-cap reference).

\begin{diffbox}
The contrast with the $\chi$ contact kernels (Section~\ref{sec:chinu}) is instructive:
those are the SM weak-kernel family (they cancel and carry
$\sim10^{-5}\cdot\mathrm{frs}$ noise), so they do keep the $\Tcm^5$
convention and do require their own gate, separate from the SM legs,
to avoid Trap~1.  ``Conservation is not a cut oracle'': a kernel is gated only
if it has cancellation noise, not merely because it is a collision term.
\end{diffbox}

\subsection{The three interpolation knobs and why they exist}
\label{sec:darkknobs}

The decay kernels expose three off-grid / feature-handling knobs.  At the kernel
level all three default to a no-op, so the explicit decay path is bit-identical to
the pinned reference; they exist to handle regimes the default interpolation
mishandles.  The production stiff solver selects the physically-correct settings
among them (last paragraph).

\paragraph{\texttt{off\_grid\_mode}: the fx3 off-grid read}
In \code{nu1_x_decay_*}, the Gauss--Laguerre tail over the parent-$X$ energy
$E_X\in[a,\infty)$ (with $E_X$, $q_3$ in $T_X$ units) samples $f_X$ at
$q_3=\sqrt{E_X^2-(m_X/T_X)^2}$, which reaches values far above the $X$ grid
edge ($q_3$ up to $\sim$90 vs
$\eps_X^{\max}=25$); 40--90\% of the nodes are off-grid even for the
genuine daughter, because the integral runs over the whole $E_X$ tail.  The
default order-5 log-space polynomial extrapolates there, and for a
sharply-falling (quadratic-in-$q$) log-occupation it \textbf{overshoots upward}
(large Lagrange weights amplify the curvature), reading $f_X$ up where it should
be $\approx0$.  Since the spurious source lands at high $\eps$ and $\rhonu$
weights by $\eps^3$, this manufactures energy, the thermal-majoron ``break''
(Section~\ref{sec:majoron}).  The fix is the physical read for a
non-relativistic parent (whose occupation is genuinely $e^{-q^2/2mT}$-
suppressed off-grid): \code{off_grid_mode="floor"} returns $f\approx0$ above the
edge.  The three modes:
\begin{itemize}[nosep]
\item \code{"default"}: order-5 window to $2\eps_X^{\max}$, then log-linear,
  then floor.  Bit-faithful to BURST; benign on a smooth explicit $f_X$.
\item \code{"linear"}: log-linear downward past the edge.  Not used for
  the non-rel parent: it decays too slowly (under-resolves the quadratic
  falloff) and still leaves a spurious tail.
\item \code{"floor"}: $f\approx0$ above the edge.  The correct read for a
  non-relativistic parent; the dominant cure for the majoron break.
\end{itemize}

\paragraph{\texttt{forward\_only\_tcm}: dropping the inverse decay at low $T$}
Below this $\Tcm$ [MeV], the inverse decay $\nu\bar\nu\to X$ is dropped (forward
$X\to\nu\bar\nu$ only).  Two reasons: physical, the inverse rate is
sourced by the frozen-FD $\nu\bar\nu$ phase space and is Boltzmann-suppressed
$\propto e^{-m_X/T}\to0$ at low $T$ (a $\nu\bar\nu$ pair rarely reaches
$\sqrt s\ge m_X$; the $X$ abundance does not enter this source), so
dropping it is accuracy-neutral where it matters; numerical, the
inverse-decay integrand is precisely the source of the off-grid
ill-conditioning past $\ln a\approx1.7$.  Mechanically the source coefficient
$A_X\to0$ and $B_X\to$ the exact forward Pauli loss $\ge0$, so the $X$ only
depletes (no inverse-decay runaway, no reliance on a monotonicity stabilizer).
Default $0.0$ never fires (bit-identical); the thermal-majoron recipe uses
$\approx0.7$.

\paragraph{\texttt{range\_limit}: the window-max cap}
\code{range_limit=True} caps an off-grid overshoot at the local interpolation
window's maximum occupation, the shape-preserving (monotone-FD) bound.  It is
off by default.  It exists because the decay deposits a genuine super-FD
feature (the $\eps_{\rm feature}$ line), which the SM-path FD-envelope cap
(\code{fd_cap_logC}; Section~\ref{sec:qkecap}) would erase, so the dark
legs cannot use the FD cap, and \code{range_limit} is the milder,
feature-preserving alternative when an off-grid overshoot must be bounded.  (It is
not used on the SM/QKE path, where it is not load-bearing; it is retained as a dark
opt-in.  The two caps are complementary, not interchangeable; see
Section~\ref{sec:qkecap}.)

\paragraph{The production-solver defaults.}
The kernel arguments above default to the bit-faithful no-ops, and the explicit
decoupling path uses them as-is; the production stiff dark-decay solver
(\code{make_dark_decay_implicit_solver}, Section~\ref{sec:darkheads}) instead
selects the physically-correct reads.  It defaults to
\code{dark_offgrid_mode="floor"} (zero the off-grid read above $\eps_{\max}$,
correct because the genuine decay feature peaks at
$\eps_{\rm feature}=(m_X/2)/\Tcm<\eps_{\max}$ and the off-grid $\nu$ energy is
FD-suppressed), keeps \code{dark_range_limit} off (the window-max cap is an
opt-in for a relativistic-parent corner), and does not apply the \code{fx3}
boson-read fence (redundant under \code{floor} and \code{monotone_x}; bit-identical
on every physical state, so the affine-identity kernel gate is unchanged to
round-off).  The SM lepton legs inside the dark run likewise default to the
parameter-free \code{sm_tail_closure="srcfix"} (Section~\ref{sec:qkecap}) rather
than the $C=10$ cap, which could clip the genuine super-FD injection.  On the
canonical thermal majoron ($m_J=5$, $g=3.5\times10^{-10}$, $5\times$
over-abundant) this gives $\dNeff=+2.43$, energy- and number-conserving, with $X$
fully depleted; a robustness sweep confirms \code{floor} stays finite and
energy-conserving across stronger coupling, heavier mass, and the relativistic
vector-$Z'$ parent.

% ===========================================================================
\section{The decay-tail solvers}
\label{sec:darksolvers}
% ===========================================================================

A decaying boson breaks the assumptions of the SM solver in the ways set out in
Section~\ref{sec:darkpicture}.  The solvers in \code{decoupling.py} and
\code{hybrid.py} repair this through a sequence of carefully-motivated devices.

\subsection{Why the single SM solve fails, and the leg structure}

The production SM solve uses $\code{atol}=10^{-30}$ (pure-relative; correct for
the stiff $\bar\nu$ tail, Section~\ref{sec:knobs}).  Once a decay feature is
injected, that same $\code{atol}=10^{-30}$ is exactly what triggers the ``28~keV
step collapse'' (failure mode 1): it chases the huge relative change in the
near-empty high-$\eps$ bins and $\dd t$ underflows.  The key realization that
makes a clean fix possible: the SM-critical tight-tolerance transport all
happens above $\sim$1 MeV, and the decay tail all happens below it, where the
neutrinos are decoupled and frozen but for redshift.  So the integration is
split in $\ln a$ with different error control per leg.

\paragraph{Two legs: \texttt{two\_leg\_decay\_solve}}
\begin{lstlisting}
two_leg_decay_solve(term, solver, state0, lna_handoff, lna_stop, n_save,
                    atol_hi, atol_tail, ctrl_kwargs, max_steps, throw,
                    tail_norm=None, adjoint=None)
\end{lstlisting}
\begin{itemize}[nosep]
\item \textbf{Phase 1} $[0,\code{lna_handoff}]$ keeps $\code{atol}=\code{atol_hi}
  =10^{-30}$, the validated SM decoupling solve.  \code{lna_handoff} comes from
  $\code{decay_handoff_tcm}=0.20$ MeV, chosen below $\nu$ decoupling
  ($\sim$1 MeV, so SM transport is essentially done) yet above the decay
  epoch (so the feature forms in phase 2).
\item \textbf{Phase 2} $[\code{lna_handoff},\code{lna_stop}]$ is the decoupled
  decay tail, with the relaxed error control below.  \code{lna_stop} is set so
  the climbing feature edge stops before leaving the grid, at
  $\eps_{\rm feature}=\code{decay_stop_safety}\cdot\eps_{\max}$ (default $0.9$);
  this fences off failure mode 2.
\end{itemize}
This is \code{make_decoupling_solver(dark=..., decay_tail=True)}; validated for the
$\sim$100 keV-decaying majoron with $m_X\lesssim1.1$ MeV
($\Neff=3.07899$).

\paragraph{Three legs: \texttt{decay\_to\_sm\_solve}}
For a heavy / early-decaying boson ($m_X\gtrsim2$, so $\eps_{\rm feature}
=m_X/2\gtrsim1$ even at $\Tcm=1$), the on-grid stop happens at $\Tcm\approx0.11$
MeV, before $e^\pm$ annihilation finishes ($\Tcm\lesssim0.02$), so the
two-leg $\Neff$ has not reached its $(11/4)^{4/3}$ asymptote.  But you cannot let
one solve keep running (it chases the feature off-grid and grinds).  The fix uses
two facts: such a boson is already depleted by the time its feature
reaches the edge, and a frozen feature does not move.  So:
\begin{lstlisting}
decay_to_sm_solve(term_dark, term_sm, solver, state0,
                  lna_decay_handoff, lna_sm_handoff, lna_max, n_save,
                  atol_hi, atol_tail, ctrl_kwargs, max_steps, throw,
                  tail_norm=None, adjoint=None)
\end{lstlisting}
phase 1 (dark, tight \code{atol}) $\to$ phase 2 (dark, down-weighted norm, edge
held on-grid) $\to$ \textbf{phase 3 SM-only}: $f_X$ is dropped from the
state (no more injection) and $e^\pm$ annihilation finishes with the deposited
feature frozen in $\eps$.  The hard $f_X$ drop is energy-exact only if the
boson is depleted at the handoff; the leftover $\rho_X$ is broadcast into the
output (\code{rho_x_final}) so any residual is reported, not silently discarded.
This is \code{make_decoupling_solver(dark=..., decay_to_sm=True)} (takes precedence over
\code{decay_tail}).  Applicability: this stitching path assumes the dark
legs (phases 1--2) are non-stiff (weak coupling / OOE) and short enough to
compile; a stiff / overabundant boson (including the $m_J=5$ thermal majoron,
which is also ``heavy/early'') instead needs the exponential-midpoint head
\code{make_dark_decay_implicit_solver} of Section~\ref{sec:darkheads}.

\subsection{Down-weight, don't floor: \texttt{make\_decay\_tail\_norm}}
\label{sec:decaytailnorm}

\begin{lstlisting}
make_decay_tail_norm(grid, eps_split=13.0, feature_weight=0.05) -> norm(error_pytree)
\end{lstlisting}
The phase-2 controller must not chase the near-empty feature bins.  Two
ways to achieve that, and the choice between them is the subtle part:

\begin{itemize}[nosep]
\item \textbf{An \code{atol} floor ($10^{-10}$)} works numerically but is
  fusion-sensitive: near the floor, tiny code differences (XLA fusion
  order, rounding) change which near-empty bins trip the controller, so
  the step path differs between executables.  The diagnostic: the same
  configuration with only the static \code{max_steps} constant changed
  (50000 vs 8000) gave $\sim$270 steps (clean) vs $>$8000 (grind), and
  \code{max_steps} cannot change the step path, so XLA must have produced a
  numerically different executable.  This is the ``fusion roulette,'' the same
  family as the SM \code{atol} stall.
\item \textbf{A down-weighted error norm (the fix)} keeps $\code{atol}=10^{-30}$
  on the populated $\eps<\code{eps_split}$ bins (full tight control, no
  negative-$f$ excursions) but weights the $\eps\ge\code{eps_split}$ bins by
  $\code{feature_weight}\ll1$ inside the RMS error norm.  Below $\sim$1 MeV the
  SM neutrinos are frozen, so those high-$\eps$ bins are purely the decay
  feature: down-weighting them in the step-size decision stops the
  over-resolution, while the spectrum still evolves exactly via the RHS.  No
  \code{atol} floor $\Rightarrow$ no near-floor flicker $\Rightarrow$
  fusion-robust (identical step counts across executables).
\end{itemize}
$\code{eps_split}=13$ is where SM occupation is $<10^{-6}$, so nothing physical
is down-weighted.  (Implementation note: the norm must be built on
\code{optx.rms_norm} over a pre-weighted pytree; a hand-rolled per-leaf RMS loop
hung the outer-\code{jit} compile of the coupled builder.)

\subsection{The energy-space safety monitor}

\begin{lstlisting}
high_eps_energy_frac(grid, occ, eps_split, n_top=10) -> (tail_frac, top_frac)
\end{lstlisting}
The step controller measures error in occupation space, but $\Neff$
consumes $\rhonu=\int\eps^3 f$, $\eps^3$-weighted, and the down-weighting is
deliberately blind to exactly the high-$\eps$ bins.  So an independent,
post-solve, energy-space diagnostic is needed to keep the compromise honest:
\code{tail_frac} is the $\rhonu$ fraction in $\eps\ge\code{eps_split}$ (the
down-weighted region; $\approx0.5\%$ for the validated smooth tail; an
$\mathcal O(0.1\text{--}1)$ value flags a point where the down-weighting is
hiding real energy), and \code{top_frac} is the fraction in the last \code{n_top}
bins (the grid-edge pile-up indicator, $\sim5\times10^{-6}$ when healthy).  These
are surfaced as \code{eps_tail_energy_frac}/\code{eps_top_energy_frac}; an MCMC
rejects points with \code{eps_tail_energy_frac}$>0.05$, catching the
$m_X\gtrsim1.2$ grid-edge blow-up ($\Neff\to24.5$) at its physical root rather
than post-hoc on $\Neff$.

\subsection{The \texttt{DirectAdjoint} compile-stall fix}

A pure compile-engineering fix for a hard blocker.  The adjoint
(Section~\ref{sec:jaxprimer}) is the strategy diffrax uses to back-propagate
gradients through an ODE solve; it sets the gradient cost and the compile
cost but does not affect the forward solution.  \textbf{Symptom:} the dark
\code{diffeqsolve} compile stochastically hangs 10--40 min CPU-bound (GPU
idle) instead of $\sim$130 s.  \textbf{Cause:} diffrax's default
\code{RecursiveCheckpointAdjoint} builds a recursively-checkpointed
\code{bounded_while_loop} (the machinery that lets reverse-mode differentiation
run through an adaptive solve with bounded memory) whose XLA compile scales with
\code{max_steps}.
\textbf{Fix:} pass \code{adjoint=diffrax.DirectAdjoint()}, which skips that loop
and compiles reliably ($\sim$125 s, 5/5 clean) with a bit-identical
forward solution (DirectAdjoint changes only how a gradient would be
taken, not the forward integration).  The cost is reverse-mode AD
memory/speed, irrelevant for a forward solve or a gradient-free (Metropolis)
MCMC.  \textbf{Caveat:} DirectAdjoint's compile cost scales with the number of
\code{SaveAt} points, so use a modest \code{n_save} ($\sim$64) on the dark path.
Exposed as \code{decay_adjoint}/\code{solve_adjoint} (default unchanged); the
dark hybrid heads hard-wire it on their SM tail.

\subsection{The stiff bulk decay and the dark heads}
\label{sec:darkheads}

For an overabundant boson the decay is \textbf{stiff}: $\Gamma\approx$
const while $H\propto T^2$, so $\Gamma/H$ grows as $T$ drops
($\sim$70 at $\Tcm=0.7$ MeV, $\sim$860 by $0.2$).  Then $B_X\cdot\code{dtmax}
\gg1$ and:
\begin{itemize}[nosep]
\item \textbf{explicit Tsit5} is stability-limited (not accuracy-limited):
  $\dd t$ collapses to $\sim2\times10^{-4}$ and 8000 steps reach only
  $\ln a\approx1.72$; it grinds during the active bulk decay, at a
  populated bin, well before any off-grid issue;
\item \textbf{implicit Kvaerno5} is unconditionally stable but compile-hangs
  $>$20 min; the Newton solve needs the Jacobian of the large dark RHS graph
  (interp $+$ Gauss--Laguerre kernels), which a lower-order implicit does not
  fix.
\end{itemize}
The principled cure exploits the exact per-bin affine structure of the
decay RHS, $\dd f/\dd\ln a=A-Bf$.  Because $\mathrm{net}_X=A_X-B_X f_X$ is
linear in $f_X$ (the partner-$\nu$ and parent-$X$ integrands are frozen
sources), an exponential-Euler / exp-midpoint step
\[
  f' = f\,e^{-Bh}+\tfrac{A}{B}\big(1-e^{-Bh}\big)
\]
is the exact L-stable solution: a convex combination of $f$ and the
attractor $A/B$, hence positivity-preserving and unconditionally stable, with
no Jacobian (no implicit compile blow-up) and no stability ceiling
(no explicit grind).  This is the same exponential-midpoint device as the SM
hybrid head (Section~\ref{sec:hybridnum}), now carrying the dark legs.

The production head is the last of three designs.  The first two were tried and
removed (last present at BURST commit \code{430dfb3}); the table records why they fail:
\begin{center}\small
\begin{tabular}{@{}R{0.30\textwidth}R{0.64\textwidth}@{}}
\toprule
\textbf{head} & \textbf{status / why} \\
\midrule
plain per-bin expMid head (removed) & Physically wrong on the overabundant majoron.  A fixed-step per-bin expMid with no Pauli $f\le1$ enforcement over-deposits the sharp super-FD feature, so the two daughter occupations $f_2{+}f_3>1$ drive the $X$-loss Pauli factor $B_X=1-f_2-f_3$ negative ($B_X<0$); the $X$ exp-step then extrapolates upward and the majoron never depletes (its occupation runs away). \\
Pauli-clipped sub-cycled head (removed) & Sub-cycled IMEX head with energy-renormalized ``delivery.''  Depletes $X$ but fails the conservation gate: it operator-splits the stiff $X\!\leftrightarrow\!\nu$ exchange (decays $X$ on the cold start-of-step $\nu$), over-depleting $X$ and losing the energy ($\rhonu$ $\sim$18\% low). \\
\code{make_dark_decay_implicit_solver} & \textbf{The production thermal-majoron path.}  Sub-cycles $X$ and $\nu$ together with the actual decay affine $A_\nu-B_\nu f$ (\code{delivery="kernel"}), so the kernel's own detailed balance conserves energy and the real $-B_\nu f$ loss self-limits $\nu$ at $f=A_\nu/B_\nu\le1$ (the Pauli clip never fires).  It solves the same equations as the explicit Tsit5 and converges to the right answer (rel $\rhonu$ $1.8\times10^{-5}$ at $\code{n_head}=600$).  \code{monotone_x=True} clamps $f_X$ to its start value (the overabundant majoron only depletes). \\
\bottomrule
\end{tabular}
\end{center}
The key kwargs of the production head: \code{lna_sw=3.0}, \code{n_head=600},
\code{k_sub=8} (sub-cycle count), \code{delivery="kernel"}, \code{monotone_x=True},
plus the \code{dark_offgrid_mode} / \code{dark_forward_only_tcm} off-grid knobs of
Section~\ref{sec:darkknobs}.  (The production head defaults
\code{dark_range_limit=False}, matching the kernel-level \code{range_limit=False}
of Section~\ref{sec:darkknobs}: the \code{"floor"} off-grid recipe zeroes the
off-grid reads before any window-max cap, so the cap is unnecessary;
\code{dark_range_limit=True} is retained only as an opt-in.)

\begin{diffbox}
The dark linear-decay case is easier than the $\chi$ (Section~\ref{sec:chinu})
$2\!\to\!2$ annihilation: a $1\!\leftrightarrow\!2$ rate is linear in the tracked
occupation, so \code{exp_step} is exact and no Newton / annihilation-noise
machinery is needed.  \code{make_dark_decay_implicit_solver} reuses the
sub-cycled-implicit pattern of the $\chi$ solver but not its complexity.
\end{diffbox}

% ===========================================================================
\section{Case study: the thermal majoron $\dNeff=+2.43$}
\label{sec:majoron}
% ===========================================================================

The majoron sector spans two physically distinct regimes, handled by
different solvers, and they are the same \code{make_majoron_sector},
differing only in mass and initial abundance:
\begin{itemize}[nosep]
\item a \textbf{light, out-of-equilibrium (OOE) scalar} ($m_\phi\lesssim1$ MeV,
  decaying near $\sim$100 keV), solved by the explicit two-leg decay-tail driver
  \code{make_decoupling_solver(dark=..., decay_tail=True)} (Section~\ref{sec:darksolvers}),
  giving $\Neff=3.079$;
\item the \textbf{thermal / overabundant majoron} treated in this case study
  ($m_J=5$ MeV, $5\times$ equilibrium), solved by the implicit exponential-affine
  head, giving $\dNeff=+2.43$.
\end{itemize}
So when the manual says ``the OOE scalar'' it means the light regime of this same
majoron sector, not a separate model.  The thermal majoron is the headline dark result
and the best worked example of how the design choices above combine.  Model:
\code{make_majoron_sector} with $m_J=5$ MeV, $g=3.5\times10^{-10}$, initial
abundance $5\times$ Bose--Einstein equilibrium ($\rho_X/\rho_{1\nu}\approx9$ at
the 4 MeV start, with $\rho_{1\nu}$ the energy density of a single neutrino
species, the per-species normalizer for $\dNeff$).  Run the full annihilation
window so the $(11/4)^{4/3}$ $\Neff$
applies.  The result emerged from a sequence of diagnosed wrong answers.

\subsection{Step 1: the early $\Neff\approx40$ was an $\eps^3$ artifact}

A naive solve gave $\Neff(\rm dark)=39.7$, $\dNeff=+36.6$, and a hard
energy budget proves that is unphysical.  Each majoron deposits at most its rest
mass $m_J=5$ MeV (it is deeply non-relativistic at the bulk-decay epoch,
$m/T\approx12$).  In
comoving variables ($\rho a^4$ constant for radiation),
\begin{equation}
  \dNeff = \frac{m_J\,n_X(a{=}1)\,\langle a\rangle_{\rm decay}}
                {\rho_{1\nu}a^4},
  \qquad \langle a\rangle_{\rm decay}\le a_{\rm handoff}.
\end{equation}
Using the code's own reported $X$ number moment, the maximally
redshift-favourable ceiling is $(\dNeff)_{\max}=17.4$ (everything survives to the
handoff then decays).  The code's $+36.6$ is $2.1\times$ over the hard ceiling
$\Rightarrow$ energy is being created $\Rightarrow$ unphysical.  (The internal
ledger missed it because the old ``$E_{\rm delivered}=E_X^{\rm lost}$'' check was
a tautology ($E_{\rm delivered}$ is defined from the $X$ loss),
so it tracked delivery, never total energy conservation.)

\subsection{Step 2: the real mechanism (the off-grid parent-$X$ read)}

The energy creation is the off-grid overshoot of Section~\ref{sec:darkknobs}: the
$\nu$-injection's parent-$X$ Gauss--Laguerre integral samples $f_X$ at nodes up to
$\sim$76--92 while the $X$ grid maxes at $\eps_X=25$, so 40--90\% of the nodes are
off-grid even for the genuine $\eps\approx4$ daughter, and the default
order-5 extrapolation overshoots upward into an $\eps^3$-amplified high-$\eps$
$\nu$ tail that creates energy and number past $\ln a\approx1.7$.  The overshoot is
pure artifact: the $X$ is non-relativistic ($m/T\approx8$ at the off-grid
passage, $\Tcm\sim0.7$ MeV), so its true occupation
off-grid is $e^{-q^2/2mT}\approx0$, and a $\nu$ daughter at $\eps\gg(m_J/2)/\Tcm$ is
kinematically forbidden.  (The inverse decay alone is not the culprit; the
dominant bug is this parent-$X$ read, present even in the forward injection.)

\subsection{Step 3: the fix and the validated result}

The cure has two parts, and the production solver already defaults to the first:
\begin{enumerate}[nosep]
\item \code{dark_offgrid_mode="floor"} (the solver default; it threads to the
  interp-level \code{off_grid_mode} of Section~\ref{sec:darkknobs}), reading
  $f_X=0$ above $\eps_X^{\max}$, the physically correct read for the
  non-relativistic parent.  This is the
  dominant cure: once the overshoot is gone, ``linear'' and ``floor'' give
  identical conservation (the edge value is already tiny).
\item \code{dark_forward_only_tcm=0.7}, the one non-default opt-in: below
  $\Tcm=0.7$ MeV drop the inverse decay, letting the kernel run clean to full $X$
  depletion (a secondary refinement; the un-gated floor also conserves, but
  forward$+$floor depletes $X$ to $\sim10^{-52}$ rather than $\sim10^{-6}$).
\end{enumerate}
The production recipe is
\code{make_dark_decay_implicit_solver(dark_offgrid_mode="floor",
dark_forward_only_tcm=0.7)}.  The result, certified by the true
cross-species sum-rule ledger ($E_\nu^{\rm decay}$ vs $E_X^{\rm lost}$,
$N_\nu^{\rm decay}$ vs $2N_X^{\rm lost}$, used as gates, not
enforcement):
\[
  \boxed{\ \dNeff=+2.43\ }
\]
with energy residual $4.8\times10^{-3}$, number residual $1.7\times10^{-3}$, and
$X$ fully depleted ($n_X/n_X^0\le10^{-13}$ conservation gate; $\sim10^{-52}$
actually achieved).
Supporting validations: the matched $\dNeff$ is flat across the handoff
$\code{lna_sw}\in\{2.0,2.5,3.0\}$ (the $+0.017$ absolute drift is a head SM bias
that cancels against the $g{=}0$ pin); it is threshold-stable across
$\code{dark_forward_only_tcm}\in\{0.6,0.7,0.8\}$ ($+2.438/+2.433/+2.422$).

\paragraph{Why $+2.43$, and what ``residual'' means}
The total $\dNeff=+2.43$ splits into a bulk-decay part (the $X$ that
decays before the hand-off, a firm $\approx+2.34$), and a residual part
(the leftover $X$ that decays after the hand-off).  The bulk $+2.34$ is a
lower bound, and the residual is easy to over-estimate by guessing (a naive
$+0.3$--$0.45$ would inflate the total to $\sim2.6$--$2.7$); capturing it properly
makes its contribution only $+0.09$, so the total is $2.34+0.09=2.43$.  (The two
numbers are not in tension: $+2.43$ is the answer, $+0.09$ is just the
residual-$X$ piece of it.)  The residual is small because most of the
matter-boosted leftover-$X$ energy leaks to the still-coupled plasma
during its late decay at $T\sim0.5$ MeV (the end ratio $\Tpl/\Tcm=1.43$ confirms
the plasma is heated), not to neutrinos, which is also why the
thermal-majoron $\dNeff$ sits well below the naive ``all rest-mass
$\to\nu$'' budget.

\paragraph{Known safe-threshold floor: \texttt{dark\_forward\_only\_tcm}$\gtrsim0.6$.}
At $\le0.5$ the head still conserves perfectly, but the downstream SM Tsit5
tail NaNs, the lepton-kernel boundary fragility
(Section~\ref{sec:knownissues}): a later gate lets the full kernel seed a slightly
larger super-thermal high-$\eps$
$\nu$ tail, and the thermal-majoron tail's \code{make_rhs} runs uncapped
(the FD cap would clip the genuine super-FD feature), so the SM \code{nuer2}/
$\nu\nu$ kernels amplify the seed.  This is not a regression and not FD-cap-fixable
here; $0.7$ sits comfortably inside the safe band.

\paragraph{Demo: reproducing the $+2.43$}
Build the sector and the production head once, then solve (on a GPU node, with
the \code{neutt} repository root on \code{PYTHONPATH}):
\begin{lstlisting}
import neutt                                 # FIRST: enables x64
from neutt import dark
from neutt.hybrid import make_dark_decay_implicit_solver

sec  = dark.make_majoron_sector(m_j=5.0, g=3.5e-10)            # scalar majoron, spin_dof 1
f_x0 = 5.0 * dark.dark_be_equil(sec.grid.abscissas, 5.0/4.0)   # 5x equil.; 5.0/4.0 = m_J/T_init (the m/T arg)
solve, grid = make_dark_decay_implicit_solver(                 # build ONCE, reuse the closure
        sec, lna_sw=2.5, n_head=600, k_sub=8,                  # plateau is flat over lna_sw 2.0-3.0
        dark_offgrid_mode="floor", dark_forward_only_tcm=0.7)  # the off-grid fix
out = solve(T_init=4.0, f_x0=f_x0)
print(float(out["neff_final"]))                 # -> dark N_eff (Delta N_eff = +2.43 vs the g=0 pin)
\end{lstlisting}
The same builder with \code{g=0.0} (or \code{f_x0=0}) is the $\dNeff=0$ pin that
the matched $+2.43$ is measured against; \code{out} also carries the cross-species
ledger and the \code{eps_tail_energy_frac} monitor of
Section~\ref{sec:darksolvers}.  A light, $\sim$100 keV-decaying scalar
($m_X\lesssim1.1$, $\Neff=3.079$) instead uses the explicit decay-tail driver
\code{make_decoupling_solver(dark=sec, decay_tail=True)} of
Section~\ref{sec:darksolvers}.

% ===========================================================================
\section{The $\chi\leftrightarrow\nu$ contact-interaction sector}
\label{sec:chinu}
% ===========================================================================

This sector adds a dark-matter fermion $\chi$ that interacts with neutrinos
through a heavy mediator.  Everything below is validated and usable up to
one unfinished piece: a principled replacement for the off-grid interpolation
closures in the strong-coupling solve, which the last subsection
(Section~\ref{sec:chinufix}) sets out as the open frontier.

\subsection{Physics and the data type}
\label{sec:chinuphys}

This sector adds a \textbf{Dirac dark-matter fermion} $\chi$ coupled to neutrinos
through a heavy $Z'$ mediator.  Because $m_{Z'}$ ($\approx91$ GeV by default, the
mass of the Standard-Model $Z$) is far above the BBN energy scale, the mediator
is integrated out and the interaction is a $2\!\to\!2$ contact (four-fermion)
process: $\chi$ annihilation $\chi\bar\chi\leftrightarrow\nu\bar\nu$ and elastic
scattering $\nu\chi\to\nu\chi$.

The tree-level squared matrix elements were derived and FeynCalc-validated, and
they turn out to be exactly the SM lepton-kernel family (the same
$\nu$--$e$ collision structure used for weak decoupling), with vector/axial
couplings $C_V=1$, $C_A=0$ and the electron mass replaced by the $\chi$ mass,
\begin{equation}
  \me\;\longrightarrow\;m_\chi,
  \qquad
  G_{\rm eff}^2 \;=\; \frac{g_\chi^2\,g_\nu^2}{2\,m_{Z'}^4},
\end{equation}
where $g_\chi$, $g_\nu$ are the mediator's couplings to $\chi$ and to neutrinos.
This is why the kernels were easy to derive (they reuse the
\code{nuer}/\code{nuepma} machinery and keep the SM $\Tcm^5$ rate convention) but,
as the rest of the section shows, hard to solve.

\paragraph{The \texttt{ChiSector} data type.}
\code{dark.ChiSector} is an \code{equinox.Module} mirroring
\code{BosonDecaySector} (Section~\ref{sec:darkamp}); its coupling and mass fields
are dynamic leaves (so an $(m_\chi,K)$ scan compiles once), and only
\code{spin_dof} is static.  Build one with \code{make_chi_sector}:
\begin{lstlisting}
sec = dark.make_chi_sector(m_chi=1.0,        # chi mass [MeV]
                           m_zp=91187.6,     # heavy Z' mediator [MeV] ~ the SM Z
                           g_chi=1.0, g_nu=1.0)   # mediator couplings
# stores the per-flavour contact strength  Kflav = g_chi^2 * g_nu^2 / m_zp^4
\end{lstlisting}
\begin{center}\small
\begin{tabular}{@{}R{0.16\textwidth}R{0.74\textwidth}@{}}
\toprule
\textbf{field} & \textbf{meaning} \\
\midrule
\code{m_chi}   & $\chi$ mass [MeV] \\
\code{Kflav}   & per-flavour contact coupling $K=g_\chi^2 g_\nu^2/m_{Z'}^4$ [MeV$^{-4}$], shape (3,); the rate scales as $K\,\Tcm^5$.  It is the stored contact strength, related to the $G_{\rm eff}$ above by $K=2\,G_{\rm eff}^2$ \\
\code{tratio}  & $\Tcm/T_\chi$, the (constant) grid-unit ratio for the $\chi$ grid ($\eps=p/T_\chi$) \\
\code{grid}    & the $\chi$ $\eps$-grid (default \code{make_chi_grid()}) \\
\code{spin_dof} & static; $\chi$ dof $=\mathbf{4}$ (Dirac: particle$+$antiparticle $\times$ 2 spins, one CP-symmetric $f_\chi$) \\
\bottomrule
\end{tabular}
\end{center}

\subsection{The kernels: amplitudes, energy cuts, and why we brute-force}
\label{sec:chinukernels}

\paragraph{Amplitudes.}  The tree-level squared matrix elements for $\chi\bar\chi
\leftrightarrow\nu\bar\nu$ and $\nu\chi\to\nu\chi$ were derived with FeynCalc and
cross-checked in Mathematica.  As noted above they are exactly the
Standard-Model $\nu$--lepton kernel family at $C_V=1$, $C_A=0$, $\me\to m_\chi$, so
the coupling coefficients are the only thing that changes; everything downstream
reuses the validated SM \code{nuer}/\code{nuepma} machinery and the $\Tcm^5$ rate
convention.

\paragraph{The energy cuts, and why they exist.}  A $2\!\to\!2$ collision rate is a
multi-dimensional phase-space integral.  The classical way to evaluate it (the
way the Fortran reference does, and the way the high-accuracy reference kernel
\code{chi_nu} does) is to perform the angular integrals analytically,
which collapses the problem to a cheap one-dimensional energy quadrature.  The price
is that the integration limits change discontinuously as the bin energy
crosses the kinematic thresholds, so the integrand becomes piecewise: a case
structure of $\sim$40 ``cut'' segments, with boundary energies that have to be
worked out by hand (the $E_1$ vs.\ $m_\chi\sqrt5/2$ case split, and the
cut/transition/limit energies between regimes).  Getting this cut structure right
was a large part of the work.  On a CPU it is essential: a 1-D quadrature per
bin is cheap, whereas evaluating the full multi-dimensional integral directly, at
every bin and every time step, would be prohibitive (the Standard-Model transport
already costs $600$--$1000$ CPU-hours).  The neutrino-side $\chi$ kernels reuse the
SM lepton-kernel cuts verbatim (the kinematics are identical, with
$\code{meps}=m_\chi/\Tcm$); the $\chi$-side, where the massive particle is
the tracked one, needed its own new cut structure.

\paragraph{Why the coupled solve brute-forces instead.}  Those $\sim$40 unrolled
cut-segments make the combined right-hand-side computation graph very large,
and \code{diffrax}'s adaptive solver must compile that graph, which blew the
compile up to tens of minutes.  The production coupled solver therefore takes the
opposite approach and \textbf{brute-forces} the collision integral: it evaluates the
full four-dimensional phase-space integral directly, as a single Gauss--Legendre
quadrature over the partner momentum, two angles, and a final-state energy, with the
kinematic support imposed by masking quadrature nodes that fall outside the
allowed region rather than by analytic integration limits (\code{chi_brute}).  This
sidesteps the cut bookkeeping entirely.  The computation graph is now tiny (one
large reduction, no segments), so it compiles in the normal range; the cost moves to
more quadrature points, which the GPU (built for exactly this kind of dense,
regular arithmetic) absorbs.  The brute-force kernel is validated bin-by-bin
against the analytic cut kernel, so accuracy is not traded for the convenience.

\begin{diffbox}
A residual subtlety: a hard $0/1$ support mask makes the integral kink as
the temperature sweeps through $m_\chi$ (the kinematic boundary crosses fixed
quadrature nodes), which can collapse the adaptive step.  The cure is to rescale the
final-state-energy nodes to the exact physical centre-of-mass window each step
(\code{exact_e3}), so no node is ever out of support, smooth, no kink, and no
spurious flux from tapered nodes.
\end{diffbox}

\subsection{Weak coupling: the straightforward case}
\label{sec:chinuweak}

For realistic parameters (a light $\chi$, $m_\chi\lesssim1$ MeV, and a
small $G_{\rm eff}$ so the $\chi\leftrightarrow\nu$ rate $\Gamma\lesssim H$) the
sector is well-behaved: $\chi$ stays near equilibrium and the spectral distortion
it sources is gentle.  Pass the sector straight into the coupled BBN driver, which
co-evolves $f_\chi$ with the $\nu$ spectrum and $\Tpl$:
\begin{lstlisting}
run, info = make_coupled_bbn(transport="full",
                             chi_sector=dark.make_chi_sector(m_chi=1.0))
\end{lstlisting}
\code{chi_sector} requires \code{transport="full"} (the semi-coupled and QKE paths
have no $\chi$ block).  The integrator is selected with \code{solver=}: the default
\code{"tsit5"} (explicit) is fine in this regime, and \code{"implicit"} routes to
the IMEX solver below for added robustness.  Nothing on the abundance side changes;
$\chi$ enters only through $\rho_\chi$ in the Hubble rate and through the distortion
it leaves in $f_\nu$.

\subsection{Strong coupling and the exact-shape solver}
\label{sec:chinustrong}

The hard case is strong/stiff coupling, large $G_{\rm eff}$ (e.g.\ a
weak-scale mediator giving $K\sim100\,\GF^2$) so that $\Gamma\gg H$ by hundreds.
The difficulty is physical, not cosmetic: when $\chi$ annihilates to $\nu\bar\nu$
it deposits a genuine super-thermal feature in the neutrino spectrum (a
line / double-peak near comoving $\eps\sim m_\chi/\Tcm$) that no near-equilibrium
treatment captures.  Getting $\Neff$ and the abundances right requires
delivering the true, signed annihilation spectrum shape.  A thermal or a
two-moment (number$+$energy only) surrogate is a corner-cut on the spectral
shape, and the shape matters at the BBN floor (Section~\ref{sec:chinustatus}).
The solver therefore carries the true signed shape directly; no approximate
(thermal, two-moment, or energy-only) delivery is offered.

The strong-coupling deliverable is \code{hybrid.make_chi_exact_solver}, a
fixed-step ($\ln a$) IMEX that evolves $(f_\nu,f_\chi)$ under the true coupled
kernel net (one \code{chi_brute.s2_brute_rates} call returns both
$\mathrm{net}_\nu$ and $\mathrm{net}_\chi$), so $\nu$ and $\chi$ number and energy
are conserved to the kernel-quadrature floor.  To drive it through LINX for
abundances, pass the $\chi$ sector to \code{make_coupled_bbn(chi_sector=...,
transport="full", solver="exact")} (Section~\ref{sec:linx}).

\paragraph{Sub-cycling (\texttt{k\_sub}, default 8): why.}
The super-thermal bump must thermalize (through $\nu$--$\nu$ scattering)
and leak its energy into the $e^\pm$ plasma (through $\nu$--$e$ scattering)
on its own physical timescale.  So within each outer $\chi$ step the SM
$\nu+\Tpl$ collision evolution is sub-cycled \code{k_sub} times, recomputing the
SM rates at the actual bumped state each sub-step.  Without this, the
isolated background $\Tpl$ step cannot follow the feedback and $\Neff$ diverges.
This blow-up is \textbf{genuine numerical stiffness}, not a bookkeeping or
energy-conservation error, and not grid under-resolution: refining the
$\eps$-grid does not cure it.

\medskip
Three ingredients make the deep tail ($x\equiv m_\chi/\Tcm\gtrsim1.6$) stable, and
each is necessary: with any one removed the strong solve either saturates the
relic or blows $\Neff$ up.  They are the detailed-balance evaluation, the
rank-revealing linear solve, and the conservation-gated freeze, in turn.

\paragraph{Detailed-balance evaluation at $\Gamma\gg H$ (the \texttt{expm1} form).}
Deep in the strong-coupling tail the forward and reverse rates of every bin agree to
machine precision (that is detailed balance at $\Gamma\gg H$), so the naive
net $=\mathrm{forward}-\mathrm{reverse}$ loses its sign to round-off: the relic
abundance then collapses or saturates path-sensitively, and the noise seeds a spurious
super-thermal high-$\eps$ tail that the $\eps^3$ energy weight amplifies into a fake
$\Neff$.  The
kernel (\code{chi_brute._Ffac}) therefore evaluates the net in the
cancellation-free detailed-balance form
\[
  \mathrm{net}=-\,\mathrm{forward}\cdot\mathrm{expm1}\!\big(\log(\mathrm{reverse}/\mathrm{forward})\big),
\]
with the log-ratio summed from the individual Pauli-blocking log factors, algebraically
identical to the subtraction but accurate to machine precision exactly where it nearly
vanishes (the $\chi$-sector analogue of the SM kernels' \code{prec} normalisation of the
same forward$-$reverse cancellation).  It is gated to $|\delta|<1$ (the direct difference
is used where one leg dominates or is zero, for overflow safety) and is smooth, hence
\code{jacfwd}-safe inside the coupled Newton.  It alone takes the spurious $\Neff$ from
$\sim$58 down to $\sim$10; the next two ingredients remove the rest.

\paragraph{The default linear solve: a rank-revealing Newton (\texttt{robust=True}).}
On a super-thermal $f_\nu$ state the coupled log-occupation Jacobian is
near-rank-deficient: a positive-feedback growing mode (the same lepton-kernel
boundary amplification, Section~\ref{sec:chinufix}) sits at the backward-Euler
stability boundary, so an ordinary \code{solve} returns a meaningless result there and the step
is rejected.  The default \code{robust=True} replaces \code{solve} with a
truncated-SVD least-squares step (\code{lstsq}, \code{rcond=1e-6}) that truncates
that spurious mode at the linear-solve level and takes the minimum-norm step in the
well-conditioned subspace.  \code{rcond} is a numerical conditioning floor (insensitive
over $[10^{-8},10^{-5}]$), not a physics knob.  This is the principled,
no-amplitude-bound solver, and it reproduces $\Neff\approx3.99$, relic~$\approx0.60$
for the pinned corner $m_\chi=1$ MeV, $K=100\,\GF^2$ with no clip at all.

\paragraph{The opt-in evolved-spectrum cap (\texttt{cap\_evolved}).}
A too-loose value of the cap constant gives a cap-dependent, non-reproducible
$\Neff$: it swings $\sim0.5$ with the constant, GPU reduction non-determinism
flipping the near-cancellation (a spurious $\Neff=3.535$ is one such reading).
Fixing the cap on the evolved spectrum at $f\le e^{0.5}f^{\rm eq}_\nu
\approx1.65\,f^{\rm eq}_\nu$ (just above the physical injection peak) via
\code{cap_evolved=True}, \code{fdcap_logC=0.5} removes that: it is cap-insensitive
($\dNeff$ flat to $0.005$) and bit-reproducible, $\Neff\approx3.99$, with a
monotone coupling scan (9/9).  It is an opt-in alternative; the default
\code{robust} Newton resolves the same near-cancellation instability at the
linear-solve level with no amplitude bound.

\paragraph{Cost and the \texttt{solve\_first} fast path.}
\code{lstsq} (SVD) is $\sim33\times$ slower than \code{solve} (LU) at the
$707$-dimensional Jacobian on GPU (the SVD parallelises poorly), so robust runs
$\sim2$--$3\times$ the opt-in cap path.  \code{solve_first=True} (opt-in) recovers most
of that: it tries the cheap \code{solve} first and falls back to the SVD \code{lstsq}
only on the rare near-singular steps.  The price is a measured $\sim1\%$
accuracy loss (in the partial-progress gray zone it keeps a sub-optimal solve step
instead of the exact lstsq one).  A deliberate speed/accuracy trade; default off
keeps the exact path.

\paragraph{The regime split and the conservation-gated freeze backstop
(\texttt{deepfix="freeze"}).}
The solver partitions on $x\equiv m_\chi/\Tcm$: it runs the coupled cross-block
Newton only inside the active window $x\in[$\code{x_switch}$,$\code{x_freeze}$]$
(defaults $0.5,4.0$), pins $\chi$ to its equilibrium spectrum below the window,
and holds it above.  For very heavy $\chi$ or very strong $K$ the deep freeze tail
($x>$\code{x_freeze}) is where the near-cancellation bites hardest; the backstop
is the conservation gate \code{deepfix="freeze"}: after the Newton, a step
that violates $\nu\leftrightarrow\chi$ number or energy beyond \code{cons_tol},
or fails to converge / goes non-finite, is rejected (hold $f_\nu,f_\chi$ $\Rightarrow$
zero deposit $\Rightarrow$ trivially conserved).  Energy gating is essential:
a number-only gate lets energy-creating steps through and $\Neff$ diverges.
\code{deepfix="none"} disables it and is unstable past $x\sim1.6$.

\subsection{Status: the open relic-depletion shape systematic}
\label{sec:chinustatus}

The sector is documented but not yet closed.  The delivery shape is a real
systematic: a moment-matched (number$+$energy) surrogate that is energy-exact but
$\nu$-number short shifts the abundances by roughly $1.2$--$2.3\%$ on $\DoverH$ and
$0.3$--$2.5\%$ on $\Yp$ (growing with $m_\chi$ and coupling), right at the BBN
observational floor ($\Yp\pm0.003\approx1.2\%$, Aver et al.\ 2021;
$\DoverH\pm0.030\times10^{-5}\approx1.2\%$, Cooke et al.\ 2018).  This is why
\code{make_chi_exact_solver} carries the true signed shape rather than a surrogate.

The lesson is worth stating plainly, because it is easy to get wrong:
\textbf{``energy-exact $\Rightarrow\Neff$-exact'' is false.}  At fixed
delivered number and energy, two different delivery spectra can still differ
in $\Neff$ by $0.10$--$0.30$: at $m_\chi=1$, $K=100$ two moment-matched surrogates
give $\Neff\approx3.686$ and $\approx3.804$, both short of the exact solver's
$\approx3.99$, because the final $\Neff$ depends on where in $\eps$ the
deposited energy lands.  This is exactly why only the true signed shape is
acceptable, and why \code{make_chi_exact_solver} (which closes the number mismatch
by evolving the full coupled net) is the right answer rather than a moment-matched
surrogate.  It is the production solver, and the default \code{robust}
(rank-revealing) Newton resolves the cap dependence at the linear-solve level.  What remains ``in progress'' is
the deeper lepton-kernel source fix below, which would remove the remaining
FD-envelope/window caps and the heavy-$\chi$ ($m_\chi=3$) corner uncertainty at the
source.

\subsection{The in-progress lepton-kernel source fix}
\label{sec:chinufix}

The strong-$\chi$ path has three cap devices: the FD-envelope cap
(\code{fd_cap_logC}, Section~\ref{sec:qkecap}), the window-max
\code{range_limit} (Section~\ref{sec:darkknobs}), and the opt-in
\code{cap_evolved} above.  All three are stopgaps for the same underlying
problem.  (The default \code{robust} Newton sidesteps \code{cap_evolved} by
truncating the spurious mode at the linear solve rather than clipping the spectrum,
but the underlying off-grid overshoot is still there for the other two and is what a
true source fix would remove.)  When the lepton and
$\nu$--$\nu$ kernels (\code{nuer2}, $\nu\nu$/$\bar\nu\nu$, and the $\nu\chi\to\nu\chi$
scatter leg) need an occupation just past the top of the energy grid, the
off-grid order-5 log-occupation extrapolation in \code{interp.py} can
overshoot upward into an unphysical super-thermal spike (the large Lagrange
weights amplify the curvature), which those kernels then amplify, the documented
lepton-kernel boundary fragility (Section~\ref{sec:knownissues}), H-theorem-violating
on non-smooth spectra.

The principled cure is to replace that off-grid extrapolation
with a physical relativistic-Fermi--Dirac decay closure: above the
top grid node, force $\log f$ to fall off linearly in $\eps$ at (no shallower
than) the local slope, a monotone, H-theorem-respecting tail that cannot
overshoot upward.  Because the SM decoupling spectrum is already sub-thermal
off-grid, this leaves the SM results bit-identical, while removing the need
for the remaining caps and the residual $m_\chi=3$ strong-coupling limitation in one
stroke.  It is already exposed opt-in (\code{fd_decay}/\code{monotone_read}, default
off and SM-bit-identical); making it the default, validated across the $\chi$ grid,
is the active development frontier for the sector.

\begin{diffbox}
\textbf{Validation contract.}  There is no trustworthy Fortran end-to-end oracle
for the $\chi$ sector, so it is gated by a brute-force collision-integral oracle
(\code{chi_brute.s2_brute_rates}; \code{tests/test_chi_s2_kernels.py}) plus exact
conservation invariants (detailed balance at a common temperature, and the
cross-species $\nu\leftrightarrow\chi$ number and energy ledgers) plus the $K\to0$
reduction to the validated SM transport.
\end{diffbox}

% ===========================================================================
\section{Validation of the dark sectors}
\label{sec:darkval}
% ===========================================================================

The dark sectors follow the contract of Section~\ref{sec:darkvalidation}: no
trustworthy Fortran end-to-end oracle, so unit oracles $+$ exact invariants $+$
the classical reduction.  The concrete gates:
\begin{itemize}[nosep]
\item \textbf{Decay kernels:} a brute-force direct evaluation of the
  $1\!\leftrightarrow\!2$ collision integral; the dilute-limit reduction to the
  analytic widths $\Gamma_{\rm vec}=g^2 m/24\pi$, $\Gamma_{\rm sc}=g^2 m/8\pi$;
  the cross-species number/energy ledger.  Pytest: \code{tests/test_dv_decay.py},
  \code{tests/test_dark_decay_coupled.py}, and the data types in
  \code{tests/test_dark_foundations.py}.
\item \textbf{$\chi\leftrightarrow\nu$ contact kernels:} a brute-force oracle
  (\code{chi_brute.s2_brute_rates}), detailed balance at a common temperature, the
  cross-species $\nu\leftrightarrow\chi$ number/energy ledgers, and the $K\to0$
  reduction to the validated SM transport (Section~\ref{sec:chinu}).  Pytest:
  \code{tests/test_chi_s2_kernels.py}, the gated coupled RHS in
  \code{tests/test_chi_s2_coupled.py}, and the strong-coupling solver in
  \code{tests/test_chi_s2_exact.py} (all three of its tests are \code{slow}).
\item \textbf{Conservation:} the solvers' own ledgers (\code{num_rel_max}/
  \code{rho_rel_max} of \code{make_chi_exact_solver},
  \code{decay_residual_frac}/\code{number_residual_frac} of
  \code{make_dark_decay_implicit_solver}) and (decay) the
  \code{eps_tail_energy_frac}/\code{eps_top_energy_frac} energy-space monitors.
\item \textbf{SM reduction:} with the sectors off every solver is the SM one; a
  zero-coupling decay sector leaves the RHS bitwise unchanged
  (\code{tests/test_dark_decay_coupled.py}).
\end{itemize}
The flavour-mixing QKE sector is validated separately (it is a Standard-Model
refinement, Section~\ref{sec:qke}): \code{gfortran} unit oracles for the
weak potential ($m^2_{\rm weak}$, rel $10^{-16}$) and the coherent commutator (rel
$10^{-15}$), the matrix exponential against fine RK4 and analytic two-flavour
oscillation, and the mixing-off capstone that reduces the whole machinery
bit-for-bit to the classical transport.

\paragraph{Canonical dark numbers}
\begin{center}\small
\begin{tabular}{@{}R{0.40\textwidth}R{0.55\textwidth}@{}}
\toprule
\textbf{result} & \textbf{value} \\
\midrule
OOE scalar (light majoron, $m_X\!\lesssim\!1$, $\sim$100 keV decay) & $\Neff=3.079$, $\Yp=0.2470$, $\DoverH=2.45\times10^{-5}$ \\
thermal/overabundant majoron ($m_J{=}5$, $5\times$BE) & $\dNeff=+2.43$ (energy/number conserving) \\
strong $\chi\leftrightarrow\nu$ ($m_\chi{=}1$, $K{=}100\,\GF^2$) & $\Neff\approx3.99$ (interim; $\sim$1--2.5\% shape systematic on $\Yp$/$\DoverH$) \\
\bottomrule
\end{tabular}
\end{center}
The $\dNeff=+2.43$ entry is a differential and QED-baseline-independent (unchanged
by the NLO default).  Of the two absolute $\Neff$, the OOE-scalar $3.079$ is at the
LO baseline and shifts by the NLO $\approx-9.5\times10^{-4}$ ($\to3.078$), while the
strong-$\chi$ $\approx3.99$ carries a $\sim$1--2.5\% shape systematic
(Section~\ref{sec:chinustatus}) that dwarfs any NLO baseline shift.  The
genuine-QKE numbers (the vacuum $\dNeff=+7.6\times10^{-4}$, the pinned comparison,
the matrix-collision shift, the $\sim10^{-5}$
coherences) live with their derivation in Section~\ref{sec:qke}.

% ===========================================================================
\section{Scanning and extending the dark sectors}
\label{sec:darkextend}
% ===========================================================================

The general guidance of Section~\ref{sec:extend} (precision, gradients, scans,
fixed structure) applies unchanged; this section adds what is specific to the
dark sectors.

\paragraph{Scanning a dark model}  The \code{BosonDecaySector} and
\code{ChiSector} fields (\code{m_x}, \code{gsq_diag}, \code{tratio}; \code{m_chi},
\code{Kflav}) are dynamic \code{equinox} leaves, so a mass/coupling scan compiles
once and re-runs; build the solver with a sector of the right shape and vary the
leaves.  Vary the initial abundance through \code{f_x0} / \code{f_chi0} (a scaled
equilibrium shape for a frozen-out or over-abundant species).

\paragraph{Adding a new BSM interaction}  The recipe mirrors how the decay and
$\chi$ sectors were added:
\begin{enumerate}[nosep]
\item \textbf{Derive the squared matrix element} $|M|^2$ for the new
  $\nu$-coupled process (the headless-FeynCalc workflow in \code{dark_amplitudes/}
  on the BURST repository's \code{dark} branch is the template, with its internal
  $+$ literature gates).
\item \textbf{Angular-reduce it} to a per-output-bin energy quadrature and write
  it as a kernel returning $(\mathrm{net},\mathrm{frs})$, following the existing
  templates: a $1\!\leftrightarrow\!2$ (inverse-)decay like \code{dv_nu}, or a
  $2\!\to\!2$ contact like the $\chi$/SM-lepton family.
\item \textbf{Wire it into the RHS with the correct gating}
  (Section~\ref{sec:darkgate}): a smooth $1\!\leftrightarrow\!2$ kernel is
  added ungated after the SM gate; a cancelling $2\!\to\!2$ kernel
  keeps the $\Tcm^5$ convention and gets its own separate gate (never
  folded into the SM denominator).  Add the new species' $\rho$ to the Hubble
  rate (the \code{rho_extra} of \code{decoupling._background}), and carry its
  occupation as a new state component.
\item \textbf{For a stiff sector}, add the affine $A-Bf$ form so the
  exponential-midpoint head can integrate it (the dark decay heads are the
  template); for an out-of-equilibrium decay, the redshifting-feature machinery
  of Section~\ref{sec:darksolvers} applies.
\item \textbf{Validate} against a brute-force evaluation of the collision
  integral and the conservation sum rules (the dark sectors have no Fortran
  oracle), and check that the sector switched off leaves the SM pins bitwise
  unchanged.
\end{enumerate}
The nearest worked example is the $\chi\leftrightarrow\nu$ contact sector
(Section~\ref{sec:chinu}): its amplitudes, kernels, and strong-coupling solver
exist and are usable; the one open piece is the principled off-grid closure that
would retire the cap band-aids (Section~\ref{sec:chinufix}).

\appendix
% ===========================================================================
\section{BURST $\leftrightarrow$ \texttt{neutt} mapping index}
\label{sec:fortranmap}
% ===========================================================================

This appendix collects the complete correspondence between \texttt{neutt} and
Fortran BURST in one place, so the two trees can be cross-read.  The inline
\textbf{[Fortran $\leftrightarrow$ JAX]} and \textbf{[Port note]} boxes give the
narrative; the tables here give the full enumeration.  The \textbf{status} column
uses:

\begin{center}\small
\begin{tabular}{@{}ll@{}}
\toprule
\textbf{status} & \textbf{meaning} \\
\midrule
\textsc{same} & bit-faithful transcription (value-identical to the snapshot oracle) \\
\textsc{diff} & same physics, deliberately different numerical route (noted why) \\
\textsc{none} & BURST feature not ported \\
\textsc{add}  & port-only addition with no BURST counterpart \\
\bottomrule
\end{tabular}
\end{center}

A note on the two Fortran trees: most rows cite the snapshot in \code{fortran_reference/},
against which the port is validated (the constants and the equilibrium occupations
\code{fd_equil_calc}, \code{be_equil_calc} are in its \code{mainvar.f90}; \code{renorm.f90} is
the same file, with the same line numbers, in both trees).  Rows that name \code{part_lib.f90}
(\code{part_make_bins}, \code{part_bin_var}, \code{part_var}), \code{bbn_therm.f90} or
\code{get_early_derivs}, or say ``refactored tree'', cite the refactored BURST tree on the
repository's \code{main} branch; the snapshot counterparts are the grid loop of
\code{trans_init.f90} (lines 168--178, without the offset options), \code{bbn_yes_therm_v2.f90},
and \code{bbn_yes_evolve_v2.f90}.

\subsection{Auxiliaries and grid}
\begin{center}\small
\begin{longtable}{@{}R{0.26\textwidth}R{0.30\textwidth}l R{0.23\textwidth}@{}}
\toprule
\textbf{\texttt{neutt}} & \textbf{Fortran} & \textbf{status} & \textbf{note} \\
\midrule\endhead
\code{constants.*} & \code{mainvar.f90}, \code{transvar.f90} & \textsc{same} & values verified exact (\code{ZETA3}: full double precision; Section~\ref{sec:constants}) \\
\code{grid.make_eps_grid} & \code{part_lib.f90::part_make_bins} & \textsc{same} & linear $\eps$ grid $+$ Boole weights \\
\code{grid.EpsGrid} & \code{part_bin_var} / \code{part_var} & \textsc{same} & data model \\
\code{grid.fermi_dirac} & \code{fd_equil_calc} & \textsc{same} & \\
\code{grid.bose_einstein} & \code{be_equil_calc} & \textsc{same} & guarded \\
\code{grid.log_occupation} & \code{interp_occfrac} clamp & \textsc{same} & $[-300,0]$ clamp \\
\code{grid.legendre_canonical} & \code{gauss.f90} (Legendre) & \textsc{diff} & host NumPy, not \code{imtqlx} \\
\code{grid.gauss_laguerre*} & \code{gauss.f90} / \code{cgqf} & \textsc{diff} & host SciPy, not \code{imtqlx} \\
\code{grid.rescale_legendre} & \code{gauss.f90::scqf} & \textsc{same} & the affine rescale (ported; the kernels apply it inline) \\
\code{grid} moments (\code{integrate}, \code{moment}, \code{number_density}, \code{energy_density}) & inline moment integrals & \textsc{same} & Boole-weighted sums \\
\code{interp.lagrange_window}/\code{apply_log_interp} & \code{interp_occfrac}, \code{polint}, \code{indpolint} & \textsc{same} & static-query path \\
\code{interp.interp_log_occ_dynamic} & \code{interp_occfrac} (3 branches) & \textsc{same}/\textsc{diff} & snapshot clamp; \code{fd_cap_logC}/\code{fd_decay}/\code{monotone_read} are \textsc{add} \\
\bottomrule
\end{longtable}
\end{center}

\subsection{Background and finite-$T$ QED}
\begin{center}\small
\begin{longtable}{@{}R{0.26\textwidth}R{0.30\textwidth}l R{0.23\textwidth}@{}}
\toprule
\textbf{\texttt{neutt}} & \textbf{Fortran} & \textbf{status} & \textbf{note} \\
\midrule\endhead
\code{thermo.rho_gamma}/\code{pressure_gamma} & \code{bbn_therm.f90} & \textsc{same} & \\
\code{thermo.rho_electron}/\code{pressure_electron} & \code{glag4}/\code{glag7} leading terms & \textsc{diff} & 100-node default vs.\ BURST's 25 \\
\code{thermo.rho_plasma}/\code{pressure_plasma} & \code{bbn_therm.f90} & \textsc{same} & bare mass; QED added separately \\
\code{thermo.drho_plasma_dT} & analytic $f(1-f)$ & \textsc{same} & cross-checked vs.\ autodiff \\
\code{thermo.get_eg_mass_renorm} & \code{renorm.f90:29--114} & \textsc{same} & $\delta m_e^2,\delta m_\gamma^2$ \\
\code{thermo.get_e_mass_pdep} & \code{renorm.f90:119--282} & \textsc{same} & momentum-dependent self-energy \\
\code{thermo.m_e_eff} & \code{renorm.f90} (kernel mass set bare at \code{trans_evolve.f90:97}) & \textsc{same} & thermal $m_{\rm eff}$ exposed, but kernels use the bare mass \\
\code{thermo.qed_eos_corrections} & \code{glag1--7} (momentum-indep.) & \textsc{same} & additive $\delta\rho/\delta P$ \\
\code{thermo.qed_eos_pdep_corrections} & \code{glag4/5/7} pdep (2nd loop) & \textsc{same} & \\
\code{thermo.qed_plasma_deltas} & \code{bbn_therm.f90} EoS block & \textsc{same} & one-call form for the RHS \\
\code{thermo.qed_eos_e3_pressure}/\code{qed_eos_e3_deltas} & --- (LINX \code{explicit_P3}) & \textsc{add} & NLO $\mathcal O(e^3)$ pressure and its $\delta\rho$, $\delta C_{\rm pl}$; default on (\code{qed_nlo}) \\
(\code{glag6}, $\dd\rho_e/\dd\phi_e$ pdep) & \code{bbn_therm.f90} & \textsc{none} & EoS-derivative term unported; $\phi_e$ itself is carried (settable, default $0$) \\
\code{thermo.hubble} & \code{bbn_yes_evolve_v2.f90:156} (\code{hubcst}) & \textsc{same} & MeV$^4$ units (no \code{mev4tocgs}) \\
\code{thermo.evolve_background} & \code{get_early_derivs} path & \textsc{diff} & $\ln a$ ODE (diffrax) vs.\ time RK \\
\bottomrule
\end{longtable}
\end{center}

\subsection{Transport kernels and assembly}
\begin{center}\small
\begin{longtable}{@{}R{0.26\textwidth}R{0.30\textwidth}l R{0.23\textwidth}@{}}
\toprule
\textbf{\texttt{neutt}} & \textbf{Fortran (snapshot)} & \textbf{status} & \textbf{note} \\
\midrule\endhead
\code{transport.nunu} & \code{trans_nunu.f90} & \textsc{same} & grid-only static tensor \\
\code{transport.nunubar} & \code{trans_nunubar.f90} & \textsc{same} & scatter $+$ annih channels \\
\code{transport.nuer1} & \code{trans_nuer1.f90} & \textsc{same}/\textsc{diff} & background-dependent (union$+$mask); default \code{lag_tail_fix=True} switches the wide-regime tail to Gauss--Laguerre ($\Neff$ shift $\le8.4\times10^{-7}$; \textsc{same} at \code{lag_tail_fix=False}) \\
\code{transport.nuer2} & \code{trans_nuer2.f90} & \textsc{same}/\textsc{diff} & background-dependent; default \code{lag_tail_fix=True} replaces the wide R1 segment at the lowest bins by a Gauss--Laguerre rule ($\Neff$ shift $<10^{-9}$; \textsc{same} at \code{lag_tail_fix=False}) \\
\code{transport.nuepma} & \code{trans_nuepma.f90} & \textsc{same}/\textsc{diff} & background-dependent; default \code{lag_tail_fix=True} lowers R2's genout Legendre$\to$Laguerre switch to \code{lag_tail_gap}, so the wide segment just above \code{pcase1} takes the Gauss--Laguerre tail ($\Neff$ shift with \code{nuer1} $\le8.4\times10^{-7}$; \textsc{same} at \code{lag_tail_fix=False}) \\
\code{transport.nuer1/nuer2/nuepma} \code{*_rates_sched} & --- & \textsc{add} & launch-collapsed XLA rewrite: \code{select=False} (\code{"union"}) reproduces the reference to $\le10^{-15}$, \code{select=True} to $\sim$$6.7\times10^{-7}$ in $\Neff$ \\
\code{assemble.assemble_rhs} & \code{trans_evolve.f90:195--208} & \textsc{same} & combined gate (snapshot, not the refactored tree) \\
\code{assemble.assemble_split_rhs} & --- & \textsc{add} & $A-Bf$ affine form (for the hybrid head) \\
\code{assemble.compute_prec_trans} & \code{trans_init.f90:268--286} & \textsc{same} & $+$ \code{optimization_barrier} (\textsc{add}, runtime identity) \\
\code{assemble.equilibrium_occ} & \code{trans_init.f90:185--196} & \textsc{same} & \\
\code{assemble.compute_sumrules} & BURST sumrule1--5 & \textsc{same} & diagnostic \\
\code{assemble} \code{gate_width}, \code{fd_cap_logC} & --- & \textsc{add} & robustness knobs, default-off (no-ops) \\
\code{assemble} \code{OCC_FLOOR} & --- & \textsc{add} & always-on positivity clamp (identity on physical $f>10^{-300}$) \\
(combined vs.\ per-channel gate) & \code{trans_evolve.f90} (refactored tree) & \textsc{diff} & port follows snapshot's combined gate \\
\bottomrule
\end{longtable}
\end{center}

\subsection{Solver, driver, and coupling}
\begin{center}\small
\begin{longtable}{@{}R{0.26\textwidth}R{0.30\textwidth}l R{0.23\textwidth}@{}}
\toprule
\textbf{\texttt{neutt}} & \textbf{Fortran} & \textbf{status} & \textbf{note} \\
\midrule\endhead
\code{decoupling.make_decoupling_solver} & \code{bbn_yes} driver $+$ \code{trans_evolve} & \textsc{diff} & $\ln a$ / diffrax; $\{T_{\rm pl},f\}$ state; built once, jitted \code{solve} \\
\code{decoupling.make_rhs} & \code{bbn_yes_evolve_v2.f90} $+$ \code{trans_evolve.f90} & \textsc{same} & temperature ODE $+$ gated collision RHS \\
\code{decoupling.make_split_rates} & --- & \textsc{add} & affine RHS for the hybrid head \\
\code{decoupling.nanrobust_rms_norm} & --- & \textsc{add} & diffrax NaN-freeze guard \\
\code{hybrid.make_hybrid_solver} & --- & \textsc{add} & exp-midpoint head $+$ Tsit5 tail \\
solver: \code{Tsit5}/\code{Kvaerno5} & Cash--Karp RK5(4) & \textsc{diff} & adaptive diffrax solvers \\
\code{transport_adapter.*} (LINX) & --- & \textsc{add} & transport$\to$LINX seam (the whole module) \\
\code{linx.coupled_bbn.make_coupled_bbn} & BURST internal BBN $+$ \code{sdratio} & \textsc{diff}/\textsc{none} & abundances via LINX; no recombination \\
nuclear network & \code{nuc_*}, NSE, solver & \textsc{none} & LINX provides abundances \\
recombination / sound horizon & \code{sdratio*} & \textsc{none} & not ported \\
sterile $\nu$ ($n_\nu>3$) & slots $4,\dots,n_\nu$ & \textsc{none} & SM transport is 3 flavours \\
MPI / \texttt{\#ifdef prllel} & \code{mpi_allreduce} & \textsc{none} & replaced by on-GPU contraction \\
$\phi_e(t)$ co-evolution & \code{phie} ODE & \textsc{diff} & ported as the algebraic charge-neutrality $\phi(\Tpl)$ (\code{thermo.solve_phie} via \code{phie="auto"}, default $0$); the glag6 EoS-derivative closure is unported \\
\bottomrule
\end{longtable}
\end{center}

\subsection{Flavour mixing (QKE)}

The QKE sector is new code with \code{qburst} as its lineage, not a transcription:
the end-to-end solve of \code{qburst} is not a trustworthy oracle
(Section~\ref{sec:qke}), so the sector is validated against unit oracles,
conservation invariants, and the mixing-off reduction rather than the snapshot.

\begin{center}\footnotesize
\begin{longtable}{@{}R{0.27\textwidth}R{0.27\textwidth}l R{0.26\textwidth}@{}}
\toprule
\textbf{\texttt{neutt}} & \textbf{origin / oracle} & \textbf{status} & \textbf{note} \\
\midrule\endhead
\code{transport.qke.coherent} & qburst \code{trans_init} / \code{get_coherent} & \textsc{add} & PMNS, $m^2_{\rm weak}$, $H$ and $-i[H,\rho]$; \code{gfortran} unit oracles (\code{oracle_qke_mixing}, \code{oracle_qke_coherent}) \\
\code{transport.qke.nonadiabatic} & qburst corotating step (\code{get_off_diag_amps} oracle) & \textsc{add} & coherent-step propagators (\code{expm}, Strang, \code{hybrid}) and the moving eigenframe \\
\code{transport.qke.collide}, \code{collide_nue}, \code{collide_nunu} & classical \code{nuepma}; \code{nuer1}/\code{nuer2}; \code{nunu}/\code{nunubar} & \textsc{same}/\textsc{add} & matrix generalizations (shared tensors in \code{collide_matrix}); on a diagonal $\rho$ each reduces to its classical kernel and sources no coherence; the off-diagonal is new \\
\code{assemble.decoherence_rate} & --- & \textsc{add} & ungated $g+l$, the scalar coherence damping $\Gamma_a$ \\
\code{transport.qke.full_coupled}, \code{solve} & qburst (diverges e2e) & \textsc{add} & the density-matrix solver, standalone Strang driver; \code{mixing=False} reduces bit-identically to the classical solve \\
\bottomrule
\end{longtable}
\end{center}

\subsection{Summary of port-only additions}

Things \texttt{neutt} has that BURST does not (the BURST-faithful gold gates are
preserved because each is default-off, a runtime identity, or reversible with a
single flag):
\begin{itemize}[nosep]
\item the NLO $\mathcal O(e^3)$ finite-$T$ QED correction
  (Section~\ref{sec:nloqed}), the one default-on new-physics
  addition: it moves $\Neff^{\rm SM}$ from the LO $3.0443$ to $3.0433$, and
  \code{qed_nlo=False} recovers the LO value (BURST's own $3.0442$ also needs \code{SIN2THW_BURST});
\item the \code{lag_tail_fix} wide-regime lepton-kernel tail quadrature correction
  (Section~\ref{sec:lepkernels}), default on, but numerical: it fixes a
  latent Fortran-BURST quadrature bug in \code{nuer1}, \code{nuer2} and
  \code{nuepma} ($\Neff$ shift $\le8.4\times10^{-7}$), and
  \code{lag_tail_fix=False} recovers the bit-faithful BURST form;
\item the kernel-tensor \code{einsum} reformulation and the union-grid$+$mask
  scheme (Section~\ref{sec:kerneltensor});
\item the scheduled (\code{_sched}) launch-collapsed kernels and \code{n_chunks}
  partner-bin chunking (performance only; the \code{"union"} schedule and
  \code{n_chunks=1} are bit-faithful, the \code{select} schedule agrees to
  $\sim$$6.7\times10^{-7}$ in $\Neff$);
\item the affine split $A-Bf$ and the exponential-midpoint hybrid solver
  (Sections~\ref{sec:makerhs}, \ref{sec:hybrid});
\item the \code{optimization_barrier} gate guard, \code{nanrobust_rms_norm}, and
  the \code{sqrt}/FD/denominator \code{where}-guards (correctness/robustness);
\item the opt-in \code{gate_width} smooth gate and \code{fd_cap_logC} FD-envelope
  cap (default no-ops), and the always-on \code{OCC_FLOOR} positivity clamp (an
  identity on physical occupations);
\item full \code{jit}/\code{grad}/\code{vmap} differentiability and GPU batching;
\item the build-once solver factories (\code{make_*_solver} returning a jitted
  \code{solve}) and the compile banner (Section~\ref{sec:knobs});
\item the QKE flavour-mixing sector and its off-grid tail closures \code{fd_decay}
  and \code{monotone_read} (Section~\ref{sec:qke}), which leave the classical solve
  untouched (\code{mixing=False} is bit-identical to it);
\item the entire LINX coupling layer (Section~\ref{sec:linx}).
\end{itemize}

% ===========================================================================
\section{Glossary of symbols}
\label{sec:glossary}
% ===========================================================================

Symbols are grouped by topic; code names are in typewriter.  Energies are in MeV,
energy densities and pressures in MeV$^4$, and rates in s$^{-1}$ unless noted.  Where
one letter carries two meanings in different sections, both are listed.

\begin{center}\small
\begin{longtable}{@{}>{\raggedright\arraybackslash}p{0.27\textwidth}p{0.66\textwidth}@{}}
\toprule
\textbf{symbol} & \textbf{meaning} \\
\midrule\endhead
\multicolumn{2}{@{}l}{\textit{Temperatures, clocks, and the grid}} \\
$\Tcm$ (\code{T_cm}, \code{tcmev}) & comoving temperature parameter (redshift clock, $\Tcm=\Tcm^{\rm init}e^{-\ln a}\propto1/a$); not a physical $\Tnu$ \\
$\Tpl=\Tg$ (\code{T_pl}, \code{tplmev}) & physical plasma / photon temperature \\
$\Tnu$ & physical neutrino temperature (thermal-fluid approximation only); $f^{\rm eq}$ is Fermi--Dirac at $\Tnu=\Tcm$ \\
\code{tcmpl}, \code{ratio} & $\Tcm/\Tpl$ (a lepton-kernel argument), and its inverse $\Tpl/\Tcm$ returned by the solvers \\
$\ln a$, $s$ & integration variable ($e$-folds since the start); $\Tcm=\Tcm^{\rm init}e^{-\ln a}$.  Written $s\equiv\ln a$ in Sections~\ref{sec:thesystem} and~\ref{sec:qkesplit} and in the affine form $\dd f/\dd s=A-Bf$; the QKE sources $s$ and $s_{\rm coll}$ are below \\
\code{T_init}, \code{lna_max}, $T_{\rm end}$ & transport start temperature ($\Tpl=\Tcm$ there; default 10), the solve length in $\ln a$, and $T_{\rm end}=\code{T_init}\,e^{-\code{lna_max}}$ \\
\code{tcm_cutoff} & transport-window low edge in $\Tcm$ (BURST's 0.015); below it the rates are masked to zero and the spectrum is frozen \\
$\eps=E_\nu/\Tcm$ & comoving energy; the grid abscissa ($=p/\Tcm$ for the relativistic neutrinos); $\eps_i$, $\eps_a$ a bin's value \\
$N$, $\eps_{\max}$, $w_i$, $h$ & bin count \code{nbins} ($4n{+}1$; 101), grid extent \code{eps_range} (25), composite-Boole weights, and spacing $h=\eps_{\max}/(N-1)$ (Section~\ref{sec:grid}); the $h$ of Eq.~\eqref{eq:expstep} is instead a step in $\ln a$ \\
$p_1$, $E_2$, $E_3$, $p_4$ & the legs of a $1{+}2\to3{+}4$ process in $\Tcm$ units: $p_1=\eps_i$ the output bin, $E_2,E_3$ the integrated energies, $p_4$ the off-grid fourth leg ($p_1+E_2-E_3$ for scattering; $E_{\rm out}+E_{\rm in}-p_1$ in \code{nuepma}) \\
\midrule
\multicolumn{2}{@{}l}{\textit{Occupations, densities, and observables}} \\
$f_{\nu_i}(\eps)$, $f[3,2,N]$ & neutrino occupation number, indexed flavour ($0{:}e,1{:}\mu,2{:}\tau$) $\times$ cp $\times$ bin; equilibrium $f^{\rm eq}=1/(e^\eps+1)$ \\
cp, \code{parfact} & charge-parity index ($0{:}\nu$, $1{:}\bar\nu$) and its sign $\pm1$ \\
$\delta f$, $\delta_a$ & fractional spectral distortion $(f-f^{\rm eq})/f^{\rm eq}$; $\delta_a$ that of flavour $a$ (e.g.\ $\delta_e-\delta_\mu$) \\
$\xi=(\xi_e,\xi_\mu,\xi_\tau)$ & per-flavour $\nu$ degeneracy (lepton asymmetry); the initial spectrum is $1/(e^{\eps-\mathrm{parfact}\,\xi}+1)$ \\
$\phi_e$ (\code{phie}) & electron degeneracy $\mu_e/\Tpl$ (settable; default $0$, or \code{"auto"} for the charge-neutrality $\phi(\Tpl)$; only its derivative QED EoS term, BURST's \code{glag6}, is not ported) \\
$q_{\rm baryon}$ ($q_{\rm b}$), $n_b$, $\eta$ ($\eta_b$, \code{eta0}) & charge per baryon balanced by the $e^\pm$, $n_{e^-}-n_{e^+}=q_{\rm baryon}n_b$; baryon number density; baryon-to-photon ratio \\
$\rhonu,\rhog$ & neutrino (all flavours and cp) / photon energy density [MeV$^4$] \\
$\rho_{\rm pl},P_{\rm pl}$ & plasma ($\gamma+e^\pm$) energy density and pressure: the free-gas $\rho_\gamma+\rho_{e^\pm}$ plus the finite-$T$ QED $\delta\rho_{\rm QED}$, $\delta P$ \\
$C_{\rm pl}$ & plasma heat capacity $\dd\rho_{\rm pl}/\dd T$ [MeV$^3$] (the variable \code{C} of \code{decoupling._plasma_eos}; distinct from the collision integral $C_{\nu_i}$ and the FD-cap multiplier $C$) \\
$\rho_{\rm DR}\equiv\rho_{\rm dr}\equiv\rho_{\rm ext}$ & extra-radiation energy density, $\rho_{\rm dr}=\code{delta_neff_init}\cdot(\frac{4}{11})^{4/3}\frac78\frac{\pi^2}{15}\Tcm^4$; standalone and coupled share LINX's \code{delta_neff_init} normalization (Section~\ref{sec:dneffconv}); BURST's \code{rhors} is the same quantity without the $(4/11)^{4/3}$ \\
$r_\nu$ & $\rhonu/\rhog$ summed over flavours and cp, the transport$\to$LINX transfer variable (\code{r_nu_of_Tg}); constant above the transport start, frozen after $e^\pm$ annihilation; the late-time $\Neff=\frac87(\frac{11}{4})^{4/3}r_\nu$ \\
$\Neff$ & effective number of neutrino species (Eq.~\eqref{eq:neffdef}); SM value $3.0433$ (NLO QED, default) / $3.0443$ (LO) \\
$\dNeff$ & has three distinct uses, all read from context: (i) the input inert-dark-radiation knob \code{delta_neff_init} (LINX \code{Delt_Neff_init} convention), the extra radiation at the transport start; (ii) a physics-induced output excess $\Neff-\Neff^{\rm SM}$ (e.g.\ the genuine QKE $+7.6\times10^{-4}$, a dynamically generated shift, not an input); (iii) a solver-vs-arbiter $\Neff$ residual \\
$\Yp$, $\DoverH$, $Y_i$ & $^4$He mass fraction $\Yp=4\,Y(^4\mathrm{He})$; deuterium ratio $\DoverH=Y(\mathrm D)/Y(p)$; $Y_i=n_i/n_b$ the LINX number abundances \\
\code{eta_fac}, \code{tau_n_fac}, \code{nuclear_rates_q} & baryon-to-photon and neutron-lifetime rescales ($1=$ fiducial) and per-reaction draws $q\sim\mathcal N(0,1)$ \\
\midrule
\multicolumn{2}{@{}l}{\textit{Constants, background, and finite-$T$ QED}} \\
$\GF$, $\sin^2\theta_W$, $\alpha$, $\me$ & Fermi constant, weak mixing angle (\code{SIN2THW}; BURST's $0.23$ is \code{SIN2THW_BURST}), fine-structure constant, bare electron mass \code{M_ELECTRON} (Section~\ref{sec:constants}) \\
$M_{\rm Pl}$, $\hbar$, $m_W$, $m_Z$ & Planck mass [MeV], $\hbar$ [MeV$\cdot$s], and the $W$/$Z$ masses of the QKE potentials \\
$H$ & Hubble rate [s$^{-1}$], $H=\hbar^{-1}\sqrt{8\pi(\rho_{\rm pl}+\rhonu+\rho_{\rm dr})/3M_{\rm Pl}^2}$; in Section~\ref{sec:qke} $H$ is the flavour Hamiltonian (below) \\
$\Gamma$ & weak interaction rate, $\Gamma\sim\GF^2T^5$, and the collision loss rate in $\Gamma/H$; the QKE rates $\Gamma_a$, $\Gamma_{ab}$ are below \\
$\delta m_e^2$, $\delta m_\gamma^2$, $m_{\rm eff}$ & thermal self-energies [MeV$^2$] and $m_{\rm eff}=\sqrt{\me^2+\delta m_e^2}$ (\code{thermo.m_e_eff}; not fed to the kernels) \\
\code{glag1}--\code{glag7}, pdep & BURST's seven Gauss--Laguerre plasma-moment integrals ($\rho$, $P$, $\dd\rho/\dd T$ and their QED corrections; \code{glag6}$=\dd\rho_e/\dd\phi_e$), and their momentum-dependent second-loop part \\
$\delta P_{e^3}$, $J$ & the NLO $\mathcal O(e^3)$ interaction pressure and its Gauss--Laguerre sum (Section~\ref{sec:nloqed}) \\
\midrule
\multicolumn{2}{@{}l}{\textit{Collision rates and the precision gate}} \\
$C_{\nu_i}$, $C_i$, $C[\rho]$ & collision integral (rate) of mode $i$; in the code $C_i=\Tcm^5\,\mathrm{net}_i\,\Theta_{\rm gate}$; $C[\rho]$ its density-matrix form.  Not the heat capacity $C_{\rm pl}$ or the FD-cap multiplier $C$ \\
$F$ ($F_r$), \code{flafact} & Pauli-blocked gain$-$loss factor $=$ forward$-$reverse (Eq.~\eqref{eq:Fr}), and the per-process flavour/spin weight multiplying it \\
$S_r$, $\langle|M_r|^2\rangle$ & symmetry factor and spin-summed squared matrix element of process $r$ \\
$W$, $K$, $M$, $n_{\rm GL}$ & static weight tensor $(N,N,K)$ of the grid-only kernels, its quadrature slots per bin pair ($K=15$ \code{nunu}, $10$ \code{nunubar}), the interpolation window ($M=6$), and Gauss nodes per segment ($5$); $K$ is also the QKE frame generator \\
$\mathrm{idx}_{3,4}$, $\mathrm{lw}_{3,4}$, $L[p,a]$ & window indices and Lagrange weights of the off-grid reads of legs 3 and 4, $\log f(p)=\sum_aL[p,a]\log f_a$ \\
\code{net}, \code{frs} & per-bin collision decomposition.  \code{net}$=$forward$-$reverse rate is the physical signal that drives $\dd f/\dd\ln a$; \code{frs}$=$forward$+$reverse is the total throughput, which normalizes the cancellation in the gate and its noise floor (with \code{net} it also gives the forward and reverse rates behind the affine $A$, $B$ and the QKE $\Gamma_a$).  Near equilibrium the two nearly cancel, so $|\code{net}|\ll\code{frs}$ \\
\code{numer}, \code{denom} & $\sum$\code{net} and $\sum$\code{frs} over the five kernels, the gate's inputs \\
\code{prec}, \code{prec_trans} & equilibrium noise floor for the precision gate, $|\sum\mathrm{net}|/\sum\mathrm{frs}$ per (flavour, cp, bin) at the equilibrium spectrum (\code{PREC_FLOOR} where $\sum\mathrm{frs}\le0$) \\
\code{test_ratio}, \code{nuratiotol}, $\Theta_{\rm gate}$ & $|\sum\mathrm{net}|/(\sum\mathrm{frs}\cdot\mathrm{prec})$, its threshold $3$, and the resulting $0/1$ gate (Eq.~\eqref{eq:gate}) \\
\code{doccfracdt}, \code{drhonudt} & the gated $\dd f/\dd t$ [s$^{-1}$] and the $\nu$--plasma energy-exchange rate $\dd\rho_\nu/\dd t|_{\rm exch}$ [MeV$^4$\,s$^{-1}$] \\
\code{leg4} & the summed net rate of the three $\nu$--$e^\pm$ channels, BURST's partner-species slot~4 (the $e^\pm$), whose $\eps^3$ moment is the energy exchanged with the plasma; not the momentum leg $p_4$ (``leg 4'' below) \\
$A,B$ ($g,l$) & affine collision RHS coefficients: per $\ln a$ (\code{make_split_rates}, the decoupling/hybrid/QKE builders) $\dd f/\dd\ln a=A-Bf$ with $A=g/H$, $B=(g+l)/H$, $g,l$ the gain/loss rates [s$^{-1}$]; \code{assemble_split_rhs} returns the s$^{-1}$ form $\dd f/\dd t=A-Bf$, $A=g$, $B=g+l$ (carrying $\Tcm^5$).  Eq.~\eqref{eq:affine} writes the per-$\ln a$ form as $A=g$, $B=g+l$ with $g,l$ already divided by $H$ \\
\code{OCC_FLOOR} & always-on clamp of $f$ to $[10^{-300},1]$ before every kernel call (an identity on physical occupations) \\
\midrule
\multicolumn{2}{@{}l}{\textit{Flavour mixing (QKE)}} \\
$\rho$, $\bar\rho$ & $3\times3$ Hermitian flavour density matrix per (cp, bin), $\rho[\text{cp},i,m,m]\equiv f[m,\text{cp},i]$, its off-diagonal entries the flavour coherences ($\rho_{\rm diag}$, $\rho_{\rm off}$); $\bar\rho$ the $\bar\nu$ matrix.  Not an energy density \\
$H$ ($H_{\rm flav}$) & flavour Hamiltonian [s$^{-1}$], $H_{\rm vac}+H_{\rm mat}^{(1)}+H_{\rm mat}^{(2)}+H_{\rm si}^{(1)}+H_{\rm si}^{(2)}$ (vacuum, asymmetric and thermal MSW, $\nu$--$\nu$; the solvers pass the two $\nu$--$\nu$ terms as zero) \\
$m^2_{\rm weak}$, $U$, $M^2$ & $m^2_{\rm weak}=UM^2U^\top$ (trace removed; $-m^2_{\rm weak}$ for $\bar\nu$), the real PMNS matrix $U=R_{23}R_{13}R_{12}$ (angles $\theta_{ij}$), and $M^2=\mathrm{diag}(0,\Delta m^2_{21},\Delta m^2_{31})$ [MeV$^2$].  In Section~\ref{sec:qkesplit} $U$ also names the Strang step's unitary $Ve^{-i\Lambda\dd t}V^\dagger$ \\
$n_p$ & proton (net-$e^-$) density of $H_{\rm mat}^{(1)}$, $n_p=q_{\rm b}n_b\propto\eta_b$, scaled per call by \code{matter_asym_scale} \\
$L_{\rm coh}$, $\Lambda_{\rm off}$ & the $9\times9$ coherent Liouvillian $-i[H,\cdot]-\Lambda_{\rm off}$ and its off-diagonal-only dephasing \\
$\Gamma_a$, $\Gamma_{ab}$ & ungated collision throughput $g+l$ of flavour $a$ (\code{assemble.decoherence_rate}) and the coherence damping rate $\tfrac12(\Gamma_a+\Gamma_b)$ \\
$\omega$, $\Delta$ & largest eigenvalue splitting of $H$ ($\omega\,\dd t$ is the phase per step) and the vacuum oscillation frequency $\delta m^2/2p$ \\
$V$, $\Lambda$ & eigenvectors and eigenvalues of $H$ in the Strang unitary $Ve^{-i\Lambda\dd t}V^\dagger$ (not $\Lambda_{\rm off}$, nor the PMNS $U$) \\
$K$, $\tilde R$ & moving-eigenframe generator, $e^{-iK\dd t}=V_{\rm end}V_{\rm start}^\dagger$, and the half-step rotation $\tilde R=e^{-iK\dd t/2}$ \\
$\Theta$ & non-adiabaticity measure $\|[H_{\rm vac},\rho]\|^2/(\|H_{\rm vac}\|^2\|\rho\|^2)$ (Frobenius norms); not $\Theta_{\rm gate}$ \\
$L_{oo}$, $L_{od}$, $s$ & the (off-diagonal, off-diagonal) and (off-diagonal, diagonal) Liouvillian blocks and the coherent source $s=L_{od}\rho_{\rm diag}$ (\code{offdiag_blocks}), so the quasi-steady coherence is $-L_{oo}^{-1}s$.  The sourced step's collision rate is $s_{\rm coll}=(f_{\rm coll}-f)/\dd t$ (\code{src}) \\
$A_{\rm gain}$, $B_{\rm loss}$ & Hermitian gain and loss matrices of the density-matrix collision $\dd\rho/\dd t|_{\rm coll}=A_{\rm gain}-\{B_{\rm loss},\rho\}$; on a diagonal $\rho$ it reduces to the classical kernel and $2\,\mathrm{diag}(B_{\rm loss})$ is the ungated throughput $\Gamma_a$ \\
$n_{\rm steps}$ (\code{n_steps}) & operator-split macro-steps of the QKE solve (default 6000) \\
\code{num_drift}, \code{min_f} & per-bin number drift and minimum occupation, the QKE conservation diagnostics \\
\midrule
\multicolumn{2}{@{}l}{\textit{Terms}} \\
pin & a frozen reference value the port must reproduce on demand (e.g.\ the LO oracle pin $\Neff=3.044225186$ at \code{SIN2THW_BURST}; the default NLO pin is $3.0433222$) \\
oracle & the trusted source of truth a result is checked against (here, the compiled Fortran snapshot) \\
cold / warm & a compile-paying first call vs.\ a graph-reusing later call (Section~\ref{sec:jaxprimer}) \\
traced / static & a per-call (no-recompile) vs.\ a build-time (recompile-on-change) parameter (Section~\ref{sec:jaxprimer}) \\
precision gate & per-bin test that zeroes a bin's \code{net} rate when $|\code{net}|$ falls below the equilibrium round-off floor \code{prec} (where forward and reverse nearly cancel); prevents the solver from integrating floating-point noise as a physical spectral distortion \\
gate (verb), ``gated against'' & to validate/anchor a computed result against a trusted reference (``gated against the oracle'' $=$ regression-checked for agreement).  Distinct from the precision gate, which instead masks sub-floor bin rates \\
conservation sum rules & moment integrals of the net rate (number $\int\eps^2\,\code{net}$, energy $\int\eps^3\,\code{net}$, and CP-symmetry checks) that must vanish in equilibrium; a free, continuous correctness monitor \\
affine $A-Bf$ form & writing a stiff bin's RHS as $\dd f/\dd s=A-Bf$, which an exponential-midpoint step advances in closed form (exactly); the fixed point $f^\star=A/B$ is the bin's local quasi-static equilibrium (its QSS attractor).  See $A,B$ above \\
QSS attractor & quasi-static equilibrium: the fixed point $f^\star=A/B$ of the affine RHS onto which a fast, collision-dominated bin is driven.  Tracking it lets the stiff head take large steps without losing accuracy \\
head / tail, qualified & the hybrid's fixed-step exponential-midpoint stretch $[0,\code{lna_sw})$ and adaptive Tsit5 remainder; \code{qualified} reports that the handoff passed its QSS check (\code{qss_dist}$\le$\code{qual_tol}) rather than falling back to the full adaptive solve \\
full / semi / qke & the coupled-BBN transport modes (Section~\ref{sec:fullsemi}): full co-evolution of $\Tpl$ and the spectrum, a spectrum-only solve on LINX's plasma history, and the genuine QKE solve \\
mixing-off reduction & the QKE solve with \code{mixing=False}, bit-identical to the classical transport; the anchor of the QKE validation \\
\code{tail_closure} (\code{"srcfix"}) & the default off-grid tail closure of the QKE solver, parameter-free: \code{fd_decay} carries the live grid-edge $f/f^{\rm eq}$ outward along the FD shape, \code{monotone_read} clamps the on-grid interpolant to its node bracket.  The opt-in alternative is the FD-envelope \code{"cap"} below \\
FD-envelope cap (\code{fd_cap_logC}, \code{fd_cap_C}) & caps an off-grid source occupation at a multiple $C$ of Fermi--Dirac, $f(p)\le C\,f^{\rm eq}(p)$.  The kernels' \code{fd_cap_logC}$=\log C$ defaults to $+\infty$, a bitwise no-op; under the opt-in QKE closure \code{tail_closure="cap"}, $C=$\code{fd_cap_C}, default \code{interp.FD_ENVELOPE_C}$=10$ \\
golden & a frozen reference test (input $+$ expected output) a result is regression-checked against, as opposed to a single pin value \\
arbiter & a separate, ultra-tight-tolerance reference run treated as ground truth to measure another run's error \\
L-stable & property of a stiff integrator whose step amplitude stays bounded (in fact $\to0$) even for arbitrarily fast-decaying modes, so a too-large step cannot blow up the stiff bins \\
\code{optimization_barrier} & a JAX/XLA directive marking a value opaque so the compiler may not pre-evaluate it at build time with a different floating-point summation order; needed where a near-perfect cancellation makes the answer summation-order sensitive \\
scaling-and-squaring (\code{max_squarings}) & the algorithm JAX uses for the matrix exponential $e^{M}$ (scale down, exponentiate, square back up); a very fast rotation needs enough squarings or the result silently underflows to \code{NaN} \\
order-$p$ convergence & the discretization error of a stepper falls like $1/n_{\rm steps}^{p}$; order $2$ means doubling the step count cuts the error by $\sim4\times$ \\
dispatch-bound & wall time dominated by per-call GPU kernel-launch (dispatch) overhead rather than arithmetic; the reason batching many solves with \code{vmap} amortizes the cost \\
leg 4 (momentum) & the fourth momentum leg $p_4$ of a $1{+}2\to3{+}4$ kernel ($p_4=p_1+E_2-E_3$ for scattering, $E_{\rm out}+E_{\rm in}-p_1$ in \code{nuepma}): it lands between grid points and $f(p_4)$ is read by off-grid interpolation.  Not \code{leg4}, BURST's $e^\pm$ partner slot above \\
\midrule
\multicolumn{2}{@{}l}{\textit{Lepton kernels}} \\
\code{meps} & the lepton kernels' ``mass$/\Tcm$'' argument, $\me/\Tcm$ with the bare mass \\
regimes & a lepton-kernel bin's kinematic case: A/B in \code{nuer1} (split at $p_1=\code{meps}/2$), R1--R4 in \code{nuer2} and \code{nuepma} (split at the \code{pcase} momenta); not BURST's R1/R2 angular reductions that name \code{nuer1}/\code{nuer2} \\
\code{ecut1}--\code{ecut3}, \code{etrans1,2}, \code{elim1} & kinematic cut energies [$\Tcm$ units] bounding a lepton kernel's tiles, defined per kernel (e.g.\ $\code{ecut3}=\sqrt{p_1^2+\code{meps}^2}$ in \code{nuer1}/\code{nuer2}; \code{nuer2}'s and \code{nuepma}'s $\code{ecut1}=p_1+\code{meps}^2/(4p_1)$); \code{etrans1,2} are \code{nuepma}'s inner cuts per outer node, \code{elim1} \code{nuer2}'s inner $E_2$ limit \\
tiles, \code{genout}/\code{genin} & the fixed quadrature sub-intervals (Gauss--Legendre on $[a,b]$, Gauss--Laguerre on $[a,\infty)$); \code{nuepma}'s \code{genout}/\code{genin} segments switch from Legendre to Laguerre when wider than $200$ (BURST), or than \code{lag_tail_gap} under \code{lag_tail_fix} \\
\code{n_gl}, \code{n_lag} & Gauss--Legendre / Gauss--Laguerre orders of the kernels (BURST's $5$, $5$) \\
\code{lag_tail_fix}, \code{lag_tail_gap} & the default-on wide-regime quadrature correction of \code{nuer1}, \code{nuer2} and \code{nuepma}, and the segment width ($6\,\Tcm$) above which a segment takes the Gauss--Laguerre tail (Section~\ref{sec:lepkernels}) \\
\code{use_sched}, ``sched'' & performance-only switch selecting the lepton-kernel implementation (per-bin \code{select} gather, \code{union} mask, or \code{False} reference); all give identical physics, trading speed against bit-reproducibility, never accuracy \\
\code{pcase} & the $e^\pm$-mass-threshold momenta (\code{nuepma}: $\code{meps}/2$, $\code{meps}(1+\sqrt5)/4$, $\code{meps}$; \code{nuer2}: $\code{meps}(\sqrt5-1)/4$, $\code{meps}/(2\sqrt2)$, $\code{meps}/2$) that segment a lepton-kernel bin's energy integral into the kinematic regimes R1--R4 \\
\bottomrule
\end{longtable}
\end{center}

% ===========================================================================
\section{References}
\label{sec:references}
% ===========================================================================

\begin{itemize}[nosep]
\item E.~Grohs, G.~M.~Fuller, C.~T.~Kishimoto, M.~W.~Paris, A.~Vlasenko,
  ``Neutrino energy transport in weak decoupling and big bang nucleosynthesis,''
  Phys.\ Rev.\ D \textbf{93}, 083522 (2016),
  \href{https://arxiv.org/abs/1512.02205}{arXiv:1512.02205}. (The BURST paper.)
\item E.~Grohs, G.~M.~Fuller, ``Big Bang Nucleosynthesis,'' Handbook of Nuclear
  Physics (2023), \href{https://arxiv.org/abs/2301.12299}{arXiv:2301.12299}.
\item E.~Grohs, G.~M.~Fuller, C.~T.~Kishimoto, M.~W.~Paris, ``Lepton asymmetry,
  neutrino spectral distortions, and big bang nucleosynthesis,''
  \href{https://arxiv.org/abs/1612.01986}{arXiv:1612.01986}.
\item C.~Giovanetti, M.~Lisanti, H.~Liu, S.~Mishra-Sharma, J.~T.~Ruderman,
  ``LINX: A Fast, Differentiable, and Extensible Big Bang Nucleosynthesis
  Package,'' \href{https://arxiv.org/abs/2408.14538}{arXiv:2408.14538}.
\item C.~Giovanetti, M.~Lisanti, H.~Liu, J.~T.~Ruderman, ``Joint CMB and BBN
  Constraints on Light Dark Sectors with Dark Radiation,'' Phys.\ Rev.\ Lett.\
  \textbf{129}, 021302 (2022),
  \href{https://arxiv.org/abs/2109.03246}{arXiv:2109.03246}.
\item K.~Akita, M.~Yamaguchi, ``A precision calculation of relic neutrino
  decoupling,'' \href{https://arxiv.org/abs/2005.07047}{arXiv:2005.07047};
  J.~Froustey, C.~Pitrou, M.-C.~Volpe, ``Neutrino decoupling including flavour
  oscillations and primordial nucleosynthesis,''
  \href{https://arxiv.org/abs/2008.01074}{arXiv:2008.01074};
  J.~J.~Bennett et al., ``Towards a precision calculation of $\Neff$ in the
  Standard Model II,'' \href{https://arxiv.org/abs/2012.02726}{arXiv:2012.02726}.
\item NLO QED: J.~J.~Bennett et al., ``Towards a precision calculation of the
  effective number of neutrinos $\Neff$ in the Standard Model,''
  \href{https://arxiv.org/abs/1911.04504}{arXiv:1911.04504} (the $\mathcal O(e^3)$
  interaction pressure ported here); G.~Cielo, M.~Escudero, G.~Mangano,
  O.~Pisanti, ``$\Neff$ in the Standard Model at NLO is $3.043$,''
  \href{https://arxiv.org/abs/2306.05460}{arXiv:2306.05460}.
\end{itemize}

The list above is a curated key-reading list rather than an exhaustive
bibliography; sources used only in one place are linked inline where they appear.
This manual is intended to be self-contained: every numerical choice and caveat
is explained where it is used, so no familiarity with the repository's internal
development notes is assumed.

\end{document}










